\documentclass[11pt,a4paper]{article}
\usepackage[italian]{babel}
\usepackage[margin=2.5cm]{geometry}
\usepackage{parskip}
\usepackage{amsmath, amssymb, amsthm}
\usepackage{booktabs}
\usepackage{array}
\usepackage{xcolor}
\usepackage{needspace}
\usepackage{enumitem}
\usepackage{listings}
\usepackage{tikz}
\usepackage[most]{tcolorbox}
\usepackage[hidelinks]{hyperref}
\usetikzlibrary{decorations.pathreplacing, shapes.geometric, arrows.meta, positioning}

% Niente vedove/orfane: i paragrafi non lasciano righe isolate a fine/inizio pagina
\widowpenalty=10000
\clubpenalty=10000

% Elenchi più compatti
\setlist{itemsep=2pt, parsep=0pt, topsep=4pt, partopsep=0pt}

% --- Riquadri con bordo nero, senza numero (si spezzano tra due pagine) ---
% Uso: \begin{definizione} ... \end{definizione}
%      \begin{teorema}[Nome] ... \end{teorema}  (il nome è facoltativo)
\tcbset{nota/.style={enhanced jigsaw, breakable, colback=white,
  colframe=black, boxrule=0.5pt, sharp corners}}
\NewTColorBox{definizione}{}{nota,
  before upper={\textbf{Definizione.}\ }}
\NewTColorBox{teorema}{o}{nota,
  before upper={\textbf{Teorema\IfValueT{#1}{ (#1)}.}\ }}
\NewTColorBox{proposizione}{o}{nota,
  before upper={\textbf{Proposizione\IfValueT{#1}{ (#1)}.}\ }}
\NewTColorBox{proprieta}{o}{nota,
  before upper={\textbf{Proprietà\IfValueT{#1}{ (#1)}.}\ }}
\NewTColorBox{assioma}{o}{nota, boxrule=1pt,
  before upper={\textbf{Assioma\IfValueT{#1}{ (#1)}.}\ }}

% --- Ambienti semplici (senza riquadro) ---
% Il titolo sta in cima al blocco, anche se il contenuto inizia con un riquadro o
% con codice; e non resta mai da solo in fondo a una pagina.
% Uso: \begin{esempio}[Titolo facoltativo] ... \end{esempio}
\newcommand{\newrunin}[2]{%
  \NewDocumentEnvironment{#1}{o}{%
    \par\Needspace{4\baselineskip}\addvspace{\medskipamount}\noindent
    \textbf{#2}\IfValueT{##1}{ (##1)}\textbf{.}\ \ignorespaces}%
  {\par\addvspace{\medskipamount}}}
\newrunin{esempio}{Esempio}
\newrunin{esercizio}{Esercizio}
\newrunin{osservazione}{Osservazione}

% V e F nelle tavole di verità
\newcommand{\V}{\textsf{V}}
\newcommand{\F}{\textsf{F}}

% --- Codice C ---
% Uso: \begin{codice} ... \end{codice}      codice C numerato
%      \begin{comando} ... \end{comando}    comandi da terminale
%      \begin{risultato} ... \end{risultato} output di un programma
%      \lstinline|printf("ciao");|           codice dentro al testo
\lstdefinestyle{c}{
  language=C,
  basicstyle=\ttfamily\small,
  keywordstyle=\color{blue}\bfseries,
  commentstyle=\color{gray}\itshape,
  stringstyle=\color{red!70!black},
  numbers=left,
  numberstyle=\tiny\color{gray},
  numbersep=8pt,
  columns=flexible,
  keepspaces=true,
  breaklines=true,
  showstringspaces=false,
  tabsize=4,
  morekeywords={include, define},
  % lettere accentate nei commenti
  literate={à}{{\`a}}1 {è}{{\`e}}1 {é}{{\'e}}1 {ì}{{\`i}}1
           {ò}{{\`o}}1 {ù}{{\`u}}1 {È}{{\`E}}1 {À}{{\`A}}1
}
\lstset{style=c}
\newtcblisting{codice}{listing only, nota, left=6mm,
  listing options={style=c}}
\newtcblisting{comando}{listing only, nota,
  listing options={basicstyle=\ttfamily\small, numbers=none, language={},
    columns=flexible, keepspaces=true, breaklines=true}}
\newtcblisting{risultato}{listing only, enhanced jigsaw, breakable,
  colback=black!5, colframe=black, boxrule=0.5pt, sharp corners,
  listing options={basicstyle=\ttfamily\small, numbers=none, language={},
    columns=flexible, keepspaces=true, breaklines=true}}

% --- Diagrammi di flusso (simboli standard) ---
\tikzset{
  terminale/.style = {ellipse, draw, minimum width=2.5cm, minimum height=0.9cm},
  io/.style        = {trapezium, trapezium left angle=70, trapezium right angle=110,
                      draw, minimum width=2.5cm, minimum height=0.9cm},
  processo/.style  = {rectangle, draw, minimum width=2.5cm, minimum height=0.9cm},
  decisione/.style = {diamond, draw, aspect=2, minimum width=2.5cm},
  freccia/.style   = {thick, -{Stealth}}
}

\title{Matematica discreta\\[0.4em]\large Appunti}
\author{Marco Cerbella}
\date{Università degli Studi di Perugia -- A.A. 2026/2027\\ Aggiornato al \today}

\begin{document}
\maketitle
\setcounter{tocdepth}{1}
\tableofcontents
\clearpage

\section[Logica (Lezione 1 -- 21 settembre 2026)]{Logica}
\noindent\emph{Lezione 1 -- 21 settembre 2026}
\subsection{Enunciati e connettivi}

\begin{definizione}
Un \emph{enunciato} (o \emph{proposizione}) è una frase di cui possiamo
dire con certezza se è \V\ (vera) oppure \F\ (falsa). Gli enunciati si
indicano con $p, q, r, \dots$
\end{definizione}

\begin{esempio}
Sono enunciati: ``$5 > 3$'', ``$99$ è divisibile per $3$'',
``$508$ è multiplo di $5$''. Invece ``$n$ è un numero primo'' \textbf{non}
è un enunciato, perché la sua verità dipende da $n$.
\end{esempio}

Dati due enunciati $p, q$ se ne costruiscono di nuovi con i
\emph{connettivi}:
\begin{center}
\begin{tabular}{cll}
$p \land q$ & ``$p$ e $q$'' & congiunzione \\
$p \lor q$ & ``$p$ o $q$'' & disgiunzione \\
$\neg p$ & ``non $p$'' & negazione \\
$p \to q$ & ``se $p$ allora $q$'' & implicazione \\
$p \leftrightarrow q$ & ``$p$ se e soltanto se $q$'' & coimplicazione
\end{tabular}
\end{center}

\subsubsection*{Tavola di verità}

\begin{center}
\begin{tabular}{cc|c|c|c|c}
$p$ & $q$ & $p \land q$ & $p \lor q$ & $p \to q$ & $p \leftrightarrow q$ \\ \hline
\V & \V & \V & \V & \V & \V \\
\V & \F & \F & \V & \F & \F \\
\F & \V & \F & \V & \V & \F \\
\F & \F & \F & \F & \V & \V
\end{tabular}
\end{center}

\subsection{Proprietà dei connettivi}

\begin{proposizione}
Siano $p, q, r$ enunciati. Allora:
\begin{enumerate}
    \item $p \lor q \iff q \lor p$;
    \item $p \land q \iff q \land p$;
    \item $(p \land q) \lor r \iff (p \lor r) \land (q \lor r)$;
    \item $(p \lor q) \land r \iff (p \land r) \lor (q \land r)$;
    \item $\neg(p \land q) \iff (\neg p) \lor (\neg q)$;
    \item $\neg(p \lor q) \iff (\neg p) \land (\neg q)$;
    \item $p \to q \iff (\neg p) \lor q$;
    \item $p \to q \iff \neg q \to \neg p$;
    \item $p \leftrightarrow q \iff (p \to q) \land (q \to p)$.
\end{enumerate}
\end{proposizione}

Due enunciati composti sono equivalenti ($\iff$) quando hanno la stessa
tavola di verità. Le prime due proprietà sono la commutatività di $\lor$ e
$\land$; la 3 e la 4 sono le proprietà distributive; la 5 e la 6 sono le
leggi di De Morgan; la 8 è la \emph{contronominale}.

\begin{proof}[Verifica della proprietà 7]
Confrontiamo le tavole di verità di $p \to q$ e di $(\neg p) \lor q$:
\begin{center}
\begin{tabular}{cc|c|c|c}
$p$ & $q$ & $p \to q$ & $\neg p$ & $(\neg p) \lor q$ \\ \hline
\V & \V & \V & \F & \V \\
\V & \F & \F & \F & \F \\
\F & \V & \V & \V & \V \\
\F & \F & \V & \V & \V
\end{tabular}
\end{center}
Le colonne di $p \to q$ e di $(\neg p) \lor q$ sono uguali.
\end{proof}

\begin{proof}[Verifica della proprietà 4]
Servono otto righe, una per ogni combinazione dei valori di $p, q, r$:
\begin{center}
\begin{tabular}{ccc|c|c|c|c|c}
$p$ & $q$ & $r$ & $p \lor q$ & $(p \lor q) \land r$ & $p \land r$ &
$q \land r$ & $(p \land r) \lor (q \land r)$ \\ \hline
\V & \V & \V & \V & \V & \V & \V & \V \\
\V & \V & \F & \V & \F & \F & \F & \F \\
\V & \F & \V & \V & \V & \V & \F & \V \\
\V & \F & \F & \V & \F & \F & \F & \F \\
\F & \V & \V & \V & \V & \F & \V & \V \\
\F & \V & \F & \V & \F & \F & \F & \F \\
\F & \F & \V & \F & \F & \F & \F & \F \\
\F & \F & \F & \F & \F & \F & \F & \F
\end{tabular}
\end{center}
Le colonne di $(p \lor q) \land r$ e di $(p \land r) \lor (q \land r)$
coincidono. Le altre proprietà si verificano allo stesso modo.
\end{proof}

\subsection{Formule aperte e quantificatori}

\begin{definizione}
Una \emph{formula aperta} è un'affermazione che contiene un'incognita, un
parametro o una variabile.
\end{definizione}

\begin{esempio}
\[
p(x):\ x + 3 = 0, \qquad
q(n):\ \text{``$n$ è un numero pari''}, \qquad
r(k):\ \text{``$k$ è multiplo di $5$''}.
\]
\end{esempio}

Una formula aperta non è né vera né falsa. Per passare dalla formula aperta
a un enunciato ci sono tre modi.

\begin{enumerate}
    \item \textbf{Sostituzione}: si sostituisce alla variabile un valore
          preciso. Ad esempio $p(3)$ è l'enunciato $3 + 3 = 0$ (falso),
          $p(-3)$ è $-3 + 3 = 0$ (vero), $q(8)$ è ``$8$ è un numero pari''
          (vero).

    \item \textbf{Quantificatore esistenziale} $\exists$ (``esiste''):
          \[
          \exists\, x \in \mathbb{R} : p(x)
          \]
          si legge ``esiste $x$ reale tale che $x + 3 = 0$'' ed è vero
          (basta $x = -3$). Invece
          $\exists\, x \in \mathbb{R} : x^2 = -10$ è falso.

    \item \textbf{Quantificatore universale} $\forall$ (``per ogni''):
          $\forall\, x \in \mathbb{R} : x + 3 = 0$ è falso, mentre
          $\forall\, x \in \mathbb{R} : x^2 \geq 0$ è vero.
\end{enumerate}

\begin{proposizione}[Negazione dei quantificatori]
\[
\neg\big(\forall x : p(x)\big) \iff \exists\, x : \neg p(x), \qquad
\neg\big(\exists\, x : p(x)\big) \iff \forall x : \neg p(x).
\]
Il secondo enunciato si scrive anche $\nexists\, x : p(x)$.
\end{proposizione}

Per negare una frase con un quantificatore si scambia $\forall$ con
$\exists$ e si nega la formula.

\subsection{Implicazione: condizioni necessarie e sufficienti}

L'implicazione $p \to q$ si legge in tre modi:
\begin{itemize}
    \item ``se $p$ allora $q$'';
    \item ``$p$ è condizione \emph{sufficiente} per $q$'';
    \item ``$q$ è condizione \emph{necessaria} per $p$''.
\end{itemize}

\begin{esempio}
Per ogni $n \in \mathbb{N}$: se $n$ è divisibile per $4$ (è $p$), allora
$n$ è pari (è $q$). Essere divisibile per $4$ è sufficiente per essere
pari; essere pari è necessario per essere divisibile per $4$.
\end{esempio}

\section[Insiemi (Lezione 1 -- 21 settembre 2026)]{Insiemi}
\noindent\emph{Lezione 1 -- 21 settembre 2026}
\subsection{Insiemi ed elementi}

\begin{definizione}
Un \emph{insieme} è una collezione (o famiglia) di oggetti di cui possiamo
dire con precisione se un oggetto appartiene o no alla famiglia. Un oggetto
di un insieme si chiama \emph{elemento}.
\end{definizione}

Gli insiemi si indicano di solito con lettere maiuscole ($A, B, X, S, V,
\dots$), gli elementi con lettere minuscole ($x, y, v, k, \dots$).
\begin{itemize}
    \item $x \in S$: l'elemento $x$ appartiene all'insieme $S$;
    \item $y \notin S$: $y$ non appartiene a $S$.
\end{itemize}

Un insieme si può dare in due modi.
\begin{enumerate}
    \item \textbf{Per elencazione}: si scrivono gli elementi tra graffe.
          Ad esempio $V = \{1, 3, 5, 88\}$ definisce l'insieme $V$ per
          elencazione.
    \item \textbf{Per descrizione} (o per caratteristica): ad esempio
          $X = \{\text{le lettere dell'alfabeto inglese}\}$.
\end{enumerate}

Quando la descrizione usa una formula aperta si scrive
\[
W = \{x \mid x^2 - 3x + 2 = 0\}, \qquad
A = \{x \in \mathbb{R} \mid P(x)\}, \qquad
B = \{x \in \mathbb{R} \mid x > 3\} = \{x : x > 3\},
\]
dove ``$\mid$'' (oppure ``$:$'') si legge ``tale che''.

Un insieme con un solo elemento, come $U = \{4\}$ (per cui $4 \in U$ e
$3 \notin U$), si chiama \emph{singoletto} (in inglese \emph{singleton}).

L'\emph{insieme vuoto}, indicato con $\emptyset = \{\ \}$, è l'insieme senza
elementi: per ogni $x$ vale $x \notin \emptyset$ (cioè
$\nexists\, x : x \in \emptyset$).

\subsection{Sottoinsiemi e uguaglianza}

\begin{definizione}
Siano $A, B$ insiemi. $B$ è un \emph{sottoinsieme} di $A$ se
\[
\forall\, x \in B \implies x \in A.
\]
Si scrive $B \subset A$ oppure $B \subseteq A$; equivalentemente
$A \supseteq B$ (o $A \supset B$).
\end{definizione}

\begin{definizione}
Due insiemi $S, T$ sono \emph{uguali}, $S = T$, se hanno gli stessi
elementi, cioè
\[
x \in S \iff x \in T.
\]
\end{definizione}

\begin{osservazione}
$S = T$ se e solo se $S \subseteq T$ e $T \subseteq S$.
\end{osservazione}

Dato un insieme $A$ valgono sempre $\emptyset \subseteq A$ e
$A \subseteq A$.

\subsubsection*{Insiemi numerici}

\begin{center}
\begin{tabular}{ll}
$\mathbb{N} = \{0, 1, 2, 3, \dots\}$ & numeri naturali \\
$\mathbb{Z} = \{0, \pm 1, \pm 2, \pm 3, \dots\}
              = \{\dots, -3, -2, -1, 0, 1, 2, \dots\}$ & numeri interi \\
$\mathbb{Q} = \left\{ \dfrac{a}{b} \;\middle|\; a, b \in \mathbb{Z},\
              b \neq 0 \right\}$ & numeri razionali \\
$\mathbb{R}$ & numeri reali \\
$\mathbb{C}$ & numeri complessi
\end{tabular}
\end{center}

Valgono le inclusioni
\[
\mathbb{N} \subseteq \mathbb{Z} \subseteq \mathbb{Q} \subseteq
\mathbb{R} \subseteq \mathbb{C}.
\]
Con $i$ si indica l'unità immaginaria, $i^2 = -1$. Ad esempio:
\begin{center}
\begin{tabular}{llll}
$-2 \in \mathbb{Z}$, & $\frac12 \in \mathbb{Q}$, &
$\sqrt{2}, \pi \in \mathbb{R}$, & $i \in \mathbb{C}$, \\
$-2 \notin \mathbb{N}$, & $\frac12 \notin \mathbb{Z}$, &
$\sqrt{2}, \pi \notin \mathbb{Q}$, & $i \notin \mathbb{R}$.
\end{tabular}
\end{center}

\subsection{Insieme delle parti}

\begin{definizione}
Sia $A$ un insieme. L'\emph{insieme delle parti} di $A$ è l'insieme di
tutti i sottoinsiemi di $A$:
\[
\mathcal{P}(A) = \{B \mid B \subseteq A\}.
\]
\end{definizione}

\begin{osservazione}
Per ogni insieme $A$: $\emptyset \subseteq A$ e $\emptyset \in
\mathcal{P}(A)$; $A \subseteq A$ e $A \in \mathcal{P}(A)$.
\end{osservazione}

\begin{esempio}
Sia $A = \{1, 2, 3\}$. Allora
\[
\mathcal{P}(A) = \big\{\emptyset,\ \{1\},\ \{2\},\ \{3\},\
\{1,3\},\ \{2,3\},\ \{1,2\},\ A\big\}.
\]
Bisogna distinguere bene appartenenza e inclusione:
\begin{itemize}
    \item $1 \in A$, \quad $\{1\} \subseteq A$, \quad
          $\{1, 3\} \subseteq A$;
    \item $\{1\} \in A$: \textbf{no}; \quad $\{1\} \in \mathcal{P}(A)$:
          sì;
    \item $A \subseteq \mathcal{P}(A)$: \textbf{no}; \quad
          $A \in \mathcal{P}(A)$: sì.
\end{itemize}
\end{esempio}

\begin{esempio}
$\mathcal{P}(\emptyset) = \{\emptyset\}$ e
\[
\mathcal{P}(\{\emptyset\}) = \mathcal{P}(\mathcal{P}(\emptyset))
= \{\emptyset, \{\emptyset\}\}.
\]
\end{esempio}

\subsection{Operazioni tra insiemi}

Siano $A, B$ insiemi.

\begin{definizione}
\begin{itemize}
    \item \textbf{Unione}: $A \cup B = \{x \mid x \in A \lor x \in B\}$.
    \item \textbf{Intersezione}: $A \cap B = \{x \mid x \in A \land
          x \in B\}$.
    \item \textbf{Differenza} tra $A$ e $B$:
          $A \setminus B = \{x \mid x \in A \text{ e } x \notin B\}$.
    \item \textbf{Differenza simmetrica} tra $A$ e $B$:
          $A \,\triangle\, B = (A \setminus B) \cup (B \setminus A)$.
\end{itemize}
$A$ e $B$ sono \emph{disgiunti} se $A \cap B = \emptyset$.
\end{definizione}

Esercizio: $A \,\triangle\, B = (A \cup B) \setminus (A \cap B)$.

\begin{definizione}
Se $A \subseteq U$, il \emph{complementare} di $A$ in $U$ è l'insieme
\[
A^c = U \setminus A.
\]
\end{definizione}

\begin{esempio}
Siano $A = \{1, 2, 3, 4\} = \{4, 1, 3, 2\}$ (l'ordine non conta) e
$B = \{1, 3, 6, 9, 12\}$. Allora
\begin{align*}
A \cup B &= \{1, 2, 3, 4, 6, 9, 12\}, & A \cap B &= \{1, 3\}, \\
A \setminus B &= \{2, 4\}, & B \setminus A &= \{6, 9, 12\}, \\
A \,\triangle\, B &= \{2, 4, 6, 9, 12\}.
\end{align*}
\end{esempio}

\subsection{Prodotto cartesiano}

\begin{definizione}
Siano $A, B$ insiemi. Il \emph{prodotto cartesiano} di $A$ e $B$ è
\[
A \times B = \{(a, b) \mid a \in A,\ b \in B\},
\]
l'insieme delle \emph{coppie ordinate}. Vale $(a, b) \neq (b, a)$ in
generale, e
\[
(x, y) = (x', y') \iff x = x' \text{ e } y = y'.
\]
\end{definizione}

\begin{esempio}
Siano $A = \{1, 2, 3\}$ e $B = \{v, w\}$. Allora
\begin{align*}
A \times B &= \{(1,v), (1,w), (2,v), (2,w), (3,v), (3,w)\}, \\
B \times A &= \{(v,1), (w,1), (v,2), (w,2), (v,3), (w,3)\},
\end{align*}
quindi $A \times B \neq B \times A$. Inoltre
\[
A \times A = A^2 = \{(1,1), (1,2), (1,3), (2,1), (2,2), (2,3), (3,1),
(3,2), (3,3)\}.
\]
\end{esempio}

Analogamente
\[
A \times B \times C = \{(a, b, c) \mid a \in A,\ b \in B,\ c \in C\},
\qquad A \times A \times A = A^3.
\]

\section[Relazioni di equivalenza (Lezione 2 -- 22 settembre 2026)]{Relazioni di equivalenza}
\noindent\emph{Lezione 2 -- 22 settembre 2026}
\subsection{Relazioni}

\begin{definizione}
Una \emph{relazione} $R$ (o $\rho$) su un insieme $A$ è un sottoinsieme del
prodotto cartesiano $A \times A$:
\[
R \subseteq A \times A.
\]
Se $(a, b) \in R$ si scrive $a \mathrel{R} b$.
\end{definizione}

\begin{esempio}
Sia $A = \{a, b, c\}$ e $R = \{(a,a), (a,b), (c,b)\}$. Si ha
$a \mathrel{R} b$, mentre $(b, a) \notin R$, cioè $b \not\mathrel{R} a$.
Questa relazione è transitiva, non è simmetrica e non è riflessiva
($b \not\mathrel{R} b$).
\end{esempio}

\begin{esempio}[Relazioni su $\mathbb{N}$ e $\mathbb{R}$]
\leavevmode
\begin{enumerate}
    \item Su $\mathbb{N}$ la relazione $\leq$: $(5, 8) \in\ \leq$ perché
          $5 \leq 8$, mentre $(10, 3) \notin\ \leq$ perché $10 \nleq 3$.
          È riflessiva ($a \leq a$) e transitiva
          ($a \leq b \land b \leq c \Rightarrow a \leq c$). Non è
          simmetrica: $3 \leq 5$ ma $5 \nleq 3$.

    \item Su $\mathbb{R}$ sia
          $R_1 = \{(x, y) \in \mathbb{R}^2 \mid x + 2y = 3\}
          \subseteq \mathbb{R} \times \mathbb{R}$.
          Si ha $(2, 2) \notin R_1$ perché $2 + 2 \cdot 2 = 6 \neq 3$;
          invece $1 \mathrel{R_1} 1$ (perché $1 + 2 = 3$), ma
          $2 \not\mathrel{R_1} 2$: $R_1$ \textbf{non è riflessiva}.
          Inoltre $(-1, 2) \in R_1$ perché $-1 + 2 \cdot 2 = 3$, ma
          $(2, -1) \notin R_1$ perché $2 + 2 \cdot (-1) = 0 \neq 3$:
          $R_1$ non è simmetrica.

    \item Su $\mathbb{N}$: $h \mathrel{R_2} k \iff h - k = 12$. Allora
          \[
          R_2 = \{(12, 0), (13, 1), (14, 2), \dots\}
              = \{(12 + k, k) \mid k \in \mathbb{N}\}.
          \]
          $(0, 0) \notin R_2$: non è riflessiva.

    \item Su $\mathbb{N}$: $m \mathrel{R_3} n \iff m + n = 5$. Allora
          \[
          R_3 = \{(0,5), (1,4), (2,3), (3,2), (4,1), (5,0)\}.
          \]
          Non è riflessiva ($0 \not\mathrel{R_3} 0$). È simmetrica.
          \textbf{Non} è transitiva: $1 \mathrel{R_3} 4$ e
          $4 \mathrel{R_3} 1$, ma $1 \not\mathrel{R_3} 1$.
\end{enumerate}
\end{esempio}

\begin{definizione}
Se $R$ è una relazione su $A$, la \emph{relazione opposta} $R^{op}$ si
ottiene scambiando di posto gli elementi delle coppie di $R$:
\[
(a, b) \in R \iff (b, a) \in R^{op}, \qquad
a \mathrel{R} b \iff b \mathrel{R^{op}} a.
\]
\end{definizione}

\begin{definizione}
Sia $R$ una relazione su $A$. Si dice che $R$ è:
\begin{itemize}
    \item \emph{riflessiva} se $\forall\, a \in A \quad a \mathrel{R} a$;
    \item \emph{simmetrica} se $a \mathrel{R} b \Rightarrow
          b \mathrel{R} a$;
    \item \emph{transitiva} se $a \mathrel{R} b \land b \mathrel{R} c
          \Rightarrow a \mathrel{R} c$.
\end{itemize}
\end{definizione}

\subsection{Relazioni di equivalenza}

\begin{definizione}
Una relazione $R$ su $A$ riflessiva, simmetrica e transitiva si dice
\emph{relazione di equivalenza}.
\end{definizione}

\begin{esempio}
\leavevmode
\begin{enumerate}
    \item Dato un insieme $A$, la relazione $R = A \times A$ è di
          equivalenza (tutti sono in relazione con tutti). Per ogni
          $a \in A$ si ha $[a] = A$, quindi $A/_R = \{[a]\} = \{A\}$.

    \item In $\mathbb{N}$ la relazione di uguaglianza,
          $n \mathrel{R} k \iff n = k$, è di equivalenza. Le classi sono
          $[n]_= = \{n\}$, ad esempio $[5]_= = \{5\}$, e
          $\mathbb{N}/_= = \{[0], [1], [2], \dots\}$.

    \item Sia $S$ l'insieme delle rette di un piano. Due rette
          $r, s \in S$ sono in relazione se $r$ è parallela a $s$,
          $r \parallel s$. È riflessiva ($s \parallel s$), simmetrica
          ($r \parallel s \Rightarrow s \parallel r$) e transitiva
          ($r \parallel s \land s \parallel t \Rightarrow r \parallel t$):
          è una relazione di equivalenza. La classe di $r$ è
          $[r]_{\parallel} = \{s \mid s \parallel r\}$, e se
          $r_1 \parallel r_2$ allora $[r_1]_{\parallel} =
          [r_2]_{\parallel}$.

    \item Sia $A = \{1, 2, 3, 4, 5, 6\}$ e
          \[
          \begin{split}
          R = \{&(1,1), (1,2), (1,6), (2,1), (2,6), (2,2), (6,6),\\
          &(6,1), (6,2), (3,3), (4,4), (4,5), (5,4), (5,5)\}.
          \end{split}
          \]
          È riflessiva e simmetrica. È transitiva: ad esempio
          $(1,6), (6,2) \in R \Rightarrow (1,2) \in R$ e
          $4 \mathrel{R} 5, 5 \mathrel{R} 4 \Rightarrow 4 \mathrel{R} 4$.
          Quindi è di equivalenza, con
          \[
          [1]_R = \{1, 2, 6\} = [2]_R = [6]_R, \quad [3]_R = \{3\},
          \quad [4]_R = [5]_R = \{4, 5\},
          \]
          e $A/_R = \{[1]_R, [3]_R, [4]_R\}$.
\end{enumerate}
\end{esempio}

\begin{esempio}[Relazioni che non sono di equivalenza]
\leavevmode
\begin{enumerate}
    \setcounter{enumi}{4}
    \item Su $\mathbb{Z}$: $k \mathrel{R} h \iff h + k = 10$. Non è
          riflessiva ($0 \not\mathrel{R} 0$), quindi non è di equivalenza.

    \item $A = \{a, b, c\}$ e $R = \{(a,a), (b,b), (a,c), (c,a)\}$.
          Non è riflessiva ($c \not\mathrel{R} c$). È simmetrica. Non è
          transitiva: $c \mathrel{R} a$ e $a \mathrel{R} c$, ma
          $c \not\mathrel{R} c$. (Altri casi, come
          $c \mathrel{R} a, a \mathrel{R} a \Rightarrow c \mathrel{R} a$,
          vanno bene: basta un caso che fallisce.)

    \item Su $\mathcal{P}(A)$ la relazione $X \subseteq Y$ (``essere
          sottoinsieme di'', con $X, Y \in \mathcal{P}(A)$) non è
          simmetrica.

    \item Su $\mathbb{N}$ la relazione $\leq$ non è di equivalenza: non è
          simmetrica.
\end{enumerate}
\end{esempio}

\subsection{Classi di equivalenza e insieme quotiente}

\begin{definizione}
Sia $A$ un insieme e $\rho$ una relazione di equivalenza su $A$. La
\emph{classe di equivalenza} di $a \in A$ è l'insieme
\[
[a]_\rho = \{x \in A \mid x \mathrel{\rho} a\}.
\]
L'insieme di tutte le classi di equivalenza è detto \emph{insieme
quotiente} di $A$ rispetto a $\rho$:
\[
A/_\rho = \{[a]_\rho \mid a \in A\}.
\]
\end{definizione}

\begin{osservazione}
Sia $\rho$ una relazione di equivalenza su $A$ e $a \in A$.
\begin{enumerate}
    \item $[a]_\rho \subseteq A$, quindi $[a]_\rho \in \mathcal{P}(A)$ e
          $A/_\rho = \{[a]_\rho \mid a \in A\} \subseteq \mathcal{P}(A)$.
    \item $[a]_\rho = [b]_\rho \iff a \mathrel{\rho} b \iff
          [a]_\rho \cap [b]_\rho \neq \emptyset$.
    \item Se $\rho$ è di equivalenza, $\rho^{op} = \rho$.
\end{enumerate}
\end{osservazione}

\subsection{Partizioni}

\begin{definizione}
Sia $A$ un insieme e $\mathcal{F} \subseteq \mathcal{P}(A)$ una famiglia di
sottoinsiemi di $A$. $\mathcal{F}$ è una \emph{partizione} di $A$ se:
\begin{enumerate}
    \item $\forall\, X \in \mathcal{F} \quad X \neq \emptyset$;
    \item $\forall\, X, Y \in \mathcal{F} \quad X \neq Y \Rightarrow
          X \cap Y = \emptyset$;
    \item $\displaystyle \bigcup_{X \in \mathcal{F}} X = A$.
\end{enumerate}
\end{definizione}

In parole: i pezzi non sono vuoti, sono a due a due disgiunti e ricoprono
tutto $A$.

\begin{esempio}
\leavevmode
\begin{enumerate}
    \item Se $A \neq \emptyset$, $\mathcal{F} = \{A\}$ è una partizione.
    \item $\mathcal{F} = \{\{x\} \mid x \in A\}$ è una partizione.
    \item $A = \{a, b, c, d, e\}$ e
          $\mathcal{F} = \{\{a, c, e\}, \{b\}, \{d\}\}$ è una partizione
          di $A$.
    \item In $\mathbb{Z}$ siano
          \begin{align*}
          X_0 &= \{3k \mid k \in \mathbb{Z}\} = \{0, \pm 3, \pm 6, \dots\},\\
          X_1 &= \{3k + 1 \mid k \in \mathbb{Z}\}
                = \{1, 4, -2, 7, -5, \dots\},\\
          X_2 &= \{3k + 2 \mid k \in \mathbb{Z}\}
                = \{2, 5, -1, 8, -4, \dots\}.
          \end{align*}
          Allora $\mathcal{F} = \{X_0, X_1, X_2\}$ è una partizione di
          $\mathbb{Z}$.
\end{enumerate}
\end{esempio}

\begin{proposizione}[A]
Sia $\rho$ una relazione di equivalenza su $A$. Allora l'insieme quotiente
$A/_\rho$ è una partizione di $A$.
\end{proposizione}

\begin{proof}
Verifichiamo le proprietà 1, 2 e 3 della definizione di partizione.
\begin{enumerate}
    \item $[a]_\rho \neq \emptyset$: infatti $a \mathrel{\rho} a$, cioè
          $a \in [a]_\rho$.
    \item Se $[a]_\rho \neq [b]_\rho$, vogliamo $[a]_\rho \cap [b]_\rho =
          \emptyset$. Sia, per assurdo, $x \in [a]_\rho \cap [b]_\rho$:
          allora $x \mathrel{\rho} a$ e $x \mathrel{\rho} b$. Per la
          simmetria $a \mathrel{\rho} x$ e $x \mathrel{\rho} b$, e per la
          transitività $a \mathrel{\rho} b$, cioè $a \in [b]_\rho$ e
          $b \in [a]_\rho$. Quindi $[a]_\rho = [b]_\rho$ (hanno gli stessi
          elementi), contro l'ipotesi. Dunque se $[a]_\rho \neq
          [b]_\rho$ non esiste $x \in [a]_\rho \cap [b]_\rho$, cioè
          $[a]_\rho \cap [b]_\rho = \emptyset$.
    \item $\displaystyle \bigcup_{[a]_\rho \in A/_\rho} [a]_\rho = A$:
          se $a \in A$ allora $a \in [a]_\rho$, cioè le classi
          ricoprono tutto $A$. \qedhere
\end{enumerate}
\end{proof}

Viceversa, da una partizione si ricava una relazione di equivalenza.

\begin{definizione}
Sia $\mathcal{F}$ una partizione di $A$. Si definisce la relazione
$R_\mathcal{F}$ su $A$ ponendo
\[
a \mathrel{R_\mathcal{F}} b \iff \exists\, X \in \mathcal{F} \text{ tale
che } a, b \in X.
\]
Cioè $a$ e $b$ sono in relazione se stanno nello stesso pezzo della
partizione.
\end{definizione}

\begin{proposizione}[B]
Se $\mathcal{F}$ è una partizione su $A$, la relazione $R_\mathcal{F}$ è
una relazione di equivalenza su $A$.
\end{proposizione}

\begin{proof}
Verifichiamo che $R_\mathcal{F}$ è riflessiva, simmetrica e transitiva.
\begin{enumerate}
    \item \emph{Riflessiva}: $a \mathrel{R_\mathcal{F}} a$? Poiché
          $\mathcal{F}$ è una partizione, per la proprietà 3
          $\bigcup_{X \in \mathcal{F}} X = A$. Se $a \in A$, allora
          $a \in \bigcup_{X \in \mathcal{F}} X$, cioè esiste
          $X \in \mathcal{F}$ tale che $a \in X$. Quindi
          $a \mathrel{R_\mathcal{F}} a$.
    \item \emph{Simmetrica}: se $a \mathrel{R_\mathcal{F}} b$, cioè esiste
          $X \in \mathcal{F}$ con $a, b \in X$, allora anche
          $b, a \in X$, cioè $b \mathrel{R_\mathcal{F}} a$.
    \item \emph{Transitiva}: siano $a \mathrel{R_\mathcal{F}} b$ e
          $b \mathrel{R_\mathcal{F}} c$. Esistono $X \in \mathcal{F}$ con
          $a, b \in X$ e $Y \in \mathcal{F}$ con $b, c \in Y$. Poiché
          $b \in X \cap Y \neq \emptyset$, ma per la proprietà 2 gli
          elementi della partizione sono disgiunti, deve essere $X = Y$.
          Quindi $a, b, c \in X = Y$, cioè $a, c \in X$, cioè
          $a \mathrel{R_\mathcal{F}} c$. \qedhere
\end{enumerate}
\end{proof}

\begin{teorema}[Fondamentale delle relazioni di equivalenza]
Esiste una corrispondenza 1 a 1 tra l'insieme di tutte le relazioni di
equivalenza su $A$ e l'insieme di tutte le partizioni di $A$.
\end{teorema}

\begin{proof}
Se $\rho$ è una relazione di equivalenza su $A$, allora $A/_\rho$ è una
partizione (Proposizione A). Se $\mathcal{F}$ è una partizione di $A$,
$R_\mathcal{F}$ è una relazione di equivalenza su $A$ (Proposizione B).
\end{proof}

\begin{osservazione}
Le due costruzioni sono l'una l'inversa dell'altra: la classe di $a$
rispetto a $R_\mathcal{F}$ è esattamente l'elemento di $\mathcal{F}$ che
contiene $a$, quindi $A/_{R_\mathcal{F}} = \mathcal{F}$; e due elementi
stanno nella stessa classe di $A/_\rho$ se e solo se sono in relazione
rispetto a $\rho$, quindi $R_{A/_\rho} = \rho$. Questo è ciò che rende la
corrispondenza biunivoca.
\end{osservazione}

\section[Congruenza modulo $n$ (Lezione 3 -- 28 settembre 2026)]{Congruenza modulo $n$}
\noindent\emph{Lezione 3 -- 28 settembre 2026}
\subsection{La congruenza modulo $n$}

\begin{definizione}
Siano $a, b \in \mathbb{Z}$. Si dice che $a$ \emph{divide} $b$, e si scrive
$a \mid b$, se
\[
\exists\, k \in \mathbb{Z} \text{ tale che } b = ka.
\]
Se non esiste un tale $k$ si scrive $a \nmid b$ (``$a$ non divide $b$'').
\end{definizione}

Sia $n \geq 1$ un intero fissato.

\begin{definizione}
La \emph{congruenza modulo $n$} è la relazione $\equiv_n$ su $\mathbb{Z}$
definita da
\[
a \equiv_n b \iff n \mid (b - a) \iff \exists\, k \in \mathbb{Z} \text{
tale che } b - a = kn \quad (b = kn + a).
\]
Si scrive anche $a \equiv b \pmod n$.
\end{definizione}

\begin{proposizione}
$\equiv_n$ è una relazione di equivalenza su $\mathbb{Z}$, e
\[
\mathbb{Z}/_{\equiv_n} = \mathbb{Z}_n = \{[0]_n, [1]_n, \dots, [n-1]_n\}.
\]
\end{proposizione}

\begin{proof}
Verifichiamo che $\equiv_n$ è una relazione di equivalenza.
\begin{enumerate}
    \item \emph{Riflessiva}: $a \equiv_n a$, perché $n \mid (a - a) = 0$,
          cioè $0 = k \cdot n$ con $k = 0$.
    \item \emph{Simmetrica}: se $a \equiv_n b$, allora $n \mid (b - a)$,
          cioè esiste $k \in \mathbb{Z}$ con $b - a = kn$. Posto
          $k' = -k$ si ha $a - b = -kn = k'n$, quindi $n \mid (a - b)$,
          cioè $b \equiv_n a$.
    \item \emph{Transitiva}: siano $a \equiv_n b$ e $b \equiv_n c$.
          Allora $b - a = kn$ e $c - b = hn$ per certi $k, h \in
          \mathbb{Z}$. Sommando,
          \[
          c - a = kn + hn = (k + h)n,
          \]
          quindi $n \mid (c - a)$, cioè $a \equiv_n c$. \qedhere
\end{enumerate}
\end{proof}

La classe di $a$ è
\begin{align*}
[a]_n &= \{b \in \mathbb{Z} \mid a \equiv_n b\}
       = \{b \in \mathbb{Z} \mid n \mid (b - a)\} \\
      &= \{b \in \mathbb{Z} \mid \exists\, k \in \mathbb{Z} \text{ tale
         che } b = kn + a\}
       = \{kn + a \mid k \in \mathbb{Z}\},
\end{align*}
cioè è l'insieme di tutti i numeri che differiscono da $a$ per un multiplo
di $n$. L'insieme quotiente è
$\mathbb{Z}_n = \{[0]_n, [1]_n, \dots, [n-1]_n\}$.

\begin{esempio}
Per $n = 3$ le classi $[0]_3, [1]_3, [2]_3$ sono gli insiemi $X_0, X_1,
X_2$ dell'esempio delle partizioni di $\mathbb{Z}$ (lezione 2): $\mathbb{Z}_3$ è proprio la partizione
di $\mathbb{Z}$ in multipli di $3$, multipli di $3$ più $1$ e multipli di
$3$ più $2$.
\end{esempio}

\section[Relazioni d'ordine (Lezione 3 -- 28 settembre 2026)]{Relazioni d'ordine}
\noindent\emph{Lezione 3 -- 28 settembre 2026}
\subsection{Ordini parziali e totali}

\begin{definizione}
Sia $A$ un insieme e $R$ una relazione su $A$. $R$ è
\emph{antisimmetrica} se
\[
a \mathrel{R} b \ \land\ b \mathrel{R} a \implies a = b.
\]
\end{definizione}

\begin{definizione}
Una relazione $R$ (che si indica con $\preceq$) su $A$ è una
\emph{relazione d'ordine parziale} se è riflessiva, antisimmetrica e
transitiva.
\end{definizione}

\begin{itemize}
    \item $A$ con una relazione d'ordine si dice \emph{parzialmente
          ordinato}.
    \item Si scrive $a \preceq b$; con $a \prec b$ si intende
          $a \preceq b$ e $a \neq b$.
    \item Una relazione d'ordine $\preceq$ è \emph{totale} se per ogni
          $a, b \in A$ vale $a \preceq b$ oppure $b \preceq a$.
\end{itemize}
$A$ con una relazione d'ordine totale si dice \emph{totalmente ordinato}
o \emph{catena}.

\begin{esempio}
\leavevmode
\begin{enumerate}
    \item $(\mathbb{N}, \leq)$ è riflessivo ($n \leq n$), antisimmetrico
          ($n \leq k \land k \leq n \Rightarrow n = k$) e transitivo
          ($n \leq k \land k \leq h \Rightarrow n \leq h$). È una
          relazione d'ordine \textbf{totale}: per ogni $n, k \in
          \mathbb{N}$ vale $n \leq k$ oppure $k \leq n$.

    \item $(\mathbb{Z}, \leq)$ è una relazione d'ordine totale.

    \item Sia $B = \{x, y, z\}$. Allora
          \[
          \mathcal{P}(B) = \big\{\emptyset, \{x\}, \{y\}, \{z\},
          \{x,y\}, \{x,z\}, \{y,z\}, B\big\}.
          \]
          $(\mathcal{P}(B), \subseteq)$ è una relazione d'ordine
          \textbf{parziale}. Non è totale: $\{x, y\}$ e $\{y, z\}$ non
          sono in relazione.
\end{enumerate}
\end{esempio}

\subsection{Il diagramma di Hasse}

\begin{osservazione}
Ad ogni relazione d'ordine si può associare un \emph{diagramma di Hasse}:
\begin{itemize}
    \item i nodi sono esattamente gli elementi dell'insieme;
    \item se $x \preceq y$ (con $x \neq y$) il nodo $x$ è disegnato più in
          basso del nodo $y$;
    \item due nodi $x, y$ sono collegati da un arco se \emph{non} esiste
          $z$ tale che $x \prec z$ e $z \prec y$ (cioè se $y$ ``viene
          subito dopo'' $x$).
\end{itemize}
\end{osservazione}

\begin{center}
\begin{tikzpicture}[scale=1.1, every node/.style={font=\small}]
  % catena N
  \begin{scope}
    \foreach \n/\y in {0/0, 1/1, 2/2, 3/3} {
      \fill (0,\y) circle (1.5pt) node[right=2pt] {$\n$};
    }
    \draw (0,0) -- (0,3.4);
    \node at (0,-0.7) {$(\mathbb{N}, \leq)$};
  \end{scope}

  % P(B)
  \begin{scope}[xshift=6cm]
    \node (e) at (0,0) {$\emptyset$};
    \node (x) at (-2,1.4) {$\{x\}$};
    \node (y) at (0,1.4) {$\{y\}$};
    \node (z) at (2,1.4) {$\{z\}$};
    \node (xy) at (-2,2.9) {$\{x,y\}$};
    \node (xz) at (0,2.9) {$\{x,z\}$};
    \node (yz) at (2,2.9) {$\{y,z\}$};
    \node (b) at (0,4.3) {$B$};
    \foreach \a/\b in {e/x, e/y, e/z, x/xy, x/xz, y/xy, y/yz, z/xz, z/yz,
                       xy/b, xz/b, yz/b}
      \draw (\a) -- (\b);
    \node at (0,-0.9) {$(\mathcal{P}(B), \subseteq)$};
  \end{scope}
\end{tikzpicture}
\end{center}

\begin{esempio}[continua]
\leavevmode
\begin{enumerate}
    \setcounter{enumi}{3}
    \item Sia $X = \{1, 2, 3, 4, 5, 6\}$ con la relazione ``$\mid$''
          (divide). Le coppie in relazione sono
          \begin{center}
          \begin{tabular}{llllll}
          $1 \mid 1$ & $1 \mid 2$ & $1 \mid 3$ & $1 \mid 4$ & $1 \mid 5$ &
          $1 \mid 6$ \\
                     & $2 \mid 2$ & $2 \mid 4$ & $2 \mid 6$ \\
                     & $3 \mid 3$ & $3 \mid 6$ \\
                     & $4 \mid 4$ & $5 \mid 5$ & $6 \mid 6$
          \end{tabular}
          \end{center}
          È una relazione d'ordine parziale, non totale: $4 \nmid 5$ e
          $5 \nmid 4$.

\begin{center}
\begin{tikzpicture}[scale=1.1, every node/.style={font=\small}]
  \node (1) at (1.5,0) {$1$};
  \node (2) at (0,1.4) {$2$};
  \node (3) at (1.5,1.4) {$3$};
  \node (5) at (3,1.4) {$5$};
  \node (4) at (0,2.8) {$4$};
  \node (6) at (1.5,2.8) {$6$};
  \foreach \a/\b in {1/2, 1/3, 1/5, 2/4, 2/6, 3/6}
    \draw (\a) -- (\b);
\end{tikzpicture}
\end{center}

    \item $(\mathbb{N} \setminus \{0\}, \mid)$ è una relazione d'ordine non
          totale.

    \item $\mathbb{Z} \setminus \{0\}$ con ``$\mid$'' (divide) è una
          relazione d'ordine? \textbf{No}: non è antisimmetrica, perché
          $5 \mid -5$ e $-5 \mid 5$ ma $-5 \neq 5$.
\end{enumerate}
\end{esempio}

\begin{osservazione}
La relazione opposta $R^{op}$ scambia le coppie. Se $\preceq$ è una
relazione d'ordine, allora $(\preceq)^{op} = \succeq$ è ancora una relazione
d'ordine: il suo diagramma di Hasse si ottiene capovolgendo quello di
$\preceq$. Ad esempio $(\leq)^{op} = \ \geq$ su $\mathbb{N}$, e
$(\subseteq)^{op} = \ \supseteq$ su $\mathcal{P}(B)$ (con $B$ in basso e
$\emptyset$ in alto).
\end{osservazione}

\subsection{Minimo, massimo, minimale, massimale}

\begin{definizione}
Sia $(A, \preceq)$ un insieme parzialmente ordinato. Un elemento $t \in A$ è:
\begin{itemize}
    \item \emph{minimo} se $\forall\, a \in A \quad t \preceq a$;
    \item \emph{massimo} se $\forall\, a \in A \quad a \preceq t$;
    \item \emph{minimale} se non ci sono elementi minori di $t$, cioè
          $\nexists\, a \in A$ tale che $a \prec t$;
    \item \emph{massimale} se non ci sono elementi maggiori di $t$, cioè
          $\nexists\, a \in A$ tale che $t \prec a$.
\end{itemize}
\end{definizione}

Un minimo è sempre minimale (e un massimo è sempre massimale), ma il
viceversa in generale è falso se l'ordine non è totale.

\begin{esempio}
Negli esempi precedenti:
\begin{itemize}
    \item $(\mathbb{N}, \leq)$: $0$ è minimo e minimale; non esistono
          massimo né massimali.
    \item $(\mathbb{Z}, \leq)$: non esistono minimo, massimo, minimali né
          massimali.
    \item $(\mathcal{P}(B), \subseteq)$: $\emptyset$ è minimo e minimale,
          $B$ è massimo e massimale.
    \item $(\{1, \dots, 6\}, \mid)$: $1$ è minimo e minimale; non esiste
          il massimo; $4$, $5$, $6$ sono massimali.
    \item $(\mathbb{N} \setminus \{0\}, \mid)$: $1$ è minimo e minimale;
          non esistono massimo né massimali.
\end{itemize}
\end{esempio}

\subsection{Maggioranti, minoranti, estremo superiore e inferiore}

\begin{definizione}
Sia $\preceq$ una relazione d'ordine su $A$ e sia $B \subseteq A$ un
sottoinsieme. Un elemento $w \in A$ è:
\begin{itemize}
    \item un \emph{maggiorante} di $B$ se $b \preceq w$ per ogni
          $b \in B$;
    \item un \emph{minorante} di $B$ se $w \preceq b$ per ogni $b \in B$.
\end{itemize}
\end{definizione}

\begin{esempio}
\leavevmode
\begin{enumerate}
    \item In $(\mathbb{N}, \leq)$ sia $B = \{6, 9, 12\}$.
          I minoranti di $B$ sono $\{0, 1, 2, \dots, 6\}$, i maggioranti
          sono $\{12, 13, 14, \dots\}$. Il più grande dei minoranti è
          $6 = \inf(B)$, il più piccolo dei maggioranti è
          $12 = \sup(B)$.

    \item In $(\mathbb{Z}, \leq)$ con $B = \{6, 9, 12\}$: i minoranti sono
          $\{\dots,\allowbreak -1,\allowbreak 0,\allowbreak 1,\allowbreak \dots,\allowbreak 6\}$, i maggioranti
          $\{12, 13, \dots\}$; $6 = \inf(B)$ e $12 = \sup(B)$.

    \item In $(\mathbb{N} \setminus \{0\}, \mid)$ con $B = \{6, 9, 12\}$:
          i minoranti sono i divisori comuni, $\{1, 3\}$; poiché $1 \mid 3$,
          il più grande (rispetto a $\mid$) è $3 = \inf(B)$. I maggioranti
          sono i multipli comuni, $\{36, 72, \dots\}$, cioè i multipli di
          $36$; $36$ divide tutti gli altri, quindi $36 = \sup(B)$.

    \item In $(\mathbb{Q}, \leq)$ sia
          $B = \{x \in \mathbb{Q} \mid \sqrt{2} < x < \sqrt{3}\}$.
          I minoranti sono tutti i razionali $q < \sqrt{2}$, i maggioranti
          tutti i razionali $q > \sqrt{3}$. Tra i minoranti non c'è un
          massimo e tra i maggioranti non c'è un minimo: in $\mathbb{Q}$
          non esistono né $\inf(B)$ né $\sup(B)$.
\end{enumerate}
\end{esempio}

\begin{definizione}
Sia $(A, \preceq)$ parzialmente ordinato e $B \subseteq A$.
\begin{itemize}
    \item Un elemento $t \in A$ è l'\emph{estremo inferiore} di $B$,
          $t = \inf(B)$, se $t$ è il \textbf{massimo} tra i minoranti di
          $B$, cioè
          \[
          t = \inf(B) \iff
          \begin{cases}
            t \preceq b & \forall\, b \in B, \\
            w \preceq b \ \ \forall\, b \in B \implies w \preceq t.
          \end{cases}
          \]
    \item Un elemento $t \in A$ è l'\emph{estremo superiore} di $B$,
          $t = \sup(B)$, se $t$ è il \textbf{minimo} tra i maggioranti di
          $B$, cioè
          \[
          t = \sup(B) \iff
          \begin{cases}
            b \preceq t & \forall\, b \in B, \\
            b \preceq w \ \ \forall\, b \in B \implies t \preceq w.
          \end{cases}
          \]
\end{itemize}
\end{definizione}

\section[Funzioni (Lezione 3 -- 28 settembre 2026)]{Funzioni}
\noindent\emph{Lezione 3 -- 28 settembre 2026}
\subsection{Definizione di funzione}

\begin{definizione}
Siano $A$ e $B$ due insiemi. Una \emph{funzione} $f$ da $A$ a $B$ è un
sottoinsieme del prodotto cartesiano $A \times B$ tale che
\[
\forall\, a \in A \quad \exists!\, b \in B : (a, b) \in f.
\]
\end{definizione}

In altre parole, ogni elemento di $A$ compare come prima componente in
\textbf{una e una sola} coppia di $f$. Si usano le notazioni seguenti:
\begin{itemize}
    \item $f : A \to B$ indica che $f$ è una funzione da $A$ a $B$;
    \item $a \mathrel{f} b$ oppure $f(a) = b$ significa $(a, b) \in f$;
          $b$ si dice \emph{immagine} di $a$;
    \item $A$ è il \emph{dominio} di $f$ e $B$ è il \emph{codominio} di $f$.
\end{itemize}

Il sottoinsieme di $A \times B$
\[
f = \{\, (a, f(a)) \mid a \in A,\ f(a) \in B \,\} \subseteq A \times B
\]
si chiama \emph{grafico} della funzione $f$.

\begin{osservazione}
Secondo la definizione, la funzione \emph{coincide} con il suo grafico.
Le due condizioni che la definiscono sono:
\begin{enumerate}
    \item \textbf{esistenza}: ogni $a \in A$ ha almeno un'immagine in $B$
          (nessun elemento del dominio resta ``senza freccia'');
    \item \textbf{unicità}: ogni $a \in A$ ha \emph{una sola} immagine
          (da ogni elemento del dominio parte una sola freccia).
\end{enumerate}
Invece \emph{non} si chiede che ogni elemento di $B$ sia immagine di
qualcuno, né che sia immagine di uno solo: nell'esempio~2 qui sotto $a$ e
$b$ hanno ciascuno due elementi che arrivano su di loro.
\end{osservazione}

\subsection{Esempi}

\begin{esempio}
\leavevmode
\begin{enumerate}
    \item Siano $A = \{x, y, z\}$ e $B = \{1, 2, 3, 4\}$ e sia
          $f : A \to B$ data da
          \[
          f(x) = 1, \qquad f(y) = 3, \qquad f(z) = 4.
          \]
          Quindi $x \mathrel{f} 1$, cioè $(x, 1) \in f$, e il grafico è
          $f = \{(x,1), (y,3), (z,4)\}$. L'elemento $2 \in B$ non è
          immagine di nessun elemento di $A$.

    \item Siano $B = \{1, 2, 3, 4\}$ e $C = \{a, b\}$ e sia
          $g : B \to C$ data da
          \[
          g(1) = b, \quad g(2) = b, \quad g(3) = a, \quad g(4) = a.
          \]
\end{enumerate}

\begin{center}
\begin{tikzpicture}[scale=0.9, every node/.style={font=\small}]
  % funzione f : A -> B
  \draw (0,0) ellipse (0.9 and 1.4);
  \draw (3,0) ellipse (0.9 and 1.7);
  \node at (0,1.75) {$A$};
  \node at (3,2.05) {$B$};
  \node at (1.5,2.0) {$f$};
  \foreach \n/\y in {x/0.8, y/0, z/-0.8} {
    \fill (-0.3,\y) circle (1.5pt) node[left] {$\n$};
  }
  \foreach \n/\y in {1/1.1, 2/0.35, 3/-0.4, 4/-1.15} {
    \fill (2.7,\y) circle (1.5pt) node[right] {$\n$};
  }
  \draw[->] (-0.3,0.8) -- (2.7,1.1);
  \draw[->] (-0.3,0) -- (2.7,-0.4);
  \draw[->] (-0.3,-0.8) -- (2.7,-1.15);

  % funzione g : B -> C
  \begin{scope}[xshift=8cm]
    \draw (0,0) ellipse (0.9 and 1.7);
    \draw (3,0) ellipse (0.9 and 1.0);
    \node at (0,2.05) {$B$};
    \node at (3,1.35) {$C$};
    \node at (1.5,2.0) {$g$};
    \foreach \n/\y in {1/1.1, 2/0.35, 3/-0.4, 4/-1.15} {
      \fill (-0.3,\y) circle (1.5pt) node[left] {$\n$};
    }
    \fill (2.7,0.5) circle (1.5pt) node[right] {$b$};
    \fill (2.7,-0.5) circle (1.5pt) node[right] {$a$};
    \draw[->] (-0.3,1.1) -- (2.7,0.5);
    \draw[->] (-0.3,0.35) -- (2.7,0.5);
    \draw[->] (-0.3,-0.4) -- (2.7,-0.5);
    \draw[->] (-0.3,-1.15) -- (2.7,-0.5);
  \end{scope}
\end{tikzpicture}
\end{center}

\begin{enumerate}
    \setcounter{enumi}{2}
    \item $f : \mathbb{N} \to \mathbb{N}$, \quad $f(n) = 2n$.
          Si scrive anche $n \mapsto f(n) = 2n$.

    \item $f : \mathbb{Z} \to \mathbb{N}$, \quad
          $k \mapsto |k| =
          \begin{cases}
            k & \text{se } k \geq 0, \\
            -k & \text{se } k < 0.
          \end{cases}$

    \item $h : \mathbb{R} \to \mathbb{R}$, \quad $h(x) = x^2 - 2$.

    \item $L : \mathbb{R} \to \mathbb{R}$, \quad $L(x) = 5x + 3$.

    \item $F : \mathbb{Z} \to \mathbb{Z}$, \quad $F(k) = 3$.

    \item $F : \mathbb{Z} \to \{3\}$, \quad $F(k) = 3$.
\end{enumerate}
\end{esempio}

Gli ultimi due esempi sono funzioni \emph{costanti}: tutti gli elementi del
dominio hanno la stessa immagine. Notiamo che la legge $F(k) = 3$ è la
stessa, ma le due funzioni hanno codomini diversi ($\mathbb{Z}$ nel primo
caso, $\{3\}$ nel secondo).

\begin{definizione}
Una funzione $f : A \to B$ tale che
\[
f(x_1) = f(x_2) \qquad \forall\, x_1, x_2 \in A
\]
si dice \emph{funzione costante}.
\end{definizione}

\section[Funzioni: la relazione associata (Lezione 4 -- 29 settembre 2026)]{Funzioni: la relazione associata}
\noindent\emph{Lezione 4 -- 29 settembre 2026}
Si riprendono le funzioni introdotte nella lezione 3 (definizione, dominio, codominio, grafico, funzioni costanti).

\subsection{La relazione associata a una funzione}

\begin{definizione}
Sia $f : A \to B$ una funzione. La \emph{relazione associata a $f$}, indicata
con $R_f$, è la relazione su $A$ definita da
\[
a_1 \mathrel{R_f} a_2 \iff f(a_1) = f(a_2).
\]
\end{definizione}

Due elementi di $A$ sono quindi in relazione se $f$ associa loro la stessa
immagine.

\begin{esempio}
Siano $A = \{1, 2, 3, 4\}$ e $B = \{x, y, z\}$ e sia $f : A \to B$ data da
\[
f(1) = f(2) = f(3) = x, \qquad f(4) = y.
\]
Allora $1 \mathrel{R_f} 2$, $3 \mathrel{R_f} 3$, $4 \mathrel{R_f} 4$, \dots{}
Le classi di equivalenza sono
\[
[1]_{R_f} = \{1, 2, 3\} = [2]_{R_f} = [3]_{R_f}, \qquad
[4]_{R_f} = \{4\},
\]
e l'insieme quotiente è $A/_{R_f} = \{[1]_{R_f}, [4]_{R_f}\}$.
\end{esempio}

\begin{teorema}
Sia $f : A \to B$ una funzione. La relazione $R_f$ è una relazione di
equivalenza su $A$.
\end{teorema}

\begin{proof}
Lasciata per esercizio a lezione; la verifica è
immediata, perché ``avere la stessa immagine'' eredita le proprietà
dell'uguaglianza in $B$.
\begin{enumerate}
    \item \emph{Riflessiva}: per ogni $a \in A$ vale $f(a) = f(a)$, quindi
          $a \mathrel{R_f} a$.
    \item \emph{Simmetrica}: se $a_1 \mathrel{R_f} a_2$, cioè
          $f(a_1) = f(a_2)$, allora $f(a_2) = f(a_1)$, cioè
          $a_2 \mathrel{R_f} a_1$.
    \item \emph{Transitiva}: se $a_1 \mathrel{R_f} a_2$ e
          $a_2 \mathrel{R_f} a_3$, allora $f(a_1) = f(a_2)$ e
          $f(a_2) = f(a_3)$, quindi $f(a_1) = f(a_3)$, cioè
          $a_1 \mathrel{R_f} a_3$. \qedhere
\end{enumerate}
\end{proof}

\section[Principio di induzione (Lezione 4 -- 29 settembre 2026)]{Principio di induzione}
\noindent\emph{Lezione 4 -- 29 settembre 2026}
\subsection{Il principio di induzione}

\begin{osservazione}
Sia $U \subseteq \mathbb{N}$ tale che
\begin{enumerate}
    \item $0 \in U$;
    \item $k \in U \Rightarrow k + 1 \in U$.
\end{enumerate}
Allora $U = \mathbb{N}$.
\end{osservazione}

Infatti $0 \in U$, quindi $1 \in U$, quindi $2 \in U$, e così via: ``si sale
a gradini'' e si raggiungono tutti i naturali. Se $p(n)$ è un predicato
sui naturali, si applica questa idea all'insieme $U$ dei $n$ per cui
$p(n)$ è vero.

\begin{teorema}[Principio di induzione]
Sia $p(n)$ un predicato. Supponiamo che:
\begin{enumerate}
    \item $p(0)$ sia vero \hfill [\emph{base dell'induzione}];
    \item $\forall\, k \ \big(p(k) \text{ vero} \Rightarrow p(k+1)
          \text{ vero}\big)$ \hfill [\emph{passo induttivo}].
\end{enumerate}
Allora $\forall\, n \in \mathbb{N}$ $p(n)$ è vero.
\end{teorema}

Nel passo induttivo l'ipotesi $p(k)$ si chiama \emph{ipotesi induttiva} e
$p(k+1)$ è la \emph{tesi da dimostrare}.

\begin{esempio}
Dimostrare che per ogni $n$ naturale $2 \mid n(n+1)$.

Sia $p(n)$ il predicato ``$2 \mid n(n+1)$''.

\textbf{1) Caso base.} $p(0)$: $2 \mid 0 \cdot (0 + 1) = 0$. Infatti
$2 \mid 0$ significa $0 = k \cdot 2$ con $k \in \mathbb{Z}$, ed è vero per
$k = 0$: $0 = 0 \cdot 2$.

\textbf{2) Passo induttivo.} Dobbiamo dimostrare $p(k) \Rightarrow
p(k+1)$.
\begin{itemize}
    \item \emph{Ipotesi induttiva} $p(k)$: $2 \mid k(k+1)$.
    \item \emph{Tesi da dimostrare} $p(k+1)$: $2 \mid (k+1)\big((k+1)+1\big)$.
\end{itemize}
Calcoliamo
\[
(k+1)(k+2) = (k+1)k + 2(k+1).
\]
Per l'ipotesi induttiva $2 \mid k(k+1)$, cioè $k(k+1) = 2m$ con
$m \in \mathbb{Z}$. Quindi
\[
(k+1)(k+2) = 2m + 2(k+1) = 2\,(m + k + 1),
\]
cioè $(k+1)(k+2)$ è multiplo di $2$, ovvero $2 \mid (k+1)(k+2)$, che è la
tesi da dimostrare.
\end{esempio}

\subsection{Induzione a partire da $\bar n$}

Se il predicato è vero solo da un certo naturale in poi, si parte da
quel punto.

\begin{teorema}
Sia $p(n)$ una proposizione e sia $\bar n \in \mathbb{N}$. Supponiamo che:
\begin{enumerate}
    \item $p(\bar n)$ sia vero;
    \item $\forall\, k \geq \bar n \ \big(p(k) \text{ vero} \Rightarrow
          p(k+1) \text{ vero}\big)$.
\end{enumerate}
Allora $\forall\, n \geq \bar n$ $p(n)$ è vera.
\end{teorema}

\begin{esempio}
Il predicato $p(n)$: ``$n - 5 > 0$'' è vero da $6$ in poi: si applica il
teorema con $\bar n = 6$.
\end{esempio}

\begin{esempio}
Dimostrare che per ogni $n \geq 1$
\[
1 + 2 + \dots + n = \frac{n(n+1)}{2}. \tag*{$p(n)$}
\]

\textbf{1) Caso base} $\bar n = 1$. $p(1)$: $1 = \dfrac{1 \cdot (1+1)}{2}$,
vera.

\textbf{2) Passo induttivo.}
\begin{itemize}
    \item \emph{Ipotesi induttiva} $p(k)$:
          $1 + \dots + k = \dfrac{k(k+1)}{2}$.
    \item \emph{Tesi da dimostrare} $p(k+1)$:
          $1 + \dots + k + (k+1) = \dfrac{(k+1)(k+2)}{2}$.
\end{itemize}
Si ha
\begin{align*}
1 + 2 + \dots + k + (k+1) &= (1 + \dots + k) + (k+1)\\
&= \frac{k(k+1)}{2} + (k+1) \qquad \text{(per ipotesi induttiva)}\\
&= \frac{k(k+1) + 2(k+1)}{2} = \frac{(k+1)(k+2)}{2},
\end{align*}
che è la tesi. Per induzione $p(n)$ vale per ogni $n \geq 1$.
\end{esempio}

\section[Successioni (Lezione 4 -- 29 settembre 2026)]{Successioni}
\noindent\emph{Lezione 4 -- 29 settembre 2026}
\subsection{Successioni}

\begin{definizione}
Una \emph{successione} è una funzione $f : \mathbb{N} \to X$ (oppure con
dominio $\mathbb{N} \setminus \{0\}$), dove $X$ è un insieme qualsiasi.
\end{definizione}

L'immagine $f(n)$ di $n$ si indica con $x_n \in X$ e si chiama anche
l'\emph{$n$-esimo termine} della successione. L'insieme immagine di una
successione è
\[
\operatorname{Im} f = f(\mathbb{N}) = \{x_n\}_{n \in \mathbb{N}}
\subseteq X.
\]

\begin{esempio}
\leavevmode
\begin{enumerate}
    \item $a : \mathbb{N} \setminus \{0\} \to \mathbb{Q}$,
          $n \mapsto \dfrac{1}{n}$. Si scrive $a_n = \dfrac{1}{n} = a(n)$
          e l'insieme dei termini è
          $\{a_n\} = \left\{\dfrac{1}{n}\right\}_{n \in \mathbb{N}
          \setminus \{0\}}$.

    \item $b : \mathbb{N} \to \mathbb{N}$, $n \mapsto 2n + 3$. I primi
          termini sono $b_0 = 2 \cdot 0 + 3 = 3$,
          $b_1 = 2 \cdot 1 + 3 = 5$, \dots

    \item $c : \mathbb{N} \to \mathbb{N}$, $n \mapsto 2^n$. Ad esempio
          $c_3 = 2^3 = 8$, \dots
\end{enumerate}
\end{esempio}

\section[Sommatorie, ricorsione e polinomi (Lezione 5 -- 5 ottobre 2026)]{Sommatorie, ricorsione e polinomi}
\noindent\emph{Lezione 5 -- 5 ottobre 2026}
\subsection{Sommatorie}

Sia $f \colon \mathbb{N} \to \mathbb{R}$ una successione, $f(n) = a_n$. Ad
esempio, per $a_n = n$ si ha $a_1 = 1$, $a_2 = 2$, \dots, $a_{100} = 100$.

Sommare i primi due termini significa calcolare $a_1 + a_2 = 1 + 2$; i primi
cinque, $a_1 + a_2 + a_3 + a_4 + a_5 = 1 + 2 + 3 + 4 + 5$. Per scrivere somme
con molti termini si usa il simbolo di \emph{sommatoria}:
\[
\sum_{i=1}^{n} f(i) = f(1) + f(2) + f(3) + \dots + f(n).
\]

\begin{esempio}
La somma dei primi $100$ numeri è
\[
1 + 2 + \dots + 100 = \sum_{i=1}^{100} a_i = \sum_{i=1}^{100} i .
\]
\end{esempio}

\begin{definizione}
Sia $a_n \colon \mathbb{N} \to \mathbb{R}$ una successione. Le si associa
un'altra successione
\[
s_k = a_1 + a_2 + \dots + a_k = \sum_{i=1}^{k} a_i .
\]
La successione $\{s_k\}$ si dice \emph{serie associata} alla successione
$\{a_n\}$.
\end{definizione}

\subsubsection{Proprietà della sommatoria}

\begin{proposizione}
Siano $\{a_n\}$, $\{b_n\}$ due successioni in $\mathbb{R}$ e $\bar n$ l'indice
iniziale della somma. Si hanno le seguenti proprietà:
\begin{enumerate}
    \item $\displaystyle \sum_{i=\bar n}^{n} (\lambda a_i)
          = \lambda \left( \sum_{i=\bar n}^{n} a_i \right)$, \quad
          $\lambda \in \mathbb{R}$;
    \item $\displaystyle \sum_{i=\bar n}^{n} (a_i + b_i)
          = \left( \sum_{i=\bar n}^{n} a_i \right)
          + \left( \sum_{i=\bar n}^{n} b_i \right)$;
    \item $\displaystyle \sum_{i=\bar n}^{n} a_i
          = \sum_{i=\bar n}^{k} a_i + \sum_{i=k+1}^{n} a_i$, \quad
          con $\bar n \leq k < n$;
    \item (cambio d'indice)
          $\displaystyle \sum_{i=1}^{n} a_i = \sum_{j=1}^{n} a_j
          = \sum_{k=1}^{n} a_k = \sum_{k=0}^{n-1} a_{k+1}$.
\end{enumerate}
\end{proposizione}

\subsection{Ricorsione}

\begin{definizione}
Una successione $\{a_n\}_{n \in \mathbb{N}}$ di numeri reali è
\emph{ricorsiva} se è definita specificando:
\begin{enumerate}
    \item il valore dei primi $m$ termini $a_1, a_2, \dots, a_m$;
    \item una funzione $f \colon \mathbb{R}^m \to \mathbb{R}$ tale che
          \[
            a_i = f(a_{i-1}, a_{i-2}, \dots, a_{i-m}) \qquad \forall\, i > m .
          \]
\end{enumerate}
\end{definizione}

\begin{esempio}
\begin{enumerate}
    \item Sia $a_0 = 0$ e $a_n = 1 + a_{n-1}$ per $n \geq 1$. Allora
          $a_1 = 1 + a_0 = 1$, $a_2 = 1 + a_1 = 2$, $a_3 = 1 + a_2 = 3$, \dots,
          $a_n = n$ (formula chiusa).
    \item \emph{Successione di Fibonacci}: $a_1 = 1$, $a_2 = 1$,
          $a_n = a_{n-1} + a_{n-2}$. Si ottiene $1, 1, 2, 3, 5, 8, \dots$
          (ad esempio $a_3 = a_2 + a_1 = 1 + 1 = 2$, $a_4 = a_3 + a_2 = 3$,
          $a_5 = a_4 + a_3 = 5$, $a_6 = a_5 + a_4 = 8$).
    \item Data $\{a_n\}$, la sua serie associata $\{s_n\}$ si definisce per
          ricorsione:
          \[
            \begin{cases}
              s_0 = a_0\\
              s_n = s_{n-1} + a_n
            \end{cases}
          \]
    \item $a_1 = 5$ e $a_n = 8$ per ogni $n > 1$: tutti i termini dopo il primo
          valgono $8$.
    \item \emph{Fattoriale}: $a_1 = 1$, $a_n = n \cdot a_{n-1}$. Allora
          $a_2 = 2 \cdot a_1 = 2$, $a_3 = 3 \cdot a_2 = 3 \cdot 2 \cdot 1$,
          $a_4 = 4 \cdot 3 \cdot 2 \cdot 1$, \dots, e in generale
          $a_n = n \cdot (n-1) \cdot (n-2) \cdots 3 \cdot 2 \cdot 1$.
          Si scrive $a_n = n!$ (formula esplicita), con $1! = 1$ e per
          convenzione $0! = 1$.
\end{enumerate}
\end{esempio}

\begin{definizione}
Data una successione $\{a_n\}$, una \emph{formula esplicita} o
\emph{formula chiusa} per $\{a_n\}$ è una formula $p(n)$ che esprime
direttamente il valore della successione $a_n$ a partire da $n$.

Se $\{a_n\}$ è definita per ricorsione, una formula chiusa per $\{a_n\}$ è
un'espressione $p(n)$ che verifica le condizioni iniziali e verifica la
relazione ricorsiva.
\end{definizione}

\begin{esempio}
Sia
\[
  \begin{cases}
    a_0 = 2\\
    a_n = a_{n-1} + 5
  \end{cases}
\]
Si ottiene $a_1 = 2 + 5 = 7$, $a_2 = 12$, $a_3 = 17$, $a_4 = 22$. La formula
chiusa $a_n = 5n + 2$ è quella giusta? Verifichiamo.
\begin{enumerate}
    \item \emph{Caso base} ($n = 0$): la formula dà $0 \cdot 5 + 2 = 2$, che
          coincide con il dato iniziale $a_0 = 2$.
    \item \emph{Relazione ricorsiva} $a_n = a_{n-1} + 5$: sostituendo la
          formula con $n-1$ al posto di $n$,
          \[
            a_{n-1} + 5 = \bigl( (n-1) \cdot 5 + 2 \bigr) + 5
            = 5n - 5 + 2 + 5 = 5n + 2 = a_n .
          \]
\end{enumerate}
Sono verificate entrambe: la formula chiusa è verificata.
\end{esempio}

\begin{esercizio}
Verificare che per
$\begin{cases} a_1 = 1\\ a_n = n + a_{n-1} \end{cases}$
la formula chiusa è $a_n = \dfrac{n(n+1)}{2}$.
\end{esercizio}

\emph{Svolgimento (aggiunto).} Per $n = 1$ la formula dà $\frac{1 \cdot 2}{2} = 1 = a_1$.
Inoltre
$a_{n-1} + n = \frac{(n-1)n}{2} + n = \frac{n^2 - n + 2n}{2} = \frac{n(n+1)}{2}$.

\subsubsection{Progressioni aritmetiche}

\begin{definizione}
Una successione $a_n \colon \mathbb{N} \to \mathbb{R}$ è una
\emph{progressione aritmetica} se
\[
  a_n - a_{n-1} = d, \qquad d \in \mathbb{R}.
\]
Definita in maniera ricorsiva, $a_0 = a$ e $a_n = d + a_{n-1}$, ha formula
chiusa
\[
  a_n = a + n \cdot d .
\]
\end{definizione}

\begin{esercizio}
Data $\begin{cases} a_0 = 8\\ a_n = a_{n-1} - 3 \end{cases}$, trovare la formula
chiusa e verificarla.
\end{esercizio}

\emph{Svolgimento (aggiunto).} È una progressione aritmetica con $a = 8$ e
$d = -3$, quindi $a_n = 8 - 3n$: per $n = 0$ vale $8$ e
$a_{n-1} - 3 = 8 - 3(n-1) - 3 = 8 - 3n = a_n$.

\subsection{Polinomi}

\begin{definizione}
Un \emph{polinomio nell'incognita $t$ a coefficienti in $\mathbb{R}$} è una
funzione $p \colon \mathbb{R} \to \mathbb{R}$ del tipo
\[
  p(t) = a_0 + a_1 t + a_2 t^2 + a_3 t^3 + \dots + a_n t^n
       = \sum_{i=0}^{n} a_i t^i ,
\]
dove $a_0, a_1, \dots, a_n \in \mathbb{R}$ sono i \emph{coefficienti} e $a_0$ è
il \emph{termine noto}.

Il più grande indice $k \in \mathbb{N}$ tale che $a_i = 0$ per $i > k$ e
$a_k \neq 0$ è il \emph{grado} del polinomio, $\operatorname{gr}(p) = k$. Se
$k$ non esiste, $\operatorname{gr}(p) = -\infty$.
\end{definizione}

\begin{esempio}
\begin{itemize}
    \item $o(t) = 0$ è il \emph{polinomio nullo}, con
          $\operatorname{gr}(o) = -\infty$.
    \item $q(t) = 3$ è un \emph{polinomio costante}
          ($a_0 = 3$, $a_1 = a_2 = \dots = 0$) e $\operatorname{gr}(q) = 0$.
    \item $p(t) = 5 - t + t^2$ ha $\operatorname{gr}(p) = 2$.
\end{itemize}
\end{esempio}

Si indica con $\mathbb{R}[t]$ l'insieme dei polinomi nell'incognita $t$ a
coefficienti reali, e con
\[
  \mathbb{R}_d[t] = \{ p \in \mathbb{R}[t] \mid \operatorname{gr}(p) \leq d \}
\]
quelli di grado al più $d$. Ad esempio $5 + t^2$ e $1 - 3t - t^2$ stanno in
$\mathbb{R}_2[t]$, il polinomio nullo $o$ sta in $\mathbb{R}_0[t]$, mentre
$5 + t^2 \notin \mathbb{R}_1[t]$.

\begin{osservazione}
Due polinomi $p = \sum_{i=0}^{n} a_i t^i$ e $q = \sum_{i=0}^{n} b_i t^i$ sono
uguali, $p = q$, \textbf{se e solo se} $a_i = b_i$ per ogni $i$. Ad esempio
$(x+2)^2 = x^2 + 4x + 4$.
\end{osservazione}

\subsubsection{Operazioni}

Dati $p = \sum_{i=0}^{n} a_i t^i$, $q = \sum_{i=0}^{n} b_i t^i$ e
$\lambda \in \mathbb{R}$, si definiscono:
\begin{itemize}
    \item il \emph{prodotto per uno scalare}:
          $(\lambda p) = \sum_{i=0}^{n} (\lambda a_i)\, t^i$;
    \item la \emph{somma}: $(p + q) = \sum_{i=0}^{n} (a_i + b_i)\, t^i$;
    \item il polinomio $t p = \sum_{i=0}^{n} a_i t^{i+1}
          = t a_0 + a_1 t^2 + \dots + a_n t^{n+1}
          = \sum_{i=1}^{n+1} a_{i-1} t^i$;
    \item il \emph{prodotto}:
          $p \cdot q = (a_0 + \dots + a_n t^n) \cdot q
          = a_0 q + a_1 t q + \dots + a_n t^n q$.
\end{itemize}

\begin{esempio}
Siano $p = 6 + 3t - t^2$ e $q = -6 + 5t$. Allora
\begin{align*}
p + q &= (6 - 6) + (3 + 5)\,t + (-1 + 0)\,t^2 = 8t - t^2,\\
p \cdot q &= (6 + 3t - t^2)(-6 + 5t)
  = 6(-6 + 5t) + 3t(-6 + 5t) - t^2(-6 + 5t)\\
  &= -36 + 30t - 18t + 15t^2 + 6t^2 - 5t^3
  = -36 + 12t + 21t^2 - 5t^3 .
\end{align*}
\end{esempio}

\begin{osservazione}
Per i gradi valgono
\[
  \operatorname{gr}(p + q) \leq \max\{\operatorname{gr}(p), \operatorname{gr}(q)\},
  \qquad
  \operatorname{gr}(p \cdot q) = \operatorname{gr}(p) + \operatorname{gr}(q).
\]
\end{osservazione}

\subsection{Divisione tra polinomi}

\begin{teorema}[Divisione con resto]
Siano $f, g \in \mathbb{R}[t]$ con $g \neq 0$. Allora esistono e sono unici due
polinomi $q, r \in \mathbb{R}[t]$ tali che
\[
  f = q \cdot g + r , \qquad \operatorname{gr}(r) < \operatorname{gr}(g).
\]
$q$ è il \emph{quoziente} e $r$ il \emph{resto}.
\end{teorema}

\begin{esempio}
Siano $f(t) = t^4 - t^3 + t - 1$ e $g(t) = t^2 - t$. Si ha
$t^2 \cdot g = t^4 - t^3$, quindi
\[
  f - t^2 g = t - 1 ,
\]
che ha grado $1 < \operatorname{gr}(g) = 2$: la divisione si ferma. Dunque
$q = t^2$, $r = t - 1$ e
\[
  f(t) = t^2 \cdot g + (t - 1).
\]
\end{esempio}

\begin{esempio}
Siano $f(t) = t^4 - t^3 + t - 1$ e $g(t) = t + 1$, con
$\operatorname{gr}(g) = 1$. Ad ogni passo si sottrae al resto parziale il
multiplo di $g$ che ne elimina il termine di grado massimo:
\[
\begin{array}{rcl}
  t^4 - t^3 + 0t^2 + t - 1 \;-\; t^3 g &=& -2t^3 + 0t^2 + t - 1 \\
  -2t^3 + 0t^2 + t - 1 \;-\; (-2t^2)\, g &=& 2t^2 + t - 1 \\
  2t^2 + t - 1 \;-\; (2t)\, g &=& -t - 1 \\
  -t - 1 \;-\; (-1)\, g &=& 0
\end{array}
\]
Il quoziente è $q = t^3 - 2t^2 + 2t - 1$ e il resto è $r = 0$
($\operatorname{gr}(0) = -\infty < \operatorname{gr}(g)$). Quindi
\[
  f(t) = (t^3 - 2t^2 + 2t - 1)\, g + 0 .
\]
\end{esempio}

\begin{definizione}
Siano $f, g \in \mathbb{R}[t]$ con $g \neq 0$. Si dice che $g$ \emph{divide}
$f$, e si scrive $g \mid f$, se esiste $q \in \mathbb{R}[t]$ tale che
$f = q \cdot g$ (cioè se il resto della divisione è nullo).
\end{definizione}

\section[Radici di polinomi e ricorsione lineare (Lezione 6 -- 6 ottobre 2026)]{Radici di polinomi e ricorsione lineare}
\noindent\emph{Lezione 6 -- 6 ottobre 2026}
\subsection{Radici di polinomi}

\begin{definizione}
Sia $f \in \mathbb{R}[t]$. Un numero $\alpha \in \mathbb{R}$ si dice \emph{radice}
di $f$ se $f(\alpha) = 0$.
\end{definizione}

Il legame con la divisibilità, vista nella lezione precedente, è il seguente.

\begin{teorema}[del fattore]
Sia $f \in \mathbb{R}[t]$ e sia $\alpha \in \mathbb{R}$. Allora
\[
  \alpha \text{ è una radice di } f
  \iff
  (t - \alpha) \mid f .
\]
\end{teorema}

\begin{proof}
Dividiamo $f$ per $t - \alpha$ (che ha grado $1$): per il teorema di divisione
con resto esistono $q, r \in \mathbb{R}[t]$ tali che
\[
  f(t) = q(t)\,(t - \alpha) + r(t), \qquad \operatorname{gr}(r) < 1 ,
\]
quindi $r$ è un polinomio costante. Valutando in $t = \alpha$ si ottiene
\[
  f(\alpha) = q(\alpha)\cdot 0 + r = r .
\]
Dunque $f(\alpha) = 0$ se e solo se $r = 0$, cioè se e solo se $t - \alpha$
divide $f$.
\end{proof}

\begin{osservazione}
Il resto della divisione di $f$ per $t - \alpha$ è quindi esattamente il valore
$f(\alpha)$. Se $\alpha$ è radice di $f$ si può scrivere $f = (t - \alpha)\, q$,
con $\operatorname{gr}(q) = \operatorname{gr}(f) - 1$: trovare una radice
permette di abbassare il grado del problema.
\end{osservazione}

\begin{definizione}
Sia $f \neq 0$. Una radice $\alpha$ ha \emph{molteplicità} (algebrica)
$m \geq 1$ se $(t - \alpha)^m \mid f$ ma $(t - \alpha)^{m+1} \nmid f$, cioè se
$f = (t - \alpha)^m\, g$ con $g(\alpha) \neq 0$. Le radici di molteplicità $1$
si dicono \emph{semplici}.
\end{definizione}

\subsubsection{Ricerca delle radici e scomposizione}

\begin{osservazione}
Se $f = a_n t^n + \dots + a_1 t + a_0$ ha coefficienti \emph{interi} e
$\frac{p}{q}$ (frazione ridotta) è una radice razionale di $f$, allora
$p \mid a_0$ e $q \mid a_n$. In particolare, se $f$ è monico, le radici
razionali sono intere e dividono $a_0$: basta provare i divisori di $a_0$
(con entrambi i segni).
\end{osservazione}

\begin{esempio}[radici semplici]
Scomponiamo $f(t) = t^3 - 2t^2 - 5t + 6$. Le possibili radici intere sono i
divisori di $6$: $\pm 1, \pm 2, \pm 3, \pm 6$. Si trova
\[
  f(1) = 1 - 2 - 5 + 6 = 0 ,
\]
quindi $1$ è radice e $(t - 1) \mid f$. Per trovare il quoziente si può
usare la regola di Ruffini: si riportano i coefficienti di $f$ e si procede
da sinistra a destra, moltiplicando per $\alpha = 1$ e sommando.
\[
\begin{array}{r|rrrr}
   & 1 & -2 & -5 & 6 \\
 1 &   &  1 & -1 & -6 \\ \hline
   & 1 & -1 & -6 & \;0
\end{array}
\]
L'ultimo numero è il resto ($=f(1)=0$), gli altri sono i coefficienti del
quoziente: $f = (t - 1)(t^2 - t - 6)$. Il fattore di secondo grado si
scompone cercando due numeri con somma $-1$ e prodotto $-6$:
$t^2 - t - 6 = (t - 3)(t + 2)$. Quindi
\[
  f(t) = (t - 1)(t - 3)(t + 2),
\]
e le radici sono $1$, $3$, $-2$, tutte semplici.
\end{esempio}

\begin{esempio}[radice multipla]
Sia $f(t) = t^3 - 3t^2 + 4$. I divisori di $4$ sono $\pm 1, \pm 2, \pm 4$ e
\[
  f(2) = 8 - 12 + 4 = 0 , \qquad f(-1) = -1 - 3 + 4 = 0 .
\]
Dividendo per $t - 2$ (Ruffini con $\alpha = 2$ sui coefficienti
$1, -3, 0, 4$) si ottiene $f = (t - 2)(t^2 - t - 2)$ e poi
$t^2 - t - 2 = (t - 2)(t + 1)$. Dunque
\[
  f(t) = (t - 2)^2 (t + 1) .
\]
La radice $2$ ha molteplicità $2$, la radice $-1$ ha molteplicità $1$.
\end{esempio}

\subsection{Ricorsione lineare}

\begin{definizione}
Una successione $(a_n)_{n \geq 0}$ è \emph{ricorsiva lineare omogenea di grado
$k$} (a coefficienti costanti) se esistono $c_1, \dots, c_k \in \mathbb{R}$,
con $c_k \neq 0$, tali che per ogni $n \geq k$
\[
  a_n = c_1\, a_{n-1} + c_2\, a_{n-2} + \dots + c_k\, a_{n-k} ,
\]
e se sono assegnati i primi $k$ valori $a_0, \dots, a_{k-1}$ (\emph{condizioni
iniziali}). Il polinomio
\[
  p(t) = t^k - c_1\, t^{k-1} - c_2\, t^{k-2} - \dots - c_{k-1}\, t - c_k
\]
si dice \emph{polinomio caratteristico} della ricorrenza.
\end{definizione}

``Lineare'' significa che $a_n$ è combinazione lineare dei termini precedenti,
``omogenea'' che non compaiono termini aggiuntivi che non dipendono dalla
successione.

\begin{osservazione}[perché il polinomio caratteristico]
Cerchiamo soluzioni della forma $a_n = \alpha^n$ con $\alpha \neq 0$.
Sostituendo nella ricorrenza e dividendo per $\alpha^{n-k}$ si ottiene
\[
  \alpha^k = c_1\, \alpha^{k-1} + \dots + c_k
  \iff p(\alpha) = 0 .
\]
Quindi $\alpha^n$ risolve la ricorrenza se e solo se $\alpha$ è una radice del
polinomio caratteristico. I teoremi seguenti dicono come combinare queste
soluzioni per soddisfare anche le condizioni iniziali.
\end{osservazione}

\begin{esempio}
\begin{itemize}
  \item $a_n = 3\, a_{n-1}$ ha grado $1$ e $p(t) = t - 3$; con $a_0$ dato si ha
        $a_n = a_0\, 3^n$.
  \item Fibonacci: $F_n = F_{n-1} + F_{n-2}$, con $F_0 = 0$, $F_1 = 1$. Grado
        $2$, $p(t) = t^2 - t - 1$.
  \item $a_n = 5\, a_{n-1} - 6\, a_{n-2}$ ha grado $2$ e
        $p(t) = t^2 - 5t + 6$.
  \item $a_n = 6 a_{n-1} - 11 a_{n-2} + 6 a_{n-3}$ ha grado $3$ e
        $p(t) = t^3 - 6t^2 + 11t - 6$.
\end{itemize}
\end{esempio}

\subsubsection{Teorema A: radici distinte}

\begin{teorema}[A]
Sia $(a_n)$ una successione ricorsiva lineare omogenea di grado $k$ il cui
polinomio caratteristico $p$ ha $k$ radici \emph{distinte}
$\alpha_1, \dots, \alpha_k$. Allora esistono costanti
$c_1, \dots, c_k$ tali che per ogni $n \geq 0$
\[
  a_n = c_1\, \alpha_1^{\,n} + c_2\, \alpha_2^{\,n} + \dots + c_k\, \alpha_k^{\,n} .
\]
Le costanti $c_i$ sono univocamente determinate dalle condizioni iniziali.
\end{teorema}

(Senza dimostrazione.)

\paragraph{Procedura per trovare la forma chiusa.}
\begin{enumerate}
  \item Scrivere il polinomio caratteristico $p(t)$.
  \item Trovare le sue $k$ radici distinte $\alpha_1, \dots, \alpha_k$
        (con il teorema del fattore).
  \item Scrivere $a_n = c_1 \alpha_1^{\,n} + \dots + c_k \alpha_k^{\,n}$.
  \item Imporre le condizioni iniziali $a_0, \dots, a_{k-1}$: si ottiene un
        sistema lineare di $k$ equazioni nelle incognite $c_1, \dots, c_k$.
  \item Risolvere il sistema e sostituire i valori trovati.
\end{enumerate}

\begin{esempio}
Sia $a_n = 5\, a_{n-1} - 6\, a_{n-2}$ con $a_0 = 0$, $a_1 = 1$. Il polinomio
caratteristico è $p(t) = t^2 - 5t + 6 = (t - 2)(t - 3)$, con radici distinte
$2$ e $3$. Quindi
\[
  a_n = c_1\, 2^n + c_2\, 3^n .
\]
Le condizioni iniziali danno il sistema
\[
  \begin{cases}
    a_0 = c_1 + c_2 = 0 \\
    a_1 = 2c_1 + 3c_2 = 1
  \end{cases}
  \quad\Longrightarrow\quad c_2 = 1,\ c_1 = -1 ,
\]
perciò
\[
  a_n = 3^n - 2^n .
\]
Verifica: $a_2 = 5 \cdot 1 - 6 \cdot 0 = 5 = 3^2 - 2^2$ e
$a_3 = 5 \cdot 5 - 6 \cdot 1 = 19 = 3^3 - 2^3$.
\end{esempio}

\begin{esempio}[Fibonacci]
Per $F_n = F_{n-1} + F_{n-2}$, $F_0 = 0$, $F_1 = 1$, il polinomio
caratteristico è $t^2 - t - 1$, con radici distinte
\[
  \varphi = \frac{1 + \sqrt{5}}{2} , \qquad \psi = \frac{1 - \sqrt{5}}{2} .
\]
Allora $F_n = c_1\, \varphi^n + c_2\, \psi^n$ e
\[
  \begin{cases}
    c_1 + c_2 = 0 \\
    c_1\, \varphi + c_2\, \psi = 1
  \end{cases}
  \quad\Longrightarrow\quad
  c_1 (\varphi - \psi) = 1 ,\quad \varphi - \psi = \sqrt{5}
  \quad\Longrightarrow\quad c_1 = \frac{1}{\sqrt{5}},\ c_2 = -\frac{1}{\sqrt{5}} .
\]
Si ottiene la formula di Binet:
\[
  F_n = \frac{1}{\sqrt{5}}\left[\left(\frac{1 + \sqrt{5}}{2}\right)^{n}
        - \left(\frac{1 - \sqrt{5}}{2}\right)^{n}\right] .
\]
\end{esempio}

\begin{esempio}[grado 3]
Sia $a_n = 6 a_{n-1} - 11 a_{n-2} + 6 a_{n-3}$ con $a_0 = 3$, $a_1 = 6$,
$a_2 = 14$. Il polinomio caratteristico è
$p(t) = t^3 - 6t^2 + 11t - 6$. Provando i divisori di $6$ si trova
$p(1) = 0$, e con Ruffini $p(t) = (t - 1)(t^2 - 5t + 6) = (t-1)(t-2)(t-3)$.
Le radici $1, 2, 3$ sono distinte, quindi
\[
  a_n = c_1 + c_2\, 2^n + c_3\, 3^n ,
  \qquad
  \begin{cases}
    c_1 + c_2 + c_3 = 3 \\
    c_1 + 2c_2 + 3c_3 = 6 \\
    c_1 + 4c_2 + 9c_3 = 14
  \end{cases}
\]
Sottraendo la prima equazione dalla seconda e la seconda dalla terza si ha
$c_2 + 2c_3 = 3$ e $2c_2 + 6c_3 = 8$, da cui $c_3 = 1$, $c_2 = 1$, $c_1 = 1$.
Dunque
\[
  a_n = 1 + 2^n + 3^n .
\]
\end{esempio}

\subsubsection{Teorema B: radici con molteplicità}

\begin{teorema}[B]
Sia $(a_n)$ una successione ricorsiva lineare omogenea di grado $k$ il cui
polinomio caratteristico si scompone come
\[
  p(t) = (t - \alpha_1)^{m_1} (t - \alpha_2)^{m_2} \cdots (t - \alpha_r)^{m_r},
  \qquad m_1 + m_2 + \dots + m_r = k ,
\]
con $\alpha_1, \dots, \alpha_r$ distinte e $m_i$ la molteplicità di
$\alpha_i$. Allora
\[
  a_n = \sum_{i=1}^{r}
        \bigl(c_{i,0} + c_{i,1}\, n + \dots + c_{i,m_i - 1}\, n^{m_i - 1}\bigr)\,
        \alpha_i^{\,n} ,
\]
con $k$ costanti $c_{i,j}$ univocamente determinate dalle condizioni iniziali.
\end{teorema}

(Senza dimostrazione.)

\begin{osservazione}
Ogni radice $\alpha$ di molteplicità $m$ contribuisce con $m$ termini,
\[
  \alpha^n,\ n\,\alpha^n,\ n^2\alpha^n,\ \dots,\ n^{m-1}\alpha^n ,
\]
in totale $k$ costanti, quante le condizioni iniziali. Se tutte le molteplicità
sono $1$ si ritrova il Teorema A.
\end{osservazione}

\begin{esempio}[radice doppia]
Sia $a_n = 4\, a_{n-1} - 4\, a_{n-2}$ con $a_0 = 1$, $a_1 = 0$. Il polinomio
caratteristico è $p(t) = t^2 - 4t + 4 = (t - 2)^2$: un'unica radice $2$ di
molteplicità $2$. Quindi
\[
  a_n = (c_0 + c_1\, n)\, 2^n .
\]
Dalle condizioni iniziali:
\[
  a_0 = c_0 = 1 , \qquad a_1 = 2\,(c_0 + c_1) = 0 \Rightarrow c_1 = -1 ,
\]
perciò
\[
  a_n = (1 - n)\, 2^n .
\]
Verifica: $a_2 = 4 \cdot 0 - 4 \cdot 1 = -4 = (1-2)\cdot 2^2$ e
$a_3 = 4\cdot(-4) - 4 \cdot 0 = -16 = (1-3)\cdot 2^3$.
\end{esempio}

\begin{esempio}[radici di molteplicità diversa]
Sia $a_n = 5\, a_{n-1} - 8\, a_{n-2} + 4\, a_{n-3}$ con $a_0 = 2$, $a_1 = 5$,
$a_2 = 13$. Il polinomio caratteristico è $p(t) = t^3 - 5t^2 + 8t - 4$.
Poiché $p(1) = 0$, con Ruffini si ottiene
\[
  p(t) = (t - 1)(t^2 - 4t + 4) = (t - 1)(t - 2)^2 ,
\]
cioè la radice $1$ ha molteplicità $1$ e la radice $2$ ha molteplicità $2$.
Per il Teorema B
\[
  a_n = c + (d + e\, n)\, 2^n ,
  \qquad
  \begin{cases}
    c + d = 2 \\
    c + 2d + 2e = 5 \\
    c + 4d + 8e = 13
  \end{cases}
\]
Sottraendo la prima equazione dalla seconda e la seconda dalla terza:
$d + 2e = 3$ e $2d + 6e = 8$, cioè $d + 3e = 4$. Quindi $e = 1$, $d = 1$,
$c = 1$ e
\[
  a_n = 1 + (1 + n)\, 2^n .
\]
Verifica: $a_3 = 5 \cdot 13 - 8 \cdot 5 + 4 \cdot 2 = 33 = 1 + 4 \cdot 8$.
\end{esempio}

\section[Ricorsione lineare: un esempio con radice doppia (Lezione 7 -- 12 ottobre 2026)]{Ricorsione lineare: un esempio con radice doppia}
\noindent\emph{Lezione 7 -- 12 ottobre 2026}
\subsection{Esempio con radici di molteplicità diversa}

Applichiamo il Teorema~B (radici con molteplicità) alla successione ricorsiva lineare omogenea a coefficienti costanti di grado $k = 3$
\[
\begin{cases}
a_0 = 2 \\
a_1 = -1 \\
a_2 = -8 \\
a_n = 5a_{n-1} - 3a_{n-2} - 9a_{n-3} & (n \geq 3).
\end{cases}
\]

\paragraph{Polinomio caratteristico.} L'equazione caratteristica è
\[
x^3 - 5x^2 + 3x + 9 = 0 .
\]

\paragraph{Radici.} Il polinomio è monico a coefficienti interi, quindi le radici razionali si cercano tra i divisori di $9$: $\pm 1, \pm 3, \pm 9$. Si trova che $3$ e $-1$ sono radici e
\[
(x - 3)^2 (x + 1) = 0 .
\]
\begin{itemize}
    \item $\lambda_1 = 3$ è radice con molteplicità algebrica doppia: $m_a(3) = 2$ ($= m_1$);
    \item $\lambda_2 = -1$ è radice semplice: $m_a(-1) = 1$ ($= m_2$);
    \item la somma delle molteplicità è $m_a(3) + m_a(-1) = 3 = k$, come deve essere.
\end{itemize}

\paragraph{Forma generale.} Per il Teorema~B
\[
a_n = (c_{1,0} + c_{1,1}\, n)\, 3^n + c_{2,0}\, (-1)^n .
\]

\paragraph{Condizioni iniziali.} Per trovare le costanti si impongono le condizioni iniziali (formula $=$ dato):
\[
\begin{cases}
n = 0: & c_{1,0} + c_{2,0} = 2 \\
n = 1: & (c_{1,0} + c_{1,1})\cdot 3 + c_{2,0}\,(-1) = -1 \\
n = 2: & (c_{1,0} + 2c_{1,1})\cdot 3^2 + c_{2,0}\,(-1)^2 = -8
\end{cases}
\]
Dalla prima equazione $c_{1,0} = 2 - c_{2,0}$. Sostituendo nelle altre due:
\[
\begin{cases}
6 - 3c_{2,0} + 3c_{1,1} - c_{2,0} = -1 \\
18 - 9c_{2,0} + 18c_{1,1} + c_{2,0} = -8
\end{cases}
\iff
\begin{cases}
-4c_{2,0} + 3c_{1,1} = -7 \\
-8c_{2,0} + 18c_{1,1} = -26
\end{cases}
\]
Dalla prima, $c_{1,1} = -\tfrac{7}{3} + \tfrac{4}{3} c_{2,0}$; sostituendo nella seconda
\[
-8c_{2,0} - 42 + 24c_{2,0} = -26 \iff 16 c_{2,0} = 16 \iff c_{2,0} = 1 .
\]
Quindi
\[
c_{2,0} = 1, \qquad c_{1,1} = -\tfrac{7}{3} + \tfrac{4}{3} = -1, \qquad c_{1,0} = 2 - 1 = 1 .
\]

\paragraph{Formula chiusa.}
\[
a_n = (1 - n)\, 3^n + (-1)^n .
\]

Verifica: $a_0 = 1 + 1 = 2$, $a_1 = 0 - 1 = -1$, $a_2 = (-1)\cdot 9 + 1 = -8$ e, dalla ricorrenza, $a_3 = 5(-8) - 3(-1) - 9\cdot 2 = -55 = (1-3)\cdot 27 - 1$.

\section[Calcolo combinatorio: regole di base e notazioni (Lezione 7 -- 12 ottobre 2026)]{Calcolo combinatorio: regole di base e notazioni}
\noindent\emph{Lezione 7 -- 12 ottobre 2026}
\subsection{Regole di somma e di prodotto}

Siano $A, B$ insiemi finiti con $|A| = n$ e $|B| = k$.

\begin{proprieta}[Regola della somma]
Se $A \cap B = \emptyset$ allora $|A \cup B| = n + k$.
\end{proprieta}

\begin{proprieta}[Regola del prodotto]
$|A \times B| = n \cdot k$.
\end{proprieta}

Infatti gli elementi di $A \times B$ sono le coppie $(a_i, b_j)$ e si possono disporre in una tabella con $n$ righe (una per ogni $a_i$) e $k$ colonne (una per ogni $b_j$):
\begin{center}
\begin{tabular}{c|ccc}
$A \times B$ & $b_1$ & $\cdots$ & $b_k$ \\ \midrule
$a_1$ & $(a_1, b_1)$ & $\cdots$ & $(a_1, b_k)$ \\
$\vdots$ & $\vdots$ & & $\vdots$ \\
$a_n$ & $(a_n, b_1)$ & $\cdots$ & $(a_n, b_k)$
\end{tabular}
\end{center}

\begin{esempio}
Una collana ha una pietra, che può essere Blu, Verde o Rossa, e un simbolo, che può essere Cuore, Stella, Cerchio, Quadrato o Fiore. Si possono costruire $3 \cdot 5 = 15$ collane con una pietra e un simbolo.
\end{esempio}

\subsection{Notazioni}

\begin{definizione}
Il \emph{fattoriale} di $n \in \mathbb{N}$ è
\[
n! = \begin{cases} 1 & \text{se } n = 0, \\ n \cdot (n-1)! & \text{se } n \geq 1, \end{cases}
\]
cioè $n! = n \cdot (n-1) \cdots 2 \cdot 1$.
\end{definizione}

\begin{definizione}
Per $n \geq k$ il \emph{coefficiente binomiale} ``$n$ su $k$'' è
\[
\binom{n}{k} = \frac{n!}{k!\,(n-k)!} = C_{n,k}.
\]
\end{definizione}

\subsection{Come si scelgono $k$ oggetti tra $n$}

Dati $n$ oggetti, ne scegliamo $k$: in quanti modi possiamo fare la scelta? La risposta dipende da due domande: l'\emph{ordine} conta? sono ammesse \emph{ripetizioni}? Prendiamo come esempio $X = \{x, y, z\}$.

\begin{center}
\renewcommand{\arraystretch}{1.4}
\begin{tabular}{@{}l|p{5.2cm}|p{5.2cm}@{}}
 & \textbf{Insiemi} (l'ordine non conta) & \textbf{Liste} (l'ordine è importante) \\ \midrule
\textbf{Senza ripetizioni}
 & \textbf{Combinazioni semplici}: $\{x,y\} = \{y,x\}$; $\{x,x\}$ non accettabile
 & \textbf{Disposizioni semplici}: $(x,y) \neq (y,x)$; $(x,x)$ non accettabile \\ \midrule
\textbf{Con ripetizioni}
 & \textbf{Combinazioni con ripetizioni}: $\{x,y\} = \{y,x\}$; $\{x,x\}$ accettabile
 & \textbf{Disposizioni con ripetizioni}: $(x,y) \neq (y,x)$; $(x,x)$ accettabile
\end{tabular}
\end{center}

\section[Disposizioni, permutazioni e combinazioni semplici (Lezione 7 -- 12 ottobre 2026)]{Disposizioni, permutazioni e combinazioni semplici}
\noindent\emph{Lezione 7 -- 12 ottobre 2026}
\subsection{Disposizioni semplici}

\begin{definizione}
Una \emph{disposizione semplice} di $n$ oggetti in gruppi di $k$ ($k \leq n$) è una \emph{lista ordinata} formata da $k$ oggetti scelti tra gli $n$ in modo che non ci siano ripetizioni.
\end{definizione}

\begin{proposizione}
Il numero di disposizioni semplici di $n$ oggetti in gruppi di $k$ è
\[
D_{n,k} = \underbrace{n \cdot (n-1)\cdot(n-2) \cdots (n-k+1)}_{k \text{ fattori}} .
\]
\end{proposizione}

(Per la regola del prodotto: il primo oggetto della lista si sceglie in $n$ modi, il secondo in $n-1$, \dots, il $k$-esimo in $n-k+1$.)

\begin{esempio}
Gara di nuoto con $10$ partecipanti: quanti sono i podi possibili? Sono disposizioni semplici di $10$ elementi in gruppi da $3$, quindi i podi sono $10 \cdot 9 \cdot 8 = D_{10,3} = 720$.
\end{esempio}

\subsection{Permutazioni}

\begin{osservazione}
Una disposizione semplice di $n$ oggetti in gruppi di $n$ è detta \emph{permutazione}. Il numero di permutazioni di $n$ oggetti è
\[
n! = D_{n,n}.
\]
\end{osservazione}

\begin{esempio}
Nella gara di prima ci sono $10!$ possibili classifiche.
\end{esempio}

\subsection{Combinazioni semplici}

\begin{definizione}
Una \emph{combinazione semplice} di $n$ oggetti in gruppi di $k$ ($k \leq n$) è un \emph{insieme} di $k$ elementi scelti tra gli $n$ oggetti in modo tale che non ci siano ripetizioni.
\end{definizione}

\begin{proposizione}
Il numero di combinazioni semplici di $n$ oggetti in gruppi di $k$ è dato da
\[
C_{n,k} = \binom{n}{k} = \frac{n!}{k!\,(n-k)!} = \frac{D_{n,k}}{k!} .
\]
\end{proposizione}

\begin{esempio}
Nella gara, in quanti modi $3$ atleti possono ricevere una medaglia qualsiasi (non interessa se arrivano primi o secondi)? L'ordine non conta:
\[
C_{10,3} = \binom{10}{3} = \frac{10!}{3!\,7!} = \frac{10 \cdot 9 \cdot 8}{3 \cdot 2} = 120 .
\]
\end{esempio}

\begin{osservazione}
Dato un insieme $A$ con $|A| = n$, $C_{n,k}$ è il numero di \emph{sottoinsiemi} di $A$ con $k$ elementi.
\end{osservazione}

\section[Calcolo combinatorio con ripetizioni (Lezione 7 -- 12 ottobre 2026)]{Calcolo combinatorio con ripetizioni}
\noindent\emph{Lezione 7 -- 12 ottobre 2026}
\subsection{Disposizioni con ripetizioni}

\begin{definizione}
Una \emph{disposizione con ripetizioni} di $n$ oggetti in gruppi di $k$ è una \emph{lista ordinata} di $k$ elementi scelti tra gli $n$ oggetti.
\end{definizione}

\begin{proposizione}
Il numero di disposizioni con ripetizioni di $n$ oggetti in gruppi di $k$ è
\[
D'_{n,k} = n^k .
\]
\end{proposizione}

\begin{esempio}
Quanti numeri naturali esistono con $7$ cifre tutte dispari? Gli oggetti sono le cifre dispari $\{1, 3, 5, 7, 9\}$, cioè $n = 5$. Contiamo le disposizioni con ripetizioni di $5$ oggetti in gruppi di $7$: $D'_{5,7} = 5^7 = 78\,125$.
\end{esempio}

\subsection{Permutazioni con ripetizioni}

\begin{definizione}
Dati $k$ oggetti distinti $a_1, a_2, \dots, a_k$, una \emph{permutazione} di $a_1, \dots, a_k$ in cui l'oggetto $a_1$ si ripete $n_1$ volte, $a_2$ si ripete $n_2$ volte, \dots, $a_k$ si ripete $n_k$ volte è una lista di $n = n_1 + n_2 + \dots + n_k$ oggetti in cui $a_1$ compare $n_1$ volte, \dots, $a_k$ compare $n_k$ volte.
\end{definizione}

\begin{proposizione}
Il numero di permutazioni di $k$ oggetti $a_1, \dots, a_k$ in cui $a_1$ si ripete $n_1$ volte, \dots, $a_k$ si ripete $n_k$ volte è
\[
\frac{n!}{n_1!\, n_2! \cdots n_k!} .
\]
\end{proposizione}

\begin{esempio}[Anagrammi]
\leavevmode
\begin{itemize}
    \item \emph{GUFO}: sono permutazioni di $4$ simboli distinti, quindi $4! = 4 \cdot 3 \cdot 2 = 24$ anagrammi (ad esempio UFOG, OFUG).
    \item \emph{GATTO}: sono permutazioni dei $4$ simboli G, A, T, O in cui T si ripete $2$ volte (liste di $5$ lettere); gli anagrammi sono $\dfrac{5!}{2!} = 60$ (ad esempio OTTAG, GATOT).
    \item \emph{MAMMA}: $A$ si ripete $2$ volte e $M$ si ripete $3$ volte, quindi $\dfrac{5!}{2!\,3!} = \dfrac{5 \cdot 4}{2} = 10$.
\end{itemize}
\end{esempio}

\subsection{Combinazioni con ripetizioni}

\begin{definizione}
Una \emph{combinazione con ripetizioni} di $n$ oggetti in gruppi di $k$ è un \emph{insieme} di $k$ elementi scelti tra gli $n$ oggetti (un oggetto può essere scelto più volte).
\end{definizione}

\begin{esempio}[Frutta]
Vogliamo comprare $3$ frutti e ci sono Mele, Pere, Kiwi, Arance, Banane. Rappresentiamo una scelta con \emph{stelle} (i frutti scelti) e \emph{barre} (i separatori tra i $5$ tipi di frutto). Ad esempio, scrivendo i tipi nell'ordine Mele $|$ Pere $|$ Kiwi $|$ Arance $|$ Banane:
\begin{center}
\begin{tabular}{@{}ll@{}}
$\ast \mid \ast \mid\ \mid \ast \mid\ $ & una mela, una pera, un'arancia \\
$\ast\ast\ast \mid\ \mid\ \mid\ \mid\ $ & tre mele \\
$\ \mid\ \mid \ast\ast \mid\ \mid \ast$ & due kiwi e una banana
\end{tabular}
\end{center}
Dobbiamo allora contare le permutazioni di $2$ simboli, $\ast$ e $\mid$, con ripetizioni: $\ast$ si ripete $3$ volte e $\mid$ si ripete $4$ volte. Sono
\[
\frac{7!}{3!\,4!} = 35 .
\]
\end{esempio}

\begin{proposizione}
Il numero di combinazioni con ripetizioni di $n$ oggetti in gruppi di $k$ è dato da
\[
C'_{n,k} = \frac{(n+k-1)!}{k!\,(n-1)!} = C_{n+k-1,\,k} = \binom{n+k-1}{k} .
\]
\end{proposizione}

(Nell'esempio $n = 5$ tipi e $k = 3$ frutti: $k$ stelle e $n - 1 = 4$ barre, in tutto $n + k - 1 = 7$ simboli.)

\section[Inclusione-esclusione e conteggio (Lezione 8 -- 13 ottobre 2026)]{Inclusione-esclusione e conteggio}
\noindent\emph{Lezione 8 -- 13 ottobre 2026}
\subsection{Principio di inclusione-esclusione}

\begin{proposizione}
Siano $A, B$ due insiemi finiti. Allora
\[
|A \cup B| = |A| + |B| - |A \cap B|.
\]
\end{proposizione}

\emph{Dimostrazione.} Sommando $|A|$ e $|B|$ gli elementi di $A \cap B$ vengono contati due volte. Perciò $|A \cup B| + |A \cap B| = |A| + |B|$, da cui la formula. \hfill $\square$

\begin{esempio}
Quanti numeri naturali positivi $\le 43$ sono divisibili per $4$ oppure per $5$?
Sia $A = \{n \in \mathbb{N} : n \le 43,\ 4 \mid n\}$ e $B = \{n \in \mathbb{N} : n \le 43,\ 5 \mid n\}$. Allora $|A| = 10$ e $|B| = 8$. L'intersezione è
\[
A \cap B = \{n \le 43 : 5 \mid n \text{ e } 4 \mid n\} = \{n \le 43 : 20 \mid n\} = \{20, 40\},
\]
quindi $|A \cap B| = 2$ e $|A \cup B| = 10 + 8 - 2 = 16$.
\end{esempio}

\begin{proposizione}
Siano $A, B, C$ tre insiemi finiti. Allora
\[
|A \cup B \cup C| = |A| + |B| + |C| - |A \cap B| - |A \cap C| - |B \cap C| + |A \cap B \cap C|.
\]
\end{proposizione}

\begin{esempio}[Evento sportivo]
Un evento ha tre gare e $20$ atleti partecipanti. Sia $A$ l'insieme degli atleti della gara $1$, $B$ della gara $2$, $C$ della gara $3$. Si sa che $|A| = 6$, $|B| = 10$, $|C| = 13$, $|A \cap B| = 4$, $|A \cap C| = 3$, $|B \cap C| = 5$. Quanti atleti partecipano a tutte e tre le gare?
Il numero totale di atleti è $|A \cup B \cup C| = 20$. Ricavando $|A \cap B \cap C|$ dalla formula:
\[
|A \cap B \cap C| = |A \cup B \cup C| - |A| - |B| - |C| + |A \cap B| + |A \cap C| + |B \cap C| = 20 - 6 - 10 - 13 + 4 + 3 + 5 = 3.
\]
\end{esempio}

\subsection{Contare le funzioni tra insiemi finiti}

Siano $A, B$ insiemi finiti con $|A| = n$ e $|B| = k$.

\begin{itemize}
    \item \emph{Funzioni da $A$ a $B$}: per ogni elemento di $A$ si sceglie un'immagine tra i $k$ elementi di $B$. Sono le disposizioni con ripetizioni: $D'_{k,n} = k^n$.
    \item \emph{Funzioni biunivoche} (possibili se $n = k$): sono le permutazioni di $n$ oggetti, $n! = n \cdot (n-1) \cdots 2 \cdot 1$.
    \item \emph{Funzioni iniettive} (possibili se $n \le k$): sono le disposizioni semplici, $D_{k,n} = k(k-1)\cdots(k-n+1)$.
    \item \emph{Funzioni suriettive} (possibili se $n \ge k$):
    \begin{itemize}
        \item se $|B| = 1$ c'è la sola funzione costante, che è suriettiva: $1$ funzione;
        \item se $|B| = 2$ le funzioni sono $2^n$; non sono suriettive le $2$ funzioni costanti, quindi le suriettive sono $2^n - 2$;
        \item se $|B| = 3$ le funzioni sono $3^n$. Quelle che lasciano fuori almeno un elemento di $B$ sono $3 \cdot 2^n - 3$: il termine $3 \cdot 2^n$ conta, per ciascuno dei $3$ elementi, le $2^n$ funzioni che lo escludono, ma le $3$ funzioni costanti (che escludono due elementi) sono state contate due volte, quindi se ne sottraggono $3$. Le suriettive sono
        \[
        3^n - (3 \cdot 2^n - 3) = 3^n - 3 \cdot 2^n + 3.
        \]
    \end{itemize}
\end{itemize}

\subsection{Cardinalità dell'insieme delle parti}

\begin{proposizione}
Se $A$ è un insieme finito con $|A| = n$, l'insieme delle parti $\mathcal{P}(A)$ ha cardinalità $2^n$.
\end{proposizione}

\emph{Dimostrazione} per induzione su $n$. Sia $P(n)$: se $|A| = n$ allora $|\mathcal{P}(A)| = 2^n$.

\emph{Caso base} $n = 0$: $A = \emptyset$ e $\mathcal{P}(\emptyset) = \{\emptyset\}$, che ha $1 = 2^0$ elementi.

\emph{Passo induttivo.} Ipotesi: $P(k)$, cioè un insieme con $k$ elementi ha $2^k$ sottoinsiemi. Sia $|A| = k + 1$; scegliamo $\bar a \in A$. I sottoinsiemi $X \subseteq A$ sono di due tipi:
\begin{itemize}
    \item $\bar a \notin X$: allora $X \subseteq A \setminus \{\bar a\}$, che ha $k$ elementi, e per ipotesi induttiva questi $X$ sono $2^k$;
    \item $\bar a \in X$: allora $X = \{\bar a\} \cup Y$ con $Y \subseteq A \setminus \{\bar a\}$; di nuovo i possibili $Y$ sono $2^k$.
\end{itemize}
Quindi $|\mathcal{P}(A)| = 2^k + 2^k = 2 \cdot 2^k = 2^{k+1}$. \hfill $\square$

\subsection{Esercizi di conteggio}

\begin{esempio}[Carte]
Un mazzo ha $40$ carte ($4$ semi, $10$ valori per ogni seme).
\begin{enumerate}
    \item Quante mani di $5$ carte? Combinazioni semplici: $\binom{40}{5} = \dfrac{40!}{5!\,35!}$.
    \item Quante mani con $4$ assi? Le $4$ carte sono fissate, la quinta è una delle altre $36$: $\binom{36}{1} = 36$.
    \item Quante mani con $4$ carte dello stesso valore? Si sceglie il valore ($10$ modi) e poi l'ultima carta ($36$ modi): $10 \cdot 36$.
    \item Quante mani con \emph{esattamente} $2$ assi? Si scelgono i $2$ assi e poi $3$ carte che non sono assi: $\binom{4}{2}\binom{36}{3}$.
    \item Quante mani con \emph{almeno} $2$ assi? Si sommano i casi con esattamente $2$, $3$ o $4$ assi:
    \[
    \binom{4}{2}\binom{36}{3} + \binom{4}{3}\binom{36}{2} + \binom{4}{4}\binom{36}{1}.
    \]
    \item Quante mani con $2$ coppie di valori diversi (più una quinta carta)? Si scelgono i due valori ($\binom{10}{2}$), poi $2$ carte per il primo valore ($\binom{4}{2}$), $2$ per il secondo ($\binom{4}{2}$) e la quinta carta tra le $32$ rimaste con altri valori:
    \[
    \binom{10}{2}\binom{4}{2}\binom{4}{2}\binom{32}{1}.
    \]
\end{enumerate}
\end{esempio}

\begin{esempio}[Panchina]
In quanti modi $3$ maschi e $3$ femmine possono sedersi su una panchina?
\begin{enumerate}
    \item Senza vincoli: permutazioni di $6$ oggetti, $6!$.
    \item Se i maschi stanno tutti vicini e le femmine tutte vicine: $2$ ordini dei due gruppi, $3!$ per i maschi, $3!$ per le femmine: $2 \cdot 3! \cdot 3!$.
    \item Se solo i maschi stanno vicini: i $3$ maschi formano un blocco che, insieme alle $3$ femmine, dà $4$ oggetti da permutare ($4!$), poi si ordinano i maschi nel blocco ($3!$): $4! \cdot 3!$.
    \item Se maschi e femmine si alternano: si sceglie chi inizia ($2$), poi si ordinano i maschi ($3!$) e le femmine ($3!$): $2 \cdot 3! \cdot 3!$.
\end{enumerate}
\end{esempio}

\begin{esempio}[Urna]
Un'urna contiene $20$ palline bianche e $10$ nere. Si fanno $5$ estrazioni e si vuole ottenere $3$ bianche e $2$ nere.
\begin{enumerate}
    \item \emph{Senza reimmissione}: combinazioni semplici, $\binom{20}{3}\binom{10}{2}$.
    \item \emph{Con reimmissione} (si contano gli insiemi di palline scelte): combinazioni con ripetizioni,
    \[
    C'_{20,3}\cdot C'_{10,2} = \binom{22}{3}\binom{11}{2},
    \]
    perché $C'_{n,k} = C_{n+k-1,k}$.
    \item \emph{Senza reimmissione, palline indistinguibili e si tiene conto dell'ordine} (sequenze come $BBBNN$, $NBBNB, \dots$): sono permutazioni con ripetizioni, $\dfrac{5!}{3!\,2!}$.
\end{enumerate}
\end{esempio}

\subsection{Un esercizio di induzione}

\begin{esercizio}
Dimostrare per induzione che $25 \mid 2^{5n} - 7^n$ per ogni $n \ge 1$.
\end{esercizio}

\emph{Soluzione.} Sia $P(n)$: $25 \mid 2^{5n} - 7^n$.

\emph{Caso base} $n = 1$: $2^5 - 7 = 32 - 7 = 25$ e $25 \mid 25$.

\emph{Passo induttivo.} Ipotesi: $25 \mid 2^{5k} - 7^k$, cioè esiste $h \in \mathbb{Z}$ con $2^{5k} - 7^k = 25h$, ossia $2^{5k} = 25h + 7^k$. Dobbiamo dimostrare $25 \mid 2^{5(k+1)} - 7^{k+1}$. Si ha
\begin{align*}
2^{5k+5} - 7^{k+1} &= 2^5 \cdot 2^{5k} - 7 \cdot 7^k = 2^5 (25h + 7^k) - 7 \cdot 7^k \\
&= 25 \cdot 2^5 h + 7^k (2^5 - 7) = 25 \cdot 2^5 h + 7^k \cdot 25 \\
&= 25\,(2^5 h + 7^k),
\end{align*}
quindi $25 \mid 2^{5(k+1)} - 7^{k+1}$, che era la tesi. \hfill $\square$

\section[Operazioni interne, semigruppi e monoidi (Lezione 9 -- 19 ottobre 2026)]{Operazioni interne, semigruppi e monoidi}
\noindent\emph{Lezione 9 -- 19 ottobre 2026}
\subsection{Operazioni interne e gruppoidi}

\begin{definizione}
Una \emph{operazione interna} su un insieme $A$ è una funzione
\[
  * : A \times A \to A, \qquad (x, y) \mapsto x * y = {*}(x, y).
\]
L'insieme $A$ con l'operazione interna $*$ si chiama \emph{gruppoide} e si indica con $(A, *)$.
\end{definizione}

\begin{esempio}
\begin{itemize}
  \item $(\mathbb{Z}, +)$ con $+ : \mathbb{Z}\times\mathbb{Z} \to \mathbb{Z}$, $(h, k) \mapsto h + k$, è un gruppoide; è commutativo e associativo.
  \item $(\mathbb{N}, \cdot)$ con $\cdot : \mathbb{N}\times\mathbb{N} \to \mathbb{N}$, $(a, b) \mapsto a\cdot b$, è commutativo e associativo.
  \item $(\mathbb{N}, -)$ \emph{non} è un gruppoide: $-: \mathbb{N}\times\mathbb{N} \to \mathbb{Z}$, $(a, b) \mapsto a - b$ non ha valori in $\mathbb{N}$.
  \item $(\mathbb{Z}, -)$ è un gruppoide, ma non è commutativo e non è associativo: $(5-2)-6 = 3 - 6 = -3$, mentre $5 - (2 - 6) = 5 - (-4) = 9$.
  \item $(\mathbb{N}\setminus\{0\}, :)$ \emph{non} è un gruppoide: $: (a, b) \mapsto a : b = \frac{a}{b}$ va da $\mathbb{N}\setminus\{0\} \times \mathbb{N}\setminus\{0\}$ in $\mathbb{Q}$.
  \item Sia $A$ un insieme e $\mathcal{P}(A)$ l'insieme delle sue parti: $\cup : \mathcal{P}(A)\times\mathcal{P}(A) \to \mathcal{P}(A)$, $(X, Y) \mapsto X \cup Y$ è un'operazione interna; $(\mathcal{P}(A), \cup)$ è un gruppoide, commutativo e associativo (e lo stesso vale per $\cap$).
\end{itemize}
\end{esempio}

\subsubsection{Tavola di un'operazione}

Se $A = \{x_1, \dots, x_n\}$ è finito possiamo costruire la \emph{tavola} associata all'operazione $*$: nella riga di $x_i$ e colonna di $x_j$ si scrive $x_i * x_j$.
\[
\begin{array}{c|cccc}
  * & x_1 & x_2 & \cdots & x_n \\ \hline
  x_1 & x_1 * x_1 & x_1 * x_2 & \cdots & x_1 * x_n \\
  x_2 & x_2 * x_1 & x_2 * x_2 & \cdots & x_2 * x_n \\
  \vdots & \vdots & \vdots & \ddots & \vdots \\
  x_n & x_n * x_1 & x_n * x_2 & \cdots & x_n * x_n
\end{array}
\]
Se la tavola è simmetrica rispetto alla diagonale l'operazione è commutativa.

\begin{esempio}[$A = \{a\}$]
$\mathcal{P}(A) = \{\emptyset, \{a\}\} = \{\emptyset, A\}$. Le tavole di $\cap$ e $\cup$ sono
\[
\begin{array}{c|cc}
  \cap & \emptyset & A \\ \hline
  \emptyset & \emptyset & \emptyset \\
  A & \emptyset & A
\end{array}
\qquad\qquad
\begin{array}{c|cc}
  \cup & \emptyset & A \\ \hline
  \emptyset & \emptyset & A \\
  A & A & A
\end{array}
\]
Entrambe sono simmetriche rispetto alla diagonale, quindi le operazioni sono commutative. Per $\cap$ l'elemento $A$ è l'elemento neutro, per $\cup$ lo è $\emptyset$.
\end{esempio}

\begin{esempio}[$A = \{a, b\}$, differenza insiemistica]
Sia $|\mathcal{P}(A)| = 2^{|A|} = 2^2 = 4$. La tavola di $(\mathcal{P}(A), \setminus)$ è
\[
\begin{array}{c|cccc}
  \setminus & \emptyset & \{a\} & \{b\} & A \\ \hline
  \emptyset & \emptyset & \emptyset & \emptyset & \emptyset \\
  \{a\} & \{a\} & \emptyset & \{a\} & \emptyset \\
  \{b\} & \{b\} & \{b\} & \emptyset & \emptyset \\
  A & A & \{b\} & \{a\} & \emptyset
\end{array}
\]
Non è commutativa ($\{a\}\setminus\{b\} = \{a\}$ ma $\{b\}\setminus\{a\} = \{b\}$) e non esiste l'elemento neutro.
\end{esempio}

\begin{esempio}
Sia $A = \{e, a, b\}$ con l'operazione $*$ definita dalla tavola
\[
\begin{array}{c|ccc}
  * & e & a & b \\ \hline
  e & e & a & b \\
  a & a & e & e \\
  b & b & e & e
\end{array}
\]
$*$ è un'operazione interna, definita dalla tavola; è commutativa (la tavola è simmetrica) e $e$ è l'elemento neutro. \emph{Non} vale l'associatività perché $a$ ha due inversi: $a * a = e$ e $a * b = e$. Oppure: $(b * a) * a = e * a = a$ mentre $b * (a * a) = b * e = b$, quindi l'operazione non è associativa.
\end{esempio}

\subsection{Proprietà delle operazioni interne}

\begin{definizione}
Un gruppoide $(A, *)$ è detto
\begin{itemize}
  \item \emph{commutativo} se $\forall a, b \in A$: $a * b = b * a$;
  \item \emph{associativo} se $\forall a, b, c \in A$: $(a * b) * c = a * (b * c)$.
\end{itemize}
Un gruppoide associativo $(A, *)$ è detto \emph{semigruppo}.
\end{definizione}

\begin{definizione}[Elemento neutro]
Sia $(A, *)$ un gruppoide. Un elemento $e \in A$ è \emph{elemento neutro} per $*$ se
\[
  \forall a \in A \qquad a * e = e * a = a .
\]
\end{definizione}

\begin{definizione}[Elemento invertibile]
Se esiste l'elemento neutro $e \in A$, dato $a \in A$ l'\emph{inverso} di $a$ (se esiste) è l'elemento $a'$ tale che
\[
  a * a' = a' * a = e .
\]
Allora $a$ è detto \emph{elemento invertibile}.
\end{definizione}

\begin{esempio}
\begin{itemize}
  \item $(\mathbb{N}, +)$: $0$ è l'elemento neutro; $0$ ha inverso, nessun altro numero ha inverso.
  \item $(\mathbb{Z}, +)$: $0$ è l'elemento neutro; tutti i numeri sono invertibili (l'inverso di $k$ è $-k$).
  \item $(\mathbb{N}\setminus\{0\}, +)$: non c'è elemento neutro.
  \item $(\mathbb{N}, \cdot)$: $1$ è l'elemento neutro e $1$ è invertibile.
  \item $(\mathcal{P}(A), \cap)$: $A$ è l'elemento neutro ($X \subseteq A \Rightarrow A \cap X = X = X \cap A$). Non c'è inverso: solo $A$ è invertibile.
  \item $(\mathcal{P}(A), \cup)$: $\emptyset$ è l'elemento neutro. Non c'è inverso: solo $\emptyset$ è invertibile.
\end{itemize}
\end{esempio}

\begin{esempio}[$h * k = h + k + hk$ su $\mathbb{Z}$]
Sia $* : \mathbb{Z}\times\mathbb{Z} \to \mathbb{Z}$, $(h, k) \mapsto h * k = h + k + hk$. Cerchiamo, se esiste, un elemento neutro $x \in \mathbb{Z}$. Deve valere $x * k = k$, cioè $x + k + xk = k$, quindi $x(1 + k) = 0$. Poiché $k$ è un numero qualsiasi, $x$ deve essere zero. Verifichiamo che $x = 0$ sia elemento neutro: $0 * k = 0 + k + 0\cdot k = k$ e $h * 0 = h + 0 + h\cdot 0 = h$. Dunque $0$ è elemento neutro.
\end{esempio}

\begin{esempio}[$h * k = h^2 + k^2$ su $\mathbb{Z}$]
Esiste un elemento neutro? Se $x \in \mathbb{Z}$ fosse neutro, dovrebbe valere $x * k = k$, cioè $x^2 + k^2 = k$, cioè $x^2 = k - k^2$ per ogni $k$. Ma il secondo membro varia con $k$, quindi anche $x$ varierebbe: non ci sarà mai elemento neutro. (Ad esempio $0 * k = 0^2 + k^2 = k^2 \neq k$.)
\end{esempio}

\begin{definizione}[Monoide]
Un \emph{monoide} è un semigruppo $(A, *)$ con elemento neutro.
\end{definizione}

\begin{proposizione}
\begin{enumerate}
  \item Sia $(A, *)$ un gruppoide. Se esiste l'elemento neutro, è unico.
  \item Sia $(A, *)$ un monoide. Se $a$ è invertibile, allora l'inverso di $a$ è unico.
\end{enumerate}
\end{proposizione}

\begin{proof}
1) Siano $e_1, e_2 \in A$ elementi neutri in $(A, *)$. Poiché $e_1$ è neutro, $\forall x \in A$ vale $e_1 * x = x * e_1 = x$; poiché $e_2$ è neutro, $\forall x \in A$ vale $e_2 * x = x * e_2 = x$. Prendendo $x = e_2$ nella prima e $x = e_1$ nella seconda:
\[
  e_1 = e_1 * e_2 = e_2 ,
\]
cioè l'elemento neutro è unico.

2) Sia $a$ invertibile, con $a_1'$ e $a_2'$ inversi, cioè $a * a_1' = e = a_1' * a$ e $a * a_2' = e = a_2' * a$. Allora, usando l'associatività,
\[
  a_1' = a_1' * e = a_1' * (a * a_2') = (a_1' * a) * a_2' = e * a_2' = a_2' ,
\]
cioè $a_1' = a_2'$ e l'inverso è unico.
\end{proof}

\section[Gruppi (Lezione 9 -- 19 ottobre 2026)]{Gruppi}
\noindent\emph{Lezione 9 -- 19 ottobre 2026}
\subsection{Gruppi}

\begin{definizione}
$(G, *)$ è un \emph{gruppo} se
\begin{enumerate}
  \item $*$ è associativa: $\forall a, b, c \in G$ \quad $(a * b) * c = a * (b * c)$;
  \item esiste l'elemento neutro: $\exists e \in G$ tale che $\forall a \in G$ \quad $e * a = a * e = a$;
  \item ogni elemento è invertibile: $\forall a \in G$ \quad $\exists a' \in G$ tale che $a * a' = a' * a = e$.
\end{enumerate}
Un gruppo $(G, *)$ è \emph{abeliano} se $*$ è commutativa: $\forall a, b \in G$ \quad $a * b = b * a$.
\end{definizione}

\begin{esempio}
\begin{itemize}
  \item $(\mathbb{Q}\setminus\{0\}, \cdot)$ è un gruppo commutativo: l'operazione è associativa, l'elemento neutro è $1$ e l'inverso di $\frac{a}{b}$ è $\frac{b}{a}$ (se $q \in \mathbb{Q}$ l'inverso è $q^{-1}$).
  \item $(\mathbb{Z}, +)$ è un gruppo abeliano: associativa, $0$ elemento neutro, e se $k \in \mathbb{Z}$ allora $-k$ è l'inverso.
\end{itemize}
\end{esempio}

\subsubsection{Notazione}

\begin{center}
\begin{tabular}{@{}p{3.5cm} p{4.5cm} p{4.5cm}@{}}
\toprule
$(G, *)$ gruppo & Notazione additiva $(G, +)$ & Notazione moltiplicativa $(G, \cdot)$ \\
\midrule
$e$ elemento neutro & $0$ elemento neutro & $1$ elemento neutro \\[2pt]
$a'$ inverso di $a$ & $(-a)$ \emph{opposto} di $a$ & $a^{-1}$ \emph{inverso} di $a$ \\[2pt]
$\underbrace{a * a * \dots * a}_{n\text{ volte}}$ & $\underbrace{a + \dots + a}_{n\text{ volte}} = na$ & $\underbrace{a \cdot a \cdot \ldots \cdot a}_{n\text{ volte}} = a^n$ \\
\bottomrule
\end{tabular}
\end{center}

\begin{definizione}[Gruppo ciclico]
Un gruppo $(G, *)$ è un \emph{gruppo ciclico} se esiste un elemento $g \in G$ tale che $\forall a \in G$
\[
  a = \underbrace{g * g * \dots * g}_{\text{qualche volta}}
  \quad\text{oppure}\quad a = e\ (= g^0)
  \quad\text{oppure}\quad a = g^{-1} * g^{-1} * \dots * g^{-1} .
\]
\end{definizione}

\begin{esempio}[$\mathbb{Z}_{12}$: l'aritmetica dell'orologio]
$\mathbb{Z}_{12} = \{[1], [2], \dots, [12]\}$ con la somma ad esempio $[3] + [4] = [7]$, $[8] + [5] = [1]$. Si può costruire la tavola della somma, con $[12]$ elemento neutro e ogni elemento invertibile (ad esempio $[5] + [?] = [12]$ dà $[?] = [7]$). Poiché $[2] = [1] + [1]$, $[3] = [1] + [1] + [1]$, \dots, il gruppo è ciclico, generato da $[1]$.
\end{esempio}

\begin{esempio}[$(\mathbb{Z}, +)$ è ciclico]
$(\mathbb{Z}, +)$ è un gruppo ciclico: $+1$ genera tutti i numeri. Infatti, se $h \in \mathbb{Z}$:
\begin{itemize}
  \item $h > 0$: $h = \underbrace{1 + 1 + \dots + 1}_{h\text{ volte}}$;
  \item $h = 0$: elemento neutro;
  \item $h < 0$: $h = \underbrace{(-1) + (-1) + \dots + (-1)}_{(-h)\text{ volte}}$.
\end{itemize}
\end{esempio}

\subsubsection{Simmetrie di un rettangolo}

\begin{esempio}
Consideriamo un rettangolo $ABCD$ (con $A$, $B$ in basso e $D$, $C$ in alto) e le sue simmetrie:
\begin{itemize}
  \item $a_1$: simmetria rispetto all'asse verticale, $(ABCD) \mapsto (BADC)$;
  \item $a_2$: simmetria rispetto all'asse orizzontale, $(ABCD) \mapsto (DCBA)$;
  \item $\tau$: rotazione di $180^\circ$, $(ABCD) \mapsto (CDAB)$;
  \item $\mathrm{Id}$: identità, $(ABCD) \mapsto (ABCD)$.
\end{itemize}
La composizione di simmetrie è ancora una simmetria del rettangolo. Ad esempio
\[
  \mathrm{Id} * \tau\,(ABCD) = \mathrm{Id}(CDAB) = (CDAB), \qquad
  \tau * \tau\,(ABCD) = \tau(CDAB) = (ABCD),
\]
\[
  \tau * a_1\,(ABCD) = \tau(BADC) = (DCBA) = a_2\,(ABCD).
\]
Dunque $\tau * a_1 = a_2$ (il resto della tavola della composizione viene completato nella lezione successiva).
\end{esempio}

\section[Simmetrie, anelli e campi (Lezione 10 -- 20 ottobre 2026)]{Simmetrie, anelli e campi}
\noindent\emph{Lezione 10 -- 20 ottobre 2026}
\subsection{Il gruppo delle simmetrie del rettangolo}

Riprendiamo le simmetrie del rettangolo $ABCD$: $\mathrm{Id}$, $\tau$ (rotazione di $180^\circ$), $a_1$ (simmetria rispetto all'asse verticale) e $a_2$ (simmetria rispetto all'asse orizzontale):
\[
  \begin{gathered}
  (ABCD) \xrightarrow{a_1} (BADC), \qquad
  (ABCD) \xrightarrow{a_2} (DCBA), \\
  (ABCD) \xrightarrow{\tau} (CDAB), \qquad
  (ABCD) \xrightarrow{\mathrm{Id}} (ABCD).
  \end{gathered}
\]
La tavola della composizione è
\[
\begin{array}{c|cccc}
  * & \mathrm{Id} & \tau & a_1 & a_2 \\ \hline
  \mathrm{Id} & \mathrm{Id} & \tau & a_1 & a_2 \\
  \tau & \tau & \mathrm{Id} & a_2 & a_1 \\
  a_1 & a_1 & a_2 & \mathrm{Id} & \tau \\
  a_2 & a_2 & a_1 & \tau & \mathrm{Id}
\end{array}
\]
Ad esempio $a_1 * a_2\,(ABCD) = a_1\big(a_2(ABCD)\big) = a_1(DCBA) = (CDAB)$, cioè $a_1 * a_2 = \tau$.

\begin{esempio}[Gruppo \emph{non} ciclico]
Questo è un gruppo non ciclico: vale $\tau * \tau = \mathrm{Id}$, $a_1 * a_1 = \mathrm{Id}$, $a_2 * a_2 = \mathrm{Id}$, quindi nessun elemento, composto con se stesso, genera tutti gli altri.
\end{esempio}

\begin{esempio}[Simmetrie del triangolo equilatero]
Sia $ABC$ un triangolo equilatero. Le sue simmetrie sono le tre riflessioni $a_1, a_2, a_3$ rispetto agli assi, la rotazione $\tau$ di $120^\circ$, la rotazione $\tau^2$ di $240^\circ$ e l'identità. Ad esempio
\[
  (ABC) \xrightarrow{a_1} (ACB), \qquad (ABC) \xrightarrow{\tau} (CAB), \qquad \tau * \tau = \tau^2 .
\]
La tavola della composizione ($\mathrm{Id}, a_1, a_2, a_3, \tau, \tau^2$) è stata solo impostata a lezione.
\end{esempio}

\subsection{Anelli}

\begin{definizione}[Anello]
Una struttura $(A, \oplus, \otimes)$ è un \emph{anello} se
\begin{enumerate}
  \item $(A, \oplus)$ è un gruppo abeliano (con $0$ elemento neutro per $\oplus$ e, per $a \in A$, l'opposto indicato con $-a$);
  \item valgono le \emph{distributive}:
  \[
    (a \oplus b) \otimes c = (a \otimes c) \oplus (b \otimes c), \qquad
    a \otimes (b \oplus c) = (a \otimes b) \oplus (a \otimes c);
  \]
  \item $(A, \otimes)$ è un semigruppo ($\otimes$ è associativa).
\end{enumerate}
\end{definizione}

\begin{itemize}
  \item Se inoltre esiste l'elemento neutro per $\otimes$ (indicato con $1$) l'anello è detto \emph{unitario} (anello con unità).
  \item Se $\otimes$ è commutativa, l'anello è detto \emph{anello commutativo}.
\end{itemize}

\begin{definizione}[Corpo e campo]
Sia $(A, \oplus, \otimes)$ un anello unitario.
\begin{itemize}
  \item È detto \emph{corpo} se $\forall a \in A$, $a \neq 0$, esiste l'inverso $a^{-1}$, con $a^{-1} \otimes a = a \otimes a^{-1} = 1$.
  \item È detto \emph{campo} se $\forall a \in A$, $a \neq 0$, esiste l'inverso $a^{-1}$ e inoltre $(A, \otimes)$ è commutativa.
\end{itemize}
\end{definizione}

\begin{esempio}[$2\mathbb{Z}$]
$2\mathbb{Z} = \{2k \mid k \in \mathbb{Z}\} = \{\dots, -4, -2, 0, 2, 4, \dots\}$. È $(2\mathbb{Z}, +, \cdot)$ un anello?
\begin{enumerate}
  \item $(2\mathbb{Z}, +)$ è un gruppo abeliano? $+$ è associativa; esiste l'elemento neutro $0$; $\forall 2k \in 2\mathbb{Z}$, $-2k$ è l'opposto; $+$ è commutativa. Sì.
  \item Le distributive valgono.
  \item Il prodotto è associativo e commutativo.
\end{enumerate}
Dunque $(2\mathbb{Z}, +, \cdot)$ è un anello commutativo (senza unità).
\end{esempio}

\begin{esempio}
\begin{itemize}
  \item $(\mathbb{R}, +, \cdot)$ e $(\mathbb{Q}, +, \cdot)$ sono \emph{campi}.
  \item $(\mathbb{Z}, +, \cdot)$ è un anello commutativo unitario, ma non è un campo (gli unici elementi invertibili sono $+1$ e $-1$).
  \item $\mathbb{Z}_2 = \{0, 1\}$ con le operazioni
  \[
  \begin{array}{c|cc}
    + & 0 & 1 \\ \hline
    0 & 0 & 1 \\
    1 & 1 & 0
  \end{array}
  \qquad
  \begin{array}{c|cc}
    \cdot & 0 & 1 \\ \hline
    0 & 0 & 0 \\
    1 & 0 & 1
  \end{array}
  \]
  Per $+$ l'elemento neutro è $0$, per $\cdot$ è $1$: $(\mathbb{Z}_2, +, \cdot)$ è un campo.
\end{itemize}
\end{esempio}

\section[Il campo dei numeri complessi (Lezione 10 -- 20 ottobre 2026)]{Il campo dei numeri complessi}
\noindent\emph{Lezione 10 -- 20 ottobre 2026}
\subsection{Il campo dei numeri complessi}

Consideriamo $\mathbb{R}^2 = \mathbb{R}\times\mathbb{R} = \{(x, y) \mid x, y \in \mathbb{R}\}$ con le operazioni
\begin{align*}
  (x_1, y_1) + (x_2, y_2) &= (x_1 + x_2,\ y_1 + y_2),\\
  (x_1, y_1)\cdot(x_2, y_2) &= (x_1 x_2 - y_1 y_2,\ x_1 y_2 + x_2 y_1).
\end{align*}
$(\mathbb{R}^2, +, \cdot) = \mathbb{C}$ è l'\emph{insieme dei numeri complessi}; $z = (x, y)$.

\subsubsection{Proprietà}

\begin{enumerate}
  \item \emph{$(\mathbb{C}, +)$ è un gruppo abeliano.}
  \begin{itemize}
    \item $z_1 + z_2 = z_2 + z_1$ (commutativa);
    \item $z_1 + (z_2 + z_3) = (z_1 + z_2) + z_3$ (associativa);
    \item esiste $0 = (0, 0)$ elemento neutro: $(x, y) + (0, 0) = (0, 0) + (x, y) = (x, y)$;
    \item per ogni $z = (x, y)$ esiste l'inverso (opposto) $-z = (-x, -y)$: $(x, y) + (-x, -y) = (x - x, y - y) = (0, 0) = 0$.
  \end{itemize}
  \item \emph{Distributive:} $z_1\cdot(z_2 + z_3) = z_1 z_2 + z_1 z_3$ e $(z_1 + z_2)\cdot z_3 = z_1 z_3 + z_2 z_3$.
  \item \emph{Associativa del prodotto:} $z_1\cdot(z_2\cdot z_3) = (z_1\cdot z_2)\cdot z_3$. Quindi $(\mathbb{C}, +, \cdot)$ è un anello.
  \item \emph{Commutativa del prodotto:} $z_1 \cdot z_2 = z_2 \cdot z_1$. Quindi è un anello commutativo.
  \item \emph{Elemento neutro del prodotto:} $1 = (1, 0)$ è l'unità per $\cdot$:
  \[
    (1, 0)\cdot(x, y) = (1\cdot x - 0\cdot y,\ 1\cdot y + 0\cdot x) = (x, y).
  \]
  Quindi $(\mathbb{C}, +, \cdot)$ è un anello commutativo con unità.
  \item \emph{Inverso moltiplicativo:} per ogni $z = (x, y) \neq (0, 0)$ esiste l'inverso
  \[
    z^{-1} = \left( \frac{x}{x^2 + y^2},\ -\frac{y}{x^2 + y^2} \right).
  \]
  Infatti
  \[
    (x, y)\cdot\left( \frac{x}{x^2+y^2},\ \frac{-y}{x^2+y^2} \right)
    = \left( \frac{x^2}{x^2+y^2} - \frac{y(-y)}{x^2+y^2},\ \frac{x(-y)}{x^2+y^2} + \frac{y\,x}{x^2+y^2} \right) = (1, 0) = 1 .
  \]
\end{enumerate}
Dunque $(\mathbb{C}, +, \cdot)$ è un \emph{campo}.

\subsubsection{Scrittura $a + ib$}

Per semplicità scriviamo $(1, 0) = 1$ e $(0, 1) = i$ (nuovo simbolo, \emph{unità immaginaria}). Allora
\[
  (x, y) = (x, 0) + (0, y) = x\,(1, 0) + y\,(0, 1) = x + yi .
\]
I numeri $(x, 0)$ sono i numeri reali: $(x_1, 0) + (x_2, 0) = (x_1 + x_2, 0)$ e $(x_1, 0)\cdot(x_2, 0) = (x_1 x_2, 0)$; quindi $\mathbb{R} \subseteq \mathbb{C}$.

\begin{osservazione}
$i^2 = i\cdot i = (0, 1)\cdot(0, 1) = (-1, 0) = -(1, 0) = -1$. Anche $(-i)^2 = (-i)(-i) = (0, -1)(0, -1) = (-1, 0) = -1$.
\end{osservazione}

L'equazione $x^2 + 1 = 0$ con $x \in \mathbb{R}$ non ha soluzioni. L'equazione $z^2 + 1 = 0$ con $z \in \mathbb{C}$, cioè $z^2 = -1$, ha soluzioni $z_1 = i$, $z_2 = -i$.

\begin{esempio}[Equazione di secondo grado su $\mathbb{C}$]
\[
  z^2 - 2z + 10 = 0, \qquad
  z_{1,2} = \frac{2 \pm \sqrt{4 - 40}}{2} = \frac{2 \pm \sqrt{-36}}{2} = \frac{2 \pm \sqrt{36\cdot(-1)}}{2} = \frac{2 \pm 6\sqrt{-1}}{2} ,
\]
cioè $z_1 = \dfrac{2 + 6i}{2} = 1 + 3i$ e $z_2 = \dfrac{2 - 6i}{2} = 1 - 3i$.
\end{esempio}

Possiamo scrivere i numeri complessi in questo modo: $z = a + ib$, $w = c + id$ con $a, b, c, d \in \mathbb{R}$; $a = a(1, 0)$ è la \emph{parte reale}, $b = b(0, 1)$ la \emph{parte immaginaria}. Allora
\[
  z + w = (a + c) + i(b + d),
\]
\[
  z \cdot w = (a + ib)(c + id) = ac + iad + ibc + i^2 bd = ac - bd + (ad + bc)\,i ,
\]
dove si è usato $i^2 = -1$.

\subsection{Riepilogo di esempi di strutture}

\begin{itemize}
  \item $(\mathbb{N}\setminus\{0\}, +)$: semigruppo (non monoide).
  \item $(\mathbb{N}, +)$: monoide (non gruppo).
  \item $(\mathbb{Z}, +)$: gruppo abeliano.
  \item $(2\mathbb{Z}, +, \cdot)$: anello commutativo non unitario.
  \item $(\mathbb{Z}, +, \cdot)$: anello unitario commutativo (non campo).
  \item $(\mathbb{R}, +, \cdot)$, $(\mathbb{Q}, +, \cdot)$, $(\mathbb{C}, +, \cdot)$: campi; anche $(\mathbb{Z}_2, +, \cdot)$ è un campo.
\end{itemize}

\section[Proprietà degli anelli e domini di integrità (Lezione 10 -- 20 ottobre 2026)]{Proprietà degli anelli e domini di integrità}
\noindent\emph{Lezione 10 -- 20 ottobre 2026}
\subsection{Proprietà degli anelli}

Sia $(A, \oplus, \otimes)$ un anello e $a, b, c \in A$. Valgono le seguenti proprietà:
\begin{enumerate}
  \item $a \oplus b = a \oplus c \iff b = c$ (legge di cancellazione per $\oplus$);
  \item $a \otimes 0 = 0 \otimes a = 0$;
  \item $(-a)\otimes b = a \otimes (-b) = -(a \otimes b)$;
  \item $(-a)\otimes(-b) = a \otimes b$;
  \item $(a \ominus b)\otimes c = a \otimes c \ominus b \otimes c$ e $a \otimes (b \ominus c) = a \otimes b \ominus a \otimes c$ (dove $a \ominus b = a \oplus (-b)$);
  \item $(na)\otimes b = n\,(a \otimes b) = a \otimes (nb)$ per $n > 0$, dove $na = \underbrace{a \oplus a \oplus \dots \oplus a}_{n\text{ volte}}$;
  \item $(na)\otimes(mb) = (nm)(a \otimes b)$ per $n, m \in \mathbb{Z}$ (con $ka = \underbrace{(-a)\oplus\dots\oplus(-a)}_{|k|\text{ volte}}$ per $k$ negativo);
  \item se esiste l'unità $1 \in A$, allora $\forall a \in A$: $-a = (-1)\otimes a$.
\end{enumerate}

\subsection{Domini di integrità e divisori dello zero}

\begin{definizione}[Legge di annullamento del prodotto]
Sia $(A, \oplus, \otimes)$ un anello. Vale la \emph{legge di annullamento del prodotto} se
\[
  \forall a, b \in A \qquad a \otimes b = 0 \implies (a = 0 \ \text{oppure}\ b = 0).
\]
\end{definizione}

\begin{definizione}[Dominio di integrità]
Un anello $(A, \oplus, \otimes)$ è detto \emph{dominio di integrità} se vale la legge di annullamento del prodotto.
\end{definizione}

\begin{definizione}[Divisore dello zero]
Sia $(A, \oplus, \otimes)$ un anello. Un elemento $a \in A$ è detto \emph{divisore dello zero} se $a \neq 0$ e $\exists b \neq 0$, $b \in A$, tale che $a \otimes b = 0$ oppure $b \otimes a = 0$.
\end{definizione}

\begin{proposizione}
Sia $(A, \oplus, \otimes)$ un anello. Sono equivalenti:
\begin{enumerate}
  \item in $A$ vale la legge di annullamento del prodotto;
  \item in $A$ non ci sono divisori dello zero;
  \item l'insieme $A\setminus\{0\}$ è stabile per $\otimes$ (cioè $\otimes$ è un'operazione interna a $A\setminus\{0\}$).
\end{enumerate}
\end{proposizione}

\begin{proof}
$1) \Rightarrow 2)$. Per assurdo sia $a$ un divisore dello zero, cioè $a \neq 0$ ed esiste $b \neq 0$ tale che $a \otimes b = 0$. Per la legge di annullamento del prodotto allora $a = 0$ (impossibile) oppure $b = 0$ (impossibile): assurdo. Dunque non esistono divisori dello zero.

$2) \Rightarrow 3)$. Siano $a, b \in A\setminus\{0\}$; possiamo avere $a \otimes b = 0$? Se fosse $a \otimes b = 0$ allora $a \neq 0$, $b \neq 0$ e $a$ sarebbe un divisore dello zero, ma non esistono per ipotesi. Quindi $a \otimes b \neq 0$, cioè $a \otimes b \in A\setminus\{0\}$, cioè $\otimes$ è interna ad $A\setminus\{0\}$.

$3) \Rightarrow 1)$. Siano $a, b \in A$ tali che $a \otimes b = 0$. Per assurdo siano $a \neq 0$ e $b \neq 0$, cioè $a, b \in A\setminus\{0\}$; poiché $\otimes$ è interna a $A\setminus\{0\}$ si avrebbe $a \otimes b \neq 0$: impossibile. Quindi $a = 0$ oppure $b = 0$, cioè vale la legge di annullamento del prodotto.
\end{proof}

\begin{osservazione}
\begin{itemize}
  \item $(\mathbb{Z}, +, \cdot)$, $(\mathbb{Q}, +, \cdot)$, $(\mathbb{R}, +, \cdot)$ sono domini di integrità.
  \item In un campo (o in un corpo) vale la legge di annullamento del prodotto.
  \item In un dominio di integrità vale: se $a \neq 0$, allora $a \otimes b = a \otimes c \iff b = c$ (legge di cancellazione del prodotto).
\end{itemize}
\end{osservazione}

\begin{proof}[Dimostrazione della cancellazione]
$\Leftarrow$ Vale sempre. $\Rightarrow$ Da $a \otimes b = a \otimes c$ segue $0 = a \otimes c - a \otimes b = a \otimes (c - b)$. Per la legge di annullamento del prodotto $a = 0$ oppure $c - b = 0$; ma $a \neq 0$, quindi non può succedere la prima, dunque $c - b = 0$, cioè $c = b$.
\end{proof}

\subsubsection{Notazione per i campi}

Si indica un campo con $(\mathbb{K}, +, \cdot)$: $(\mathbb{K}, +, \cdot)$ è un anello commutativo con unità, e posto $\mathbb{K}^* = \mathbb{K}\setminus\{0\}$, ogni $k \in \mathbb{K}^*$ ha inverso $k^{-1}$. In un campo l'inverso moltiplicativo è unico.

\begin{esempio}[Divisori dello zero in $\mathbb{Z}_{12}$]
In $\mathbb{Z}_{12}$ la tavola del prodotto è
\[
\begin{array}{c|cccccc}
  \cdot & [1] & [2] & [3] & [4] & \cdots & [12] \\ \hline
  \relax[1] & [1] & [2] & [3] & [4] & \cdots & [12] \\
  \relax[2] & [2] & [4] & [6] & [8] & \cdots & [12] \\
  \relax[3] & [3] & [6] & [9] & [12] & \cdots & [12]
\end{array}
\]
$(\mathbb{Z}_{12}, +, \cdot)$ è un anello commutativo con unità; $[12]$ è l'elemento neutro per la somma, cioè $[0]$. Si ha $[3]\cdot[4] = [12] = [0]$, ma $[3] \neq [12]$ e $[4] \neq [12]$: non vale la legge di annullamento del prodotto e $[3]$, $[4]$ sono divisori dello zero.

Inoltre $[2]\cdot[1] = [2] = [2]\cdot[7]$ ma $[1] \neq [7]$: non vale la cancellazione del prodotto.
\end{esempio}

\section[Interi, divisibilità e divisione con resto (Lezione 11 -- 26 ottobre 2026)]{Interi, divisibilità e divisione con resto}
\noindent\emph{Lezione 11 -- 26 ottobre 2026}
\subsection{Gli interi come struttura algebrica}

$\mathbb{Z} = \{0, \pm 1, \pm 2, \pm 3, \dots\}$ con la somma e il prodotto di interi.
\begin{itemize}
  \item $(\mathbb{Z}, +)$ è un gruppo abeliano: $0$ elemento neutro, l'opposto di $k \in \mathbb{Z}$ è $-k \in \mathbb{Z}$.
  \item $(\mathbb{Z}, \cdot)$ è un semigruppo commutativo, con $+1$ elemento neutro per $\cdot$ ($1$ è l'unità); $+1$ e $-1$ sono invertibili per $\cdot$.
  \item Valgono le distributive.
\end{itemize}
Quindi $(\mathbb{Z}, +, \cdot)$ è un anello commutativo con unità. Vale la legge di annullamento del prodotto: è un dominio di integrità. In particolare, se $k \neq 0$ e $k\cdot x = k\cdot y \Rightarrow x = y$ ($k, x, y \in \mathbb{Z}$): vale la legge di cancellazione del prodotto. Indicheremo $(\mathbb{Z}, +, \cdot)$ semplicemente con $\mathbb{Z}$.

\subsection{Divisibilità}

\begin{definizione}
Siano $a, b \in \mathbb{Z}$. Si dice che $a$ \emph{divide} $b$, e si scrive $a \mid b$, se $\exists k \in \mathbb{Z}$ tale che $b = k\cdot a$. Se $a$ non divide $b$ ($\nexists k$ tale che $b = ka$) si scrive $a \nmid b$.
\end{definizione}

\begin{esempio}
$5 \nmid 7$; \quad $3 \mid 33$ perché $3\cdot 11 = 33$; \quad $-4 \mid 20$ perché $\exists k = -5$ tale che $(-4)\cdot(-5) = 20$; \quad $7 \mid -21$ perché $\exists -3 \in \mathbb{Z}$ tale che $7\cdot(-3) = -21$.
\end{esempio}

\begin{osservazione}
\begin{itemize}
  \item $\forall a \in \mathbb{Z}$: $a \mid 0$, perché $0\cdot a = 0$.
  \item $\forall a \in \mathbb{Z}$: $1 \mid a$, perché $a\cdot 1 = a$.
  \item $0 \mid b \Rightarrow \exists k$ tale che $0\cdot k = b \Rightarrow b = 0$: cioè $0$ divide solo zero.
  \item $\forall a$: $a \mid a$, perché $a\cdot 1 = a$.
  \item La relazione $\mid$ su $\mathbb{N}$ è una relazione d'ordine parziale. Su $\mathbb{Z}$ non è una relazione d'ordine perché non è antisimmetrica: se $a \mid b \wedge b \mid a$ allora $a = b$ oppure $a = -b$ (ad esempio $5 \mid -5 \wedge -5 \mid 5$ ma $-5 \neq 5$).
  \item $a \mid b \wedge a \mid c \Rightarrow a \mid (b + c)$; inoltre $a \mid (b + c) \wedge a \mid b \Rightarrow a \mid c$.
\end{itemize}
\end{osservazione}

\begin{proposizione}[Lemma]
Se $a \mid b$ e $a \mid c$ allora $a \mid (\alpha b + \beta c)$ per ogni $\alpha, \beta \in \mathbb{Z}$.
\end{proposizione}

\begin{proof}
$a \mid b$ significa $\exists k_1 \in \mathbb{Z}$ tale che $b = k_1\cdot a$; $a \mid c$ significa $\exists k_2 \in \mathbb{Z}$ tale che $c = k_2\cdot a$. Allora
\[
  \alpha b + \beta c = \alpha k_1 a + \beta k_2 a = (\alpha k_1 + \beta k_2)\cdot a = h\cdot a ,
\]
con $h = \alpha k_1 + \beta k_2 \in \mathbb{Z}$. Cioè $a \mid (\alpha b + \beta c)$.
\end{proof}

\subsection{Divisione con resto}

\begin{proposizione}
Siano $a, b \in \mathbb{Z}$, $b \neq 0$. Allora esistono e sono unici $q, r \in \mathbb{Z}$ tali che
\[
  a = q\cdot b + r, \qquad 0 \le r < |b| ,
\]
dove $q$ è il \emph{quoziente} e $r$ il \emph{resto}.
\end{proposizione}

\begin{proof}[Dimostrazione dell'esistenza]
Sia $T = \{k \in \mathbb{Z} \mid k = a - mb,\ m \in \mathbb{Z},\ k \ge 0\} \subseteq \mathbb{N}$.

$T \neq \emptyset$? Se $a \ge 0$, allora $a \in T$ (con $m = 0$). Se $a < 0$, esiste $\bar m \in \mathbb{Z}$ tale che $\bar m b \le a$ (ad esempio $\bar m = -|a|\,b$, perché $\bar m b = -|a|\,b^2 \le -|a| = a$), cioè $0 \le a - \bar m b$, cioè $a - \bar m b \in T$. Quindi $T \neq \emptyset$.

$\mathbb{N}$ è bene ordinato, quindi $T$ ha minimo. Sia $r = \min T$. Poiché $r \in T$, $r = a - qb$ per qualche $q \in \mathbb{Z}$, cioè $a = qb + r$, e $r \ge 0$. Resta da far vedere che $r < |b|$.

Per assurdo sia $r \ge |b|$, cioè $r = |b| + r_0$ con $r_0 \ge 0$. Allora
\[
  0 \le r_0 = r - |b| = a - qb - |b| = \begin{cases} a - b(q + 1) & b > 0 \\ a - b(q - 1) & b < 0 \end{cases}
\]
cioè $r_0 \in T$ e $r_0 = r - |b| < r$, contro la minimalità di $r$: assurdo. Dunque $r < |b|$.
\end{proof}

\begin{proof}[Dimostrazione dell'unicità]
Supponiamo che ci siano due coppie $q, r$ e $q', r'$ tali che $a = qb + r = q'b + r'$ con $0 \le r, r' < |b|$. Supponiamo $r \le r'$. Allora
\[
  0 \le r' - r = (a - q'b) - (a - qb) = b(q - q') .
\]
Passando ai valori assoluti:
\[
  0 \le |q - q'|\,|b| = |r' - r| = r' - r \le r' < |b| ,
\]
cioè $|q - q'|\,|b| < |b|$, cioè $|q - q'| < 1$, cioè $|q - q'| = 0$, cioè $q' = q$. Da $qb + r = q'b + r'$ segue allora $r' = r$.
\end{proof}

\begin{esempio}
\begin{itemize}
  \item $a = 15$, $b = 6$: $15 = 6\cdot 2 + 3$ ($q = 2$, $r = 3$).
  \item $a = 15$, $b = -6$: $15 = (-6)\cdot(-2) + 3$.
  \item $a = -15$, $b = 6$: $-15 = 6\cdot(-2) - 3 = 6\cdot(-2) - 3 + 6 - 6 = 6\cdot(-3) + 3$, cioè quoziente $-3$ e resto $3$.
  \item $a = -15$, $b = -6$: $-15 = (-6)\cdot 2 - 3 = (-6)\cdot 3 + 3$.
  \item $a = -24$, $b = 7$: $-24 = 7\cdot(-4) + 4$.
  \item $a = 24$, $b = -7$: $24 = (-7)\cdot(-3) + 3$.
\end{itemize}
\end{esempio}

\subsection{Classi di resto modulo 7}

$a \equiv b \bmod 7$ se $7 \mid a - b$, cioè se $\exists k \in \mathbb{Z}$ tale che $a - b = 7k$, cioè $a = 7k + b$. Le classi sono
\begin{align*}
  [0]_7 &= \{0, \pm 7, \pm 14, \dots\},\\
  [1]_7 &= \{1,\ 7 + 1,\ 14 + 1, \dots, 7k + 1, \dots, -7 + 1,\ -14 + 1, \dots\},\\
  [2]_7 &= \{2,\ 7 + 2, \dots, 7k + 2, \dots, -7 + 2,\ -14 + 2, \dots\},\\
  [-1]_7 &= \{-1 = (-1)\cdot 7 + 6,\ 6,\ 6 + 7k, \dots\} = [6]_7 .
\end{align*}
Le classi di congruenza modulo $n$ di $\mathbb{Z}_n = \mathbb{Z}/{\equiv_n}$ sono esattamente $n$, cioè i resti possibili della divisione per $n$: $[0], [1], \dots, [n-1]$.

\section[MCD, mcm e algoritmo di Euclide (Lezione 11 -- 26 ottobre 2026)]{MCD, mcm e algoritmo di Euclide}
\noindent\emph{Lezione 11 -- 26 ottobre 2026}
\subsection{Massimo comune divisore}

\begin{definizione}
Siano $a, b \in \mathbb{Z}$ non entrambi nulli. Un elemento $d \in \mathbb{Z}$ è \emph{un massimo comun divisore} tra $a$ e $b$ se
\begin{enumerate}
  \item[i)] $d \mid a$ e $d \mid b$;
  \item[ii)] $\forall d' \in \mathbb{Z}$ tale che $d' \mid a$ e $d' \mid b$ si ha $d' \mid d$.
\end{enumerate}
\end{definizione}

\begin{osservazione}
In generale, se $d$ è un massimo comun divisore tra $a$ e $b$, anche $-d$ è un massimo comun divisore tra $a$ e $b$. Chiamiamo $\operatorname{MCD}(a, b)$ quello \emph{positivo}.
\end{osservazione}

\begin{esempio}
$\operatorname{MCD}(48, 15) = 3$, \quad $\operatorname{MCD}(3, 7) = 1$, \quad $\operatorname{MCD}(-12, 36) = 12$.
\end{esempio}

\begin{proprieta}
\begin{itemize}
  \item $\operatorname{MCD}(a, b) = \operatorname{MCD}(b, a) = \operatorname{MCD}(|a|, |b|)$.
  \item $\operatorname{MCD}(ab, ac) = |a|\operatorname{MCD}(b, c)$.
  \item $\operatorname{MCD}(0, a) = |a|$: infatti $d \mid 0$ per ogni $d$ (ogni $d$ divide $0$), e $d \mid a$; quindi $d$ è un divisore di $a$ e il più grande possibile è $|a|$.
  \item $\operatorname{MCD}(0, 0)$ \emph{non è definito}.
\end{itemize}
\end{proprieta}

\begin{definizione}[Numeri coprimi]
Due numeri $a, b$ sono \emph{coprimi} (o \emph{relativamente primi}, o \emph{primi tra loro}) se $\operatorname{MCD}(a, b) = 1$.
\end{definizione}

\begin{esempio}
$3, 5$ sono coprimi; $11$ e $36$ sono coprimi.
\end{esempio}

\subsection{Minimo comune multiplo}

\begin{definizione}
Siano $a, b \in \mathbb{Z}$. Un elemento $m \in \mathbb{Z}$ è un \emph{minimo comune multiplo} di $a$ e $b$ se
\begin{enumerate}
  \item[i)] $a \mid m$ e $b \mid m$;
  \item[ii)] $\forall m' \in \mathbb{Z}$ tale che $a \mid m'$ e $b \mid m'$ si ha $m \mid m'$.
\end{enumerate}
\end{definizione}

Se $m$ è un minimo comune multiplo di $a$ e $b$, anche $-m$ lo è; indichiamo con $\operatorname{mcm}(a, b)$ quello \emph{positivo}.

\begin{esempio}
$\operatorname{mcm}(4, 6) = 12$, \quad $\operatorname{mcm}(-7, 9) = 63$.
\end{esempio}

\begin{proprieta}
$\operatorname{mcm}(a, b) = \operatorname{mcm}(b, a) = \operatorname{mcm}(|a|, |b|)$ e $\operatorname{mcm}(0, b) = 0$.
\end{proprieta}

\begin{proprieta}[Formula]
Per $a, b$ non entrambi nulli:
\[
  \operatorname{MCD}(a, b)\cdot\operatorname{mcm}(a, b) = |a|\cdot|b| .
\]
\end{proprieta}

\subsection{Algoritmo euclideo delle divisioni successive}

\begin{teorema}[Algoritmo di Euclide]
Siano $a, b \in \mathbb{Z}$, $a \ge b > 0$. Consideriamo queste uguaglianze (divisioni con resto):
\begin{align*}
  a &= q_1 b + r_1, & 0 &< r_1 < b,\\
  b &= q_2 r_1 + r_2, & 0 &< r_2 < r_1,\\
  r_1 &= q_3 r_2 + r_3, & 0 &< r_3 < r_2,\\
  &\ \,\vdots\\
  r_i &= q_{i+2}\, r_{i+1} + r_{i+2}, & 0 &< r_{i+2} < r_{i+1},\\
  &\ \,\vdots\\
  r_{n-3} &= q_{n-1}\, r_{n-2} + r_{n-1}, & 0 &< r_{n-1} < r_{n-2},\\
  r_{n-2} &= q_n\, r_{n-1} + r_n, & 0 &< r_n < r_{n-1},\\
  r_{n-1} &= q_{n+1}\, r_n .
\end{align*}
Allora $r_n = \operatorname{MCD}(a, b)$ (l'ultimo resto non nullo).
\end{teorema}

\begin{proof}
L'algoritmo termina in un numero finito di passi perché $0 < r_n < r_{n-1} < \dots < r_1 < b$.

\emph{Condizione i)} del MCD: dall'ultima uguaglianza $r_{n-1} = q_{n+1} r_n$, cioè $r_n \mid r_{n-1}$. Poi
\[
  r_{n-2} = r_{n-1} q_n + r_n = q_{n+1} r_n q_n + r_n \ \Rightarrow\ r_n \mid r_n(q_{n+1} q_n + 1) = r_{n-2}.
\]
Per il lemma, da $r_{n-3} = r_{n-2} q_{n-1} + r_{n-1}$ segue $r_n \mid r_{n-3}$. Risalendo tutte le uguaglianze si ottiene $r_n \mid r_i$ per ogni resto. Dalla seconda uguaglianza $b = r_1 q_2 + r_2$ e dal lemma segue $r_n \mid b$; dalla prima $a = q_1 b + r_1$ segue $r_n \mid a$. Questa è la condizione i) del MCD.

\emph{Condizione ii)}: sia $d$ un altro divisore comune di $a$ e $b$. Da $d \mid a = q_1 b + r_1$ e $d \mid b$ segue, per il lemma, $d \mid r_1$; poi $d \mid b = r_1 q_2 + r_2$ e $d \mid r_1$ implicano $d \mid r_2$; $d \mid r_1 = r_2 q_3 + r_3$ implica $d \mid r_3$, e così via. Quindi $d \mid r_i$ per ogni $i$, in particolare $d \mid r_{n-2}$ e $d \mid r_{n-1}$, da cui, per il lemma, $d \mid r_n = r_{n-2} - q_n r_{n-1}$. Cioè $d \mid r_n$: è la condizione ii) del MCD.
\end{proof}

\begin{esempio}[$\operatorname{MCD}(824, 376)$]
\begin{align*}
  824 &= 376\cdot 2 + 72,\\
  376 &= 72\cdot 5 + 16,\\
  72 &= 16\cdot 4 + 8,\\
  16 &= 8\cdot 2 + 0 .
\end{align*}
Dunque $\operatorname{MCD}(824, 376) = 8$ (l'ultimo resto non nullo).
\end{esempio}

\section[Numeri primi e anello delle classi di resto (Lezione 13 -- 2 novembre 2026)]{Numeri primi e anello delle classi di resto}
\noindent\emph{Lezione 13 -- 2 novembre 2026}
\subsection{Conseguenze del teorema di fattorizzazione}

\begin{proposizione}
I numeri primi sono infiniti.
\end{proposizione}

\begin{proof}
Per assurdo siano finiti: $p_1, \dots, p_k$. Consideriamo $m = p_1 p_2 \cdots p_k + 1$. Per ogni $p_i$ si ha $p_i \nmid m$ (il resto della divisione di $m$ per $p_i$ è $1$), ma $m$ ha un divisore primo per il teorema di fattorizzazione: assurdo.
\end{proof}

\begin{proposizione}
Sia $p$ un numero primo ($p > 0$). Allora $\sqrt{p}$ è irrazionale.
\end{proposizione}

\begin{proof}
Dobbiamo dimostrare che $\sqrt{p} \notin \mathbb{Q}$. Per assurdo sia $\sqrt{p} \in \mathbb{Q}$, cioè $\exists a, b \in \mathbb{Z}$, $b \neq 0$, con $\sqrt{p} = \frac{a}{b}$. Allora $b\sqrt{p} = a$ e, elevando al quadrato, $b^2 p = a^2$. Nella fattorizzazione di $a^2$ il primo $p$ ha esponente \emph{pari}; ma nella fattorizzazione di $b^2 p$ il primo $p$ ha esponente \emph{dispari}: assurdo. Quindi $\sqrt{p} \in \mathbb{R}$ ma $\sqrt{p} \notin \mathbb{Q}$.
\end{proof}

\subsection{$\mathbb{Z}_n$: l'anello delle classi di resto}

Siano $a, b \in \mathbb{Z}$, $n > 0$. Si ha $a \equiv b \bmod n \iff n \mid a - b \iff \exists k \in \mathbb{Z}$ tale che $a - b = nk$.

\[
  \mathbb{Z}_n = \{[0]_n, [1]_n, \dots, [n-1]_n\} = \mathbb{Z}/{\equiv_n}
\]
con
\begin{align*}
  [0]_n &= \{h \in \mathbb{Z} \mid 0 \equiv h \bmod n\} = \{kn \mid k \in \mathbb{Z}\},\\
  [1]_n &= \{h \in \mathbb{Z} \mid 1 \equiv h \bmod n\} = \{kn + 1 \mid k \in \mathbb{Z}\},\\
  &\ \,\vdots\\
  [n-1]_n &= \{h \in \mathbb{Z} \mid n-1 \equiv h \bmod n\} = \{kn + n - 1 \mid k \in \mathbb{Z}\}.
\end{align*}

Definiamo le operazioni $+ : \mathbb{Z}_n\times\mathbb{Z}_n \to \mathbb{Z}_n$ e $\cdot : \mathbb{Z}_n\times\mathbb{Z}_n \to \mathbb{Z}_n$ con
\[
  [a]_n + [b]_n = [a + b]_n, \qquad [a]_n\cdot[b]_n = [a\cdot b]_n .
\]

\begin{proposizione}[Lemma]
Siano $a, b, c, d \in \mathbb{Z}$. Se $a \equiv b \bmod n$ e $c \equiv d \bmod n$ allora
\[
  a + c \equiv b + d \bmod n \qquad\text{e}\qquad a\cdot c \equiv b\cdot d \bmod n .
\]
\end{proposizione}

(Il lemma garantisce che le operazioni sono ben definite, cioè non dipendono dai rappresentanti scelti.)

\begin{proposizione}
$(\mathbb{Z}_n, +, \cdot)$ è un anello commutativo con unità. In particolare $[0]_n$ è l'elemento neutro per $+$ in $\mathbb{Z}_n$; l'opposto di $[a]_n$ è la classe $[-a]_n = [n - a]_n$; $[1]_n$ è l'unità per $\cdot$.
\end{proposizione}

\begin{esempio}[in $\mathbb{Z}_{15}$]
\begin{align*}
  [4]_{15} + [10]_{15} &= [4 + 10]_{15} = [14]_{15},\\
  [10]_{15} + [13]_{15} &= [23]_{15} = [8]_{15} \qquad (23 = 1\cdot 15 + 8),\\
  [3]_{15}\cdot[8]_{15} &= [24]_{15} = [9]_{15} \qquad (24 = 1\cdot 15 + 9).
\end{align*}
\end{esempio}

\subsubsection{Tavole di $\mathbb{Z}_2$, $\mathbb{Z}_3$, $\mathbb{Z}_4$}

$\mathbb{Z}_2 = \{[0], [1]\}$ (mod 2):
\[
\begin{array}{c|cc} + & 0 & 1 \\ \hline 0 & 0 & 1 \\ 1 & 1 & 0 \end{array}
\qquad
\begin{array}{c|cc} \cdot & 0 & 1 \\ \hline 0 & 0 & 0 \\ 1 & 0 & 1 \end{array}
\]

$\mathbb{Z}_3 = \{[0]_3, [1]_3, [2]_3\}$ (mod 3):
\[
\begin{array}{c|ccc} + & 0 & 1 & 2 \\ \hline 0 & 0 & 1 & 2 \\ 1 & 1 & 2 & 0 \\ 2 & 2 & 0 & 1 \end{array}
\qquad
\begin{array}{c|ccc} \cdot & 0 & 1 & 2 \\ \hline 0 & 0 & 0 & 0 \\ 1 & 0 & 1 & 2 \\ 2 & 0 & 2 & 1 \end{array}
\]
Si ha ad esempio $[2]_3\cdot[2]_3 = [2\cdot 2]_3 = [4]_3 = [1]_3$ (perché $4 = 3\cdot 1 + 1$).

$\mathbb{Z}_4 = \{[0], [1], [2], [3]\}$:
\[
\begin{array}{c|cccc} + & 0 & 1 & 2 & 3 \\ \hline 0 & 0 & 1 & 2 & 3 \\ 1 & 1 & 2 & 3 & 0 \\ 2 & 2 & 3 & 0 & 1 \\ 3 & 3 & 0 & 1 & 2 \end{array}
\qquad
\begin{array}{c|cccc} \cdot & 0 & 1 & 2 & 3 \\ \hline 0 & 0 & 0 & 0 & 0 \\ 1 & 0 & 1 & 2 & 3 \\ 2 & 0 & 2 & 0 & 2 \\ 3 & 0 & 3 & 2 & 1 \end{array}
\]
$\mathbb{Z}_4$ ha divisori dello zero: $[2]_4\cdot[2]_4 = [0]_4$. Inoltre $[3]$ è invertibile in $\mathbb{Z}_4$: $[3]_4\cdot[3]_4 = [1]_4$.

\section[Invertibili e divisori dello zero modulo $n$ (Lezione 13 -- 2 novembre 2026)]{Invertibili e divisori dello zero modulo $n$}
\noindent\emph{Lezione 13 -- 2 novembre 2026}
\subsection{Divisibilità, invertibili e divisori dello zero in $\mathbb{Z}_n$}

\begin{definizione}
Siano $[a]_n, [b]_n \in \mathbb{Z}_n$.
\begin{itemize}
  \item Diciamo che $[a]_n$ \emph{divide} $[b]_n$ ($[a]_n \mid [b]_n$) se $\exists [c]_n \in \mathbb{Z}_n$ tale che $[b]_n = [a]_n\cdot[c]_n$.
  \item $[a]_n \in \mathbb{Z}_n$ è \emph{invertibile} se $\exists [a']_n \in \mathbb{Z}_n$ tale che $[a]_n\cdot[a']_n = [1]_n$.
  \item $[a]_n \in \mathbb{Z}_n$, $[a]_n \neq [0]_n$, è un \emph{divisore dello zero} se $\exists [b]_n \neq [0]_n$ tale che $[a]_n\cdot[b]_n = [0]_n$.
\end{itemize}
\end{definizione}

\begin{osservazione}
In generale $\mathbb{Z}_n$ non è un dominio di integrità, cioè in $\mathbb{Z}_n$ non vale la legge di annullamento del prodotto: ad esempio $[2]_6\cdot[3]_6 = [6]_6 = [0]_6$.

Non sempre vale la legge di cancellazione del prodotto. Ad esempio in $\mathbb{Z}_{12}$: $[5]\cdot[8] = [40] = [4]$ e $[8]\cdot[8] = [64] = [4]$ ($64 = 12\cdot 5 + 4$), quindi $[5]\cdot[8] = [8]\cdot[8]$ ma $[5] \neq [8]$.

Però $\mathbb{Z}_2$, $\mathbb{Z}_3$, $\mathbb{Z}_5$ sono campi.
\end{osservazione}

\subsection{Proprietà delle congruenze}

\begin{proprieta}[Corollario]
Sia $n > 0$ e $\forall a, b, c \in \mathbb{Z}$. Se $a \equiv b \bmod n$, allora
\[
  a + c \equiv b + c \bmod n, \qquad a\cdot c \equiv b\cdot c \bmod n, \qquad a^i \equiv b^i \bmod n .
\]
In termini di classi: $[a]_n = [b]_n \Rightarrow [a + c]_n = [b + c]_n$, $[ac]_n = [bc]_n$, $[a^i]_n = [b^i]_n$.
\end{proprieta}

\begin{proprieta}
\begin{itemize}
  \item $ac \equiv bc \bmod n$ e $\operatorname{MCD}(c, n) = 1 \Rightarrow a \equiv b \bmod n$ (cioè $[ac]_n = [bc]_n$ e $\operatorname{MCD}(c, n) = 1 \Rightarrow [a]_n = [b]_n$).
  \item $ac \equiv bc \bmod n$ e $\operatorname{MCD}(c, n) = d \Rightarrow a\cdot\frac{c}{d} \equiv b\cdot\frac{c}{d} \bmod \frac{n}{d} \Rightarrow a \equiv b \bmod \frac{n}{d}$.
  \item $a \equiv b \bmod n$ e $p \mid n \Rightarrow a \equiv b \bmod p$.
  \item $a \equiv b \bmod r$ e $a \equiv b \bmod s \Rightarrow a \equiv b \bmod \operatorname{mcm}(r, s)$.
\end{itemize}
\end{proprieta}

\begin{osservazione}[Rappresentanti di $\mathbb{Z}_n$]
Si può rappresentare $\mathbb{Z}_n$ con qualunque insieme di $n$ interi consecutivi. Ad esempio
\[
  \mathbb{Z}_6 = \{[0], [1], [2], [3], [4], [5]\} = \{[7], [8], [9], [10], [11], [12]\} = \{[-2], [-1], [0], [1], [2], [3]\},
\]
in generale $\mathbb{Z}_n = \{[c]_n, [c + 1]_n, \dots, [c + (n-1)]_n\}$. Per esempio in $\mathbb{Z}_6$: $[0]_6 = \{6k \mid k \in \mathbb{Z}\} = [6]_6 = [12]_6 = \{\dots, -6, 0, 6, 12, \dots\}$; $[1]_6 = \{6k + 1\} = [7]_6 = \{\dots, -11, -5, 1, 7, 13, \dots\}$; $[2]_6 = \{6k + 2\} = [8]_6 = [-4]_6 = \{\dots, -10, -4, 2, 8, \dots\}$.
\end{osservazione}

\begin{proposizione}
Sia $c$ tale che $\operatorname{MCD}(c, n) = 1$. Allora
\[
  \mathbb{Z}_n = \{[c\cdot 0], [c\cdot 1], [c\cdot 2], \dots, [c(n-1)]\} .
\]
\end{proposizione}

\begin{esempio}
In $\mathbb{Z}_6$, con $c = 5$ e $\operatorname{MCD}(5, 6) = 1$:
\[
  \mathbb{Z}_6 = \{[5\cdot 0], [5\cdot 1], [5\cdot 2], [5\cdot 3], [5\cdot 4], [5\cdot 5]\} = \{[0], [5], [10], [15], [20], [25]\}.
\]
\end{esempio}

\subsection{Invertibili e divisori dello zero}

\begin{proposizione}
Sia $[a]_n \in \mathbb{Z}_n$, $[a]_n \neq [0]_n$. Allora
\[
  [a]_n \text{ è un divisore dello zero} \iff [a]_n \text{ non è invertibile.}
\]
\end{proposizione}

\begin{proof}
$\Rightarrow$ Sia $[a]$ un divisore dello zero, cioè $\exists [b] \neq [0]$ tale che $[a]\cdot[b] = [0]$. Se $[a]$ fosse invertibile, esisterebbe $[a']$ con $[a]\cdot[a'] = [1]$ e
\[
  [a']\cdot[a]\cdot[b] = [a']\cdot[0] \ \Rightarrow\ [1]\cdot[b] = [b] = [0] ,
\]
ma $[b] \neq [0]$: assurdo.

$\Leftarrow$ Sia $[a]$ non invertibile, cioè $\forall [b] \in \{[0], \dots, [n-1]\}$ si ha $[a]\cdot[b] \neq [1]$. Allora gli $n$ prodotti $[a]\cdot[b]$ assumono al più $n - 1$ valori distinti, quindi devono esserci due classi $[b_1] \neq [b_2]$ tali che $[a]\cdot[b_1] = [a]\cdot[b_2]$, cioè $[ab_1] = [ab_2]$, cioè $ab_1 \equiv ab_2 \bmod n$, cioè $\exists k \in \mathbb{Z}$ tale che $ab_1 - ab_2 = kn$. Cioè $a(b_1 - b_2) \equiv 0 \bmod n$, cioè $[a]\cdot[b_1 - b_2] = [0]$, con $[a] \neq [0]$ per ipotesi e $[b_1 - b_2] \neq [0]$ perché $b_1 \neq b_2$ (come classi). Quindi $[a]$ è un divisore dello zero.
\end{proof}

\begin{proposizione}
In $\mathbb{Z}_n$ la classe $[c]_n$ è invertibile se e solo se $\operatorname{MCD}(c, n) = 1$.
\end{proposizione}

\begin{proof}
$\Rightarrow$ Sia $[c]$ invertibile: $\exists [c'] \in \mathbb{Z}_n$ con $[c]\cdot[c'] = [1]$, cioè $[c\cdot c'] = [1]$, cioè $c\cdot c' - 1 = kn$ per qualche $k$, cioè $c\cdot c' - kn = 1$. Allora $\operatorname{MCD}(c, n) \mid (c\cdot c' - kn) = 1$, cioè $\operatorname{MCD}(c, n) = 1$.

$\Leftarrow$ Sia $\operatorname{MCD}(c, n) = 1$. Per il teorema di Bézout esistono $\alpha, \beta \in \mathbb{Z}$ tali che $1 = \operatorname{MCD}(c, n) = \alpha c + \beta n$, cioè $\alpha c - 1 = -\beta n$, cioè $\alpha c \equiv 1 \bmod n$, cioè $[\alpha]_n\cdot[c]_n = [1]_n$. Quindi $[\alpha]$ è l'inverso di $[c]$ modulo $n$.
\end{proof}

\begin{esempio}[Inverso di $[13]_{30}$ in $\mathbb{Z}_{30}$]
Si usa l'algoritmo di Euclide \emph{esteso} tra $13$ e $30$. Si costruisce la tabella con colonne $q$, $r$, $s$, $t$ (si parte con $r = 30, 13$; $s = 1, 0$; $t = 0, 1$; e ad ogni passo $q$ è il quoziente, $r$ il resto, $s_{\text{nuovo}} = s_{\text{prec-prec}} - q\cdot s_{\text{prec}}$, e analogamente per $t$):
\[
\begin{array}{c|rrr}
  q & r & s & t \\ \hline
  - & 30 & 1 & 0 \\
  - & 13 & 0 & 1 \\
  2 & 4 & 1 & -2 \\
  3 & 1 & -3 & 7 \\
  4 & 0 & 13 & -30
\end{array}
\]
L'ultimo resto non nullo è $1 = \operatorname{MCD}(30, 13)$, con $s = -3$ e $t = 7$:
\[
  1 = 30\cdot(-3) + 13\cdot 7 .
\]
Quindi $[7]_{30}$ è l'inverso di $[13]_{30}$. Verifica: $[7]\cdot[13] = [91] = [1]_{30}$, perché $91 = 3\cdot 30 + 1$.
\end{esempio}

\section[$\mathbb{Z}_p$ e piccolo teorema di Fermat (Lezione 14 -- 3 novembre 2026)]{$\mathbb{Z}_p$ e piccolo teorema di Fermat}
\noindent\emph{Lezione 14 -- 3 novembre 2026}
\subsection{$\mathbb{Z}_p$ con $p$ primo}

\begin{osservazione}
Sia $p$ un numero primo. Allora $[c] \in \mathbb{Z}_p = \{[0], [1], \dots, [p-1]\}$ e, se $[c] \neq [0]$, vale
\[
  \operatorname{MCD}(c, p) = 1 \iff [c] \text{ è invertibile.}
\]
Quindi $(\mathbb{Z}_p, +, \cdot)$ è un \emph{campo}.
\end{osservazione}

\begin{proposizione}
Sia $p$ primo, $p > 0$, e $x, y \in \mathbb{Z}$. Allora
\[
  (x + y)^p \equiv x^p + y^p \pmod p .
\]
\end{proposizione}

\begin{proof}
Per lo sviluppo del binomio di Newton,
\[
  (x + y)^p = \sum_{k=0}^{p} \binom{p}{k} x^k y^{p-k} = x^p + y^p + \sum_{k=1}^{p-1}\binom{p}{k} x^k y^{p-k},
\]
dove si è usato $\binom{p}{p} = 1 = \binom{p}{0}$. Per $1 \le k \le p-1$,
\[
  \binom{p}{k} = \frac{p!}{k!\,(p-k)!} = \frac{p\,(p-1)\cdots(p-k+1)\,(p-k)!}{k!\,(p-k)!} = \frac{p\,(p-1)\cdots(p-k+1)}{k!} .
\]
Il numero $p$ compare al numeratore, ma non divide $k!$ (con $k < p$ e $p$ primo), quindi $p \mid \binom{p}{k}$. Dunque, modulo $p$, tutti i termini della sommatoria sono $\equiv 0$ e
\[
  (x + y)^p \equiv x^p + y^p \pmod p . \qedhere
\]
\end{proof}

\begin{osservazione}
In termini di classi: $\big([x + y]_p\big)^p = \big[(x + y)^p\big]_p = \big[x^p + y^p\big]_p$.
\end{osservazione}

\subsection{Piccolo teorema di Fermat}

\begin{teorema}[Piccolo teorema di Fermat]
Sia $a \in \mathbb{Z}$ e $p$ primo. Allora
\[
  a^p \equiv a \pmod p .
\]
\end{teorema}

\begin{proof}
Sia dapprima $a \ge 0$: si procede per induzione su $a$ (con $p$ primo, $p > 0$).

\emph{Caso base:} $a = 0$: $0^p \equiv 0 \pmod p$.

\emph{Passo induttivo.} Supponiamo vero $a^p \equiv a \pmod p$ (ipotesi induttiva) e dimostriamo $(a+1)^p \equiv a + 1 \pmod p$. Per la proposizione precedente e per l'ipotesi induttiva,
\[
  (a + 1)^p \equiv a^p + 1^p \equiv a + 1 \pmod p ,
\]
quindi il risultato vale per ogni $a \ge 0$.

Se $a < 0$, allora $(-a) > 0$ e vale il risultato per i numeri non negativi. Si ha
\[
  0 = 0^p = \big(a + (-a)\big)^p \equiv a^p + (-a)^p \equiv a^p + (-a) = a^p - a \pmod p ,
\]
cioè $0 \equiv a^p - a \pmod p$, cioè $a \equiv a^p \pmod p$.
\end{proof}

\begin{teorema}[Corollario]
Se $\operatorname{MCD}(a, p) = 1$ allora $a^{p-1} \equiv 1 \pmod p$.
\end{teorema}

\begin{esempio}
Il resto di $8^3$ diviso $3$: per il piccolo teorema di Fermat $8^3 \equiv 8 \equiv 2 \pmod 3$.

$12\,440^3 \equiv\, ?\pmod 3$: poiché $12\,440 = 3\cdot 4146 + 2$, si ha $12\,440^3 \equiv 12\,440 \equiv 2 \pmod 3$.
\end{esempio}

\begin{esempio}[$4^{194} \bmod 7$]
Usiamo il corollario: $\operatorname{MCD}(4, 7) = 1$, quindi $4^6 \equiv 1 \pmod 7$. Dividiamo $194$ per $6$: $194 = 6\cdot 32 + 2$. Allora
\[
  4^{194} = 4^{6\cdot 32 + 2} = \left(4^6\right)^{32}\cdot 4^2 \equiv 1^{32}\cdot 4^2 = 16 \equiv 2 \pmod 7 .
\]
\end{esempio}

\begin{osservazione}
Il corollario ha l'ipotesi $\operatorname{MCD}(a, p) = 1$. Per $a = 6$, $p = 3$: $6^3 \equiv 6 \pmod 3$ (piccolo teorema di Fermat) è vero, ma poiché $\operatorname{MCD}(6, 3) \neq 1$, la congruenza $6^2 \equiv 1 \pmod 3$ è \emph{falsa}: $6^2 = 36 \equiv 0 \pmod 3$.
\end{osservazione}

\section[Gruppo moltiplicativo di $\mathbb{Z}_n$ e funzione di Eulero (Lezione 14 -- 3 novembre 2026)]{Gruppo moltiplicativo di $\mathbb{Z}_n$ e funzione di Eulero}
\noindent\emph{Lezione 14 -- 3 novembre 2026}
\subsection{Il gruppo moltiplicativo di $\mathbb{Z}_n$}

\begin{itemize}
  \item $(\mathbb{Z}_n, \cdot)$ non è mai un gruppo, perché $[0]$ non è invertibile.
  \item $(\mathbb{Z}_p\setminus\{[0]\}, \cdot)$ con $p$ primo è un gruppo (infatti $(\mathbb{Z}_p, +, \cdot)$ è un campo).
  \item $(\mathbb{Z}_n\setminus\{[0]\}, \cdot)$ in generale non è un gruppo.
\end{itemize}
Possiamo però definire
\[
  \mathbb{Z}_n^* = \left\{ [a]_n : \operatorname{MCD}(a, n) = 1 \right\} \subset \mathbb{Z}_n .
\]
$(\mathbb{Z}_n^*, \cdot)$ è un gruppo, detto \emph{gruppo moltiplicativo di $\mathbb{Z}_n$}.

\begin{esempio}
\begin{itemize}
  \item $\mathbb{Z}_p^* = \{[1], \dots, [p-1]\}$, con $p$ primo.
  \item $\mathbb{Z}_6^* = \{[1], [5]\}$.
  \item $\mathbb{Z}_9^* = \{[1], [2], [4], [5], [7], [8]\}$.
\end{itemize}
\end{esempio}

\subsection{La funzione di Eulero}

\begin{definizione}[Funzione di Eulero]
Sia $n \ge 1$. La \emph{funzione di Eulero} $\varphi$ è
\[
  \varphi(n) = |\mathbb{Z}_n^*| = \Big|\big\{ a \mid 1 \le a \le n,\ \operatorname{MCD}(a, n) = 1 \big\}\Big| .
\]
\end{definizione}

\begin{esempio}
$\varphi(6) = 2$, \qquad $\varphi(9) = 6$.
\end{esempio}

\begin{osservazione}
$\varphi(p) = p - 1$ se $p$ è primo, e $\varphi(p^2) = p^2 - p$.
\end{osservazione}

\begin{proposizione}
Sia $p$ primo e $n \ge 1$ intero. Allora
\[
  \varphi(p^n) = p^n - p^{n-1} = p^{n-1}(p - 1) .
\]
\end{proposizione}

\begin{proof}
Tra $1$ e $p^n$ ci sono $p^n$ numeri. Quelli non coprimi con $p^n$ hanno la forma $k\cdot p$, e $k$ può avere tutti i valori da $1$ fino a $p^{n-1}$: ci sono $p^{n-1}$ possibili numeri non coprimi con $p^n$. Quindi $\varphi(p^n) = p^n - p^{n-1}$.
\end{proof}

\begin{esempio}
$\varphi(125) = \varphi(5^3) = 5^3 - 5^2 = 125 - 25 = 100$.
\end{esempio}

\subsection{La biiezione $\mathbb{Z}_{nk} \to \mathbb{Z}_n \times \mathbb{Z}_k$}

Consideriamo la funzione
\[
  \psi : \mathbb{Z}_{nk} \to \mathbb{Z}_n \times \mathbb{Z}_k, \qquad [a]_{nk} \mapsto \big([a]_n, [a]_k\big).
\]

\begin{esempio}
\begin{itemize}
  \item $\psi_1 : \mathbb{Z}_{12} \to \mathbb{Z}_3 \times \mathbb{Z}_4$:
  \[
    \begin{gathered}
    [0]_{12} \mapsto ([0]_3, [0]_4), \quad [1]_{12} \mapsto ([1]_3, [1]_4), \quad [3]_{12} \mapsto ([0]_3, [3]_4),\\
    [6]_{12} \mapsto ([0]_3, [2]_4), \quad [7]_{12} \mapsto ([1]_3, [3]_4).
    \end{gathered}
  \]
  $\psi_1$ è biunivoca.
  \item $\psi_2 : \mathbb{Z}_{12} \to \mathbb{Z}_6 \times \mathbb{Z}_2$:
  \[
    \begin{gathered}
    [0]_{12} \mapsto ([0]_6, [0]_2), \quad [1]_{12} \mapsto ([1]_6, [1]_2), \quad [3]_{12} \mapsto ([3]_6, [1]_2),\\
    [6]_{12} \mapsto ([0]_6, [0]_2), \quad [7]_{12} \mapsto ([1]_6, [1]_2).
    \end{gathered}
  \]
  $\psi_2$ \emph{non è iniettiva} (ad esempio $[1]_{12}$ e $[7]_{12}$ hanno la stessa immagine) e \emph{non è suriettiva}: non esiste $[a]_{12}$ tale che $\psi_2([a]_{12}) = ([0]_6, [1]_2)$.
\end{itemize}
\end{esempio}

\begin{osservazione}
$[a]_{nk}$ è invertibile in $\mathbb{Z}_{nk}$ se e solo se $[a]_n$ è invertibile in $\mathbb{Z}_n$ e $[a]_k$ è invertibile in $\mathbb{Z}_k$.

\emph{Dimostrazione.} $\Rightarrow$ Se $[a]_{nk}$ è invertibile in $\mathbb{Z}_{nk}$, allora $\operatorname{MCD}(a, nk) = 1$, quindi $a$ non ha fattori comuni con $n$ e non ne ha con $k$, cioè $\operatorname{MCD}(a, n) = 1$ e $\operatorname{MCD}(a, k) = 1$: $[a]_n$ è invertibile in $\mathbb{Z}_n$ e $[a]_k$ è invertibile in $\mathbb{Z}_k$. $\Leftarrow$ Viceversa, se $\operatorname{MCD}(a, n) = 1$ e $\operatorname{MCD}(a, k) = 1$, allora $a$ non ha fattori comuni né con $n$ né con $k$, quindi $\operatorname{MCD}(a, nk) = 1$.
\end{osservazione}

\begin{proposizione}
Se $\operatorname{MCD}(n, k) = 1$, allora $\psi : \mathbb{Z}_{nk} \to \mathbb{Z}_n \times \mathbb{Z}_k$, $\psi([a]_{nk}) = ([a]_n, [a]_k)$, è \emph{biunivoca}.
\end{proposizione}

\begin{proof}
Si ha $|\mathbb{Z}_{nk}| = nk$ e $|\mathbb{Z}_n\times\mathbb{Z}_k| = |\mathbb{Z}_n|\cdot|\mathbb{Z}_k| = nk$. Proviamo che $\psi$ è iniettiva. Siano $[a]_{nk}, [b]_{nk}$ tali che $\psi([a]_{nk}) = \psi([b]_{nk})$, cioè $([a]_n, [a]_k) = ([b]_n, [b]_k)$, cioè
\[
  \begin{cases} [a]_n = [b]_n \\ [a]_k = [b]_k \end{cases}
  \quad\text{cioè}\quad
  \begin{cases} a \equiv b \pmod n \\ a \equiv b \pmod k \end{cases}
  \ \Rightarrow\ a \equiv b \pmod{\operatorname{mcm}(n, k)} .
\]
Ma $\operatorname{MCD}(n, k) = 1$, quindi $\operatorname{mcm}(n, k) = nk$, cioè $a \equiv b \pmod{nk}$, cioè $[a]_{nk} = [b]_{nk}$. Quindi $\psi$ è iniettiva; dominio e codominio hanno la stessa cardinalità finita, quindi $\psi$ è biunivoca.
\end{proof}

\begin{proposizione}[Moltiplicatività di $\varphi$]
Sia $\operatorname{MCD}(n, k) = 1$. Allora
\[
  \varphi(nk) = \varphi(n)\cdot\varphi(k) .
\]
\end{proposizione}

\begin{proof}
$\varphi(nk) = |\mathbb{Z}_{nk}^*| = |\mathbb{Z}_n^*\times\mathbb{Z}_k^*| = |\mathbb{Z}_n^*|\cdot|\mathbb{Z}_k^*| = \varphi(n)\cdot\varphi(k)$, dove si è usato che, per l'osservazione precedente, la restrizione $\psi : \mathbb{Z}_{nk}^* \to \mathbb{Z}_n^*\times\mathbb{Z}_k^*$ è biunivoca.
\end{proof}

\section[Teorema di Eulero e potenze modulo $n$ (Lezione 14 -- 3 novembre 2026)]{Teorema di Eulero e potenze modulo $n$}
\noindent\emph{Lezione 14 -- 3 novembre 2026}
\subsection{Calcolo di $\varphi(n)$}

\begin{teorema}[Corollario]
Sia $n = p_1^{h_1}\cdot p_2^{h_2}\cdots p_r^{h_r}$ la fattorizzazione in primi. Allora
\[
  \varphi(n) = \varphi(p_1^{h_1})\cdot\varphi(p_2^{h_2})\cdots\varphi(p_r^{h_r}) = \left(p_1^{h_1} - p_1^{h_1 - 1}\right)\left(p_2^{h_2} - p_2^{h_2 - 1}\right)\cdots\left(p_r^{h_r} - p_r^{h_r - 1}\right).
\]
\end{teorema}

\begin{esempio}
\begin{align*}
  \varphi(1000) &= \varphi(10^3) = \varphi(5^3\cdot 2^3) = \varphi(5^3)\,\varphi(2^3)\\
  &= (5^3 - 5^2)(2^3 - 2^2) = (125 - 25)(8 - 4) = 100\cdot 4 = 400,\\
  \varphi(360) &= \varphi(2^3\cdot 3^2\cdot 5) = \varphi(2^3)\,\varphi(3^2)\,\varphi(5)\\
  &= (2^3 - 2^2)(3^2 - 3)(5 - 1) = 4\cdot 6\cdot 4 = 96 .
\end{align*}
\end{esempio}

\subsection{Teorema di Eulero}

\begin{teorema}[di Eulero]
Sia $n > 1$ e $a \in \mathbb{Z}$. Se $\operatorname{MCD}(a, n) = 1$ allora
\[
  a^{\varphi(n)} \equiv 1 \pmod n .
\]
\end{teorema}

(Per $n = p$ primo, $\varphi(p) = p - 1$ e si ritrova il corollario del piccolo teorema di Fermat.)

\subsection{Come calcolare le potenze modulo $n$}

Si vuole calcolare $a^b \bmod n$.
\begin{enumerate}
  \item Se $a > n$, si riduce la base modulo $n$: $a = n\cdot q + a'$, e $a^b \equiv (a')^b \pmod n$.
  \item Se $\operatorname{MCD}(a', n) = 1$, si applica il teorema di Eulero: si calcola $\varphi(n)$ e si divide $b$ per $\varphi(n)$: $b = q\cdot\varphi(n) + r$. Allora
  \[
    a^b \equiv (a')^b \equiv (a')^{q\varphi(n) + r} = \left((a')^{\varphi(n)}\right)^q\cdot(a')^r \equiv (a')^r \pmod n,
  \]
  perché $(a')^{\varphi(n)} \equiv 1 \pmod n$.
\end{enumerate}

\begin{esempio}[$187^{5\,867\,432} \bmod 12$]
Si riduce la base modulo $12$: $187 = 12\cdot 15 + 7$, quindi $187^{5\,867\,432} \equiv 7^{5\,867\,432} \pmod{12}$. Si verifica che $\operatorname{MCD}(7, 12) = 1$, quindi si può applicare il teorema di Eulero:
\[
  \varphi(12) = \varphi(3\cdot 2^2) = \varphi(3)\,\varphi(2^2) = (3 - 1)(2^2 - 2) = 4 .
\]
Poiché $5\,867\,432$ è divisibile per $4$ (il resto è $0$), 
\[
  187^{5\,867\,432} \equiv 7^{5\,867\,432} \equiv 7^0 \equiv 1 \pmod{12} .
\]
\end{esempio}

\begin{esempio}[$4^{48\,642\,369} - 33^{1\,759\,468} \bmod 10$]
Qui $\operatorname{MCD}(4, 10) = 2$: i numeri non sono coprimi, quindi \emph{non} si può applicare il teorema di Eulero. Si cercano le potenze di $4$ modulo $10$:
\[
  \begin{gathered}
  4^0 = 1,\quad 4^1 = 4,\quad 4^2 = 16 \equiv 6,\\
  4^3 = 4^2\cdot 4 \equiv 6\cdot 4 = 24 \equiv 4,\quad 4^4 = 4^3\cdot 4 \equiv 4\cdot 4 = 16 \equiv 6 \pmod{10}.
  \end{gathered}
\]
Le potenze pari di $4$ sono $\equiv 6$, quelle dispari $\equiv 4 \pmod{10}$. L'esponente $48\,642\,369$ è dispari, quindi $4^{48\,642\,369} \equiv 4 \pmod{10}$.

Per l'altro termine: $33 \equiv 3 \pmod{10}$ e $\operatorname{MCD}(3, 10) = 1$, quindi si può applicare Eulero con $\varphi(10) = \varphi(2)\,\varphi(5) = (2 - 1)(5 - 1) = 4$. Poiché $1\,759\,468$ è divisibile per $4$ (resto $0$),
\[
  33^{1\,759\,468} \equiv 3^{1\,759\,468} \equiv 3^0 \equiv 1 \pmod{10}.
\]
In conclusione $4^{48\,642\,369} - 33^{1\,759\,468} \equiv 4 - 1 \equiv 3 \pmod{10}$.
\end{esempio}

\section[Criteri di divisibilità e potenze modulo $n$ (Lezione 15 -- 9 novembre 2026)]{Criteri di divisibilità e potenze modulo $n$}
\noindent\emph{Lezione 15 -- 9 novembre 2026}

\subsection{Criteri di divisibilità}

Sia $n \in \mathbb{N}$ scritto in base $10$:
\[
  n = a_k\cdot 10^k + a_{k-1}\cdot 10^{k-1} + \dots + a_1\cdot 10 + a_0 = (a_k a_{k-1}\dots a_1 a_0)_{10}, \qquad a_i \in \{0, 1, \dots, 9\}.
\]

\begin{proposizione}[Criteri di divisibilità]
\begin{itemize}
  \item $n$ è divisibile per $2$ se e solo se $a_0$ è divisibile per $2$.
  \item $n$ è divisibile per $10$ se e solo se $10 \mid a_0$ (cioè $a_0 = 0$).
  \item $n$ è divisibile per $5$ se e solo se $5 \mid a_0$.
  \item $n$ è divisibile per $2^j$ se e solo se $(a_{j-1}\dots a_1 a_0)_{10}$ è divisibile per $2^j$.
  \item $n$ è divisibile per $5^j$ se e solo se $5^j \mid (a_{j-1}\dots a_1 a_0)_{10}$.
  \item $n$ è divisibile per $3$ se e solo se $a_k + a_{k-1} + \dots + a_1 + a_0$ è divisibile per $3$.
  \item $n$ è divisibile per $9$ se e solo se $a_k + \dots + a_1 + a_0$ è divisibile per $9$.
  \item $n$ è divisibile per $11$ se e solo se $11 \mid a_0 - a_1 + a_2 - a_3 + \dots + (-1)^k a_k$.
\end{itemize}
\end{proposizione}

\begin{proof}
\emph{Criterio per $2$ (e per $5$, $10$).} Poiché $10 \equiv 0 \pmod 2$ (e analogamente $10 \equiv 0 \pmod 5$ e $10 \equiv 0 \pmod{10}$),
\[
  n = a_k\cdot 10^k + \dots + a_1\cdot 10 + a_0 \equiv a_0 \pmod 2 ,
\]
quindi $n \equiv a_0 \pmod 2$.

\emph{Criterio per $3$ e per $9$.} Poiché $10 = 3\cdot 3 + 1$, cioè $10 \equiv 1 \pmod 3$ (e $10 \equiv 1 \pmod 9$),
\[
  n = a_k\cdot 10^k + \dots + 10\cdot a_1 + a_0 \equiv a_k + a_{k-1} + \dots + a_1 + a_0 \pmod 3 \ (\text{o } \bmod 9).
\]

\emph{Criterio per $11$.} Poiché $10 \equiv -1 \pmod{11}$ (infatti $-1 = -1\cdot 11 + 10$),
\[
  n = a_0 + a_1\cdot 10 + a_2\cdot 10^2 + \dots + a_k\cdot 10^k \equiv a_0 - a_1 + a_2 - a_3 + \dots + (-1)^k a_k \pmod{11}. \qedhere
\]
\end{proof}

\begin{osservazione}
Per $2^j$ e $5^j$ il criterio segue dal fatto che $10^j = 2^j 5^j \equiv 0$ modulo $2^j$ e modulo $5^j$: contano solo le ultime $j$ cifre.
\end{osservazione}

\begin{esempio}
Sia $n = 16\,079$. Si calcolano i resti per diversi divisori.
\begin{center}
\begin{tabular}{rl}
\toprule
Divisore & Resto \\
\midrule
$2$ & $1$ \\
$5$ & $9 \equiv 4 \pmod 5$ \\
$10$ & $9$ \\
$4 = 2^2$ & $79 \equiv 3 \pmod 4$ \quad ($79 = 76 + 3$) \\
$25 = 5^2$ & $79 \equiv 4 \pmod{25}$ \\
$1000 = 10^3$ & $079 = 79 \pmod{1000}$ \\
$8 = 2^3$ & $079 \equiv 7 \pmod 8$ \quad ($79 = 9\cdot 8 + 7$) \\
$3$ & $1 + 6 + 0 + 7 + 9 = 23 \equiv 2 + 3 = 5 \equiv 2 \pmod 3$ \\
$9$ & $23 \equiv 2 + 3 = 5 \pmod 9$ \\
$11$ & $9 - 7 + 0 - 6 + 1 = -3 \equiv 8 \pmod{11}$ \\
\bottomrule
\end{tabular}
\end{center}
\end{esempio}

\begin{esempio}[$4589 \equiv x \pmod{44}$]
Poiché $44 = 4\cdot 11$ con $\operatorname{MCD}(4, 11) = 1$, la funzione
\[
  \psi : \mathbb{Z}_{44} \to \mathbb{Z}_{11}\times\mathbb{Z}_4, \qquad [x]_{44} \mapsto \big([x]_{11}, [x]_4\big)
\]
è biunivoca. Si ha
\[
  \begin{cases} 4589 \equiv 89 \equiv 1 \pmod 4 \\ 4589 \equiv 9 - 8 + 5 - 4 = 2 \pmod{11} \end{cases}
\]
cioè $\psi([x]_{44}) = ([2]_{11}, [1]_4)$. I numeri congrui a $2$ modulo $11$ sono $2, 13, 24, 35, \dots$; tra quelli minori di $44$:
\[
  \begin{gathered}
  13 = 11 + 2,\ 13 \equiv 1 \pmod 4; \qquad 24 = 2\cdot 11 + 2,\ 24 \equiv 0 \pmod 4;\\
  35 = 3\cdot 11 + 2,\ 35 \equiv 3 \pmod 4 .
  \end{gathered}
\]
Quindi l'unica soluzione è $[13]_{44}$: $4589 \equiv 13 \pmod{44}$.
\end{esempio}

\subsection{Potenze modulo $n$: esponente in base $2$}

\emph{Problema.} Calcolare $a^b \bmod n$.

Si scrive l'esponente in base $2$: $b = \varepsilon_k 2^k + \varepsilon_{k-1} 2^{k-1} + \dots + \varepsilon_1 2 + \varepsilon_0$, con $\varepsilon_i \in \{0, 1\}$. Allora
\[
  a^b = a^{\varepsilon_k 2^k + \dots + \varepsilon_1 2 + \varepsilon_0} = a^{\varepsilon_k 2^k}\cdot a^{\varepsilon_{k-1}2^{k-1}}\cdots a^{\varepsilon_1 2}\cdot a^{\varepsilon_0}.
\]
Procedimento:
\begin{enumerate}
  \item si riduce l'esponente in base $2$;
  \item si calcolano $a^{2^i} \bmod n$ (elevando via via al quadrato);
  \item si calcola il prodotto (modulo $n$) dei fattori $a^{2^i}$ con $\varepsilon_i = 1$.
\end{enumerate}

\begin{esempio}[$3^{13} \bmod 100$]
Scriviamo l'esponente in base $2$ con divisioni successive:
\[
  13 = 6\cdot 2 + 1,\quad 6 = 3\cdot 2 + 0,\quad 3 = 1\cdot 2 + 1,\quad 1 = 0\cdot 2 + 1,
\]
quindi $13 = (1101)_2 = 1\cdot 2^3 + 1\cdot 2^2 + 0\cdot 2 + 1 = 2^3 + 2^2 + 1$ e
\[
  3^{13} = 3^{2^3}\cdot 3^{2^2}\cdot 3 .
\]
Calcoliamo le potenze:
\[
  3^{2} = 9, \qquad 3^{2^2} = 9\cdot 9 = 81, \qquad 3^{2^3} = 81\cdot 81 = 6561 \equiv 61 \pmod{100}.
\]
Allora
\[
  3^{13} \equiv 61\cdot 81\cdot 3 \equiv 61\cdot 243 \equiv 61\cdot 43 = 2623 \equiv 23 \pmod{100}.
\]
\end{esempio}

\begin{esercizio}
Calcolare $203^{345\,618} \bmod 100$.
\end{esercizio}

\section[Equazioni diofantee e congruenze lineari (Lezione 15 -- 9 novembre 2026)]{Equazioni diofantee e congruenze lineari}
\noindent\emph{Lezione 15 -- 9 novembre 2026}

\subsection{Equazioni diofantee}

\begin{definizione}[Equazione diofantea]
Siano $a, b, c \in \mathbb{Z}$. Un'equazione del tipo
\[
  ax + by = c
\]
della quale cerchiamo soluzioni intere $x, y \in \mathbb{Z}$ si chiama \emph{equazione diofantea}.
\end{definizione}

\begin{esempio}
\begin{enumerate}
  \item $5x + 3y = 1$. Poiché $\operatorname{MCD}(5, 3) = 1 \mid 1$ ci sono soluzioni: ad esempio $(x, y) = (-1, 2)$ e $(x, y) = (2, -3)$. Tutte le soluzioni sono
  \[
    \begin{cases} x = -1 - 3k \\ y = 2 + 5k \end{cases} \qquad k \in \mathbb{Z}.
  \]
  \item $6x + 3y = 5$: impossibile in $\mathbb{Z}$, perché $3(2x + y) = 5$ e $\operatorname{MCD}(6, 3) = 3 \nmid 5$.
\end{enumerate}
\end{esempio}

\begin{proposizione}[Esistenza di soluzioni]
Siano $a, b, c \in \mathbb{Z}$. L'equazione diofantea $ax + by = c$ ha soluzioni se e solo se $\operatorname{MCD}(a, b) \mid c$.
\end{proposizione}

\begin{proof}
$(\Rightarrow)$ Siano $\bar{x}, \bar{y} \in \mathbb{Z}$ soluzioni: $a\bar{x} + b\bar{y} = c$. Allora $\operatorname{MCD}(a, b)$ divide $a\bar{x} + b\bar{y} = c$.

$(\Leftarrow)$ Sia $\operatorname{MCD}(a, b) = d \mid c$, cioè $c = hd$. Per il teorema di Bézout esistono $\alpha, \beta \in \mathbb{Z}$ tali che $\alpha a + \beta b = d$. Moltiplicando per $h$:
\[
  h\alpha\, a + h\beta\, b = dh = c,
\]
quindi $x = h\alpha$, $y = h\beta$ è una soluzione. \qedhere
\end{proof}

\begin{proposizione}[Tutte le soluzioni]
Siano $a, b, c \in \mathbb{Z}$ e sia $(\bar{x}, \bar{y})$ una soluzione di $ax + by = c$. Allora tutte le soluzioni intere sono del tipo
\[
  \begin{cases} x' = \bar{x} - \dfrac{b}{d}\,k \\[2mm] y' = \bar{y} + \dfrac{a}{d}\,k \end{cases} \qquad k \in \mathbb{Z}, \quad d = \operatorname{MCD}(a, b).
\]
\end{proposizione}

\begin{esempio}[$6x - 8y = 10$]
Si vogliono trovare tutte le soluzioni intere. L'equazione è risolubile: $\operatorname{MCD}(6, -8) = 2 \mid 10$.

Si applica Bézout fra $6$ e $8$ con l'algoritmo di Euclide esteso (colonne $q$, $r$, $s$, $t$):
\begin{center}
\begin{tabular}{c|r|r|r}
$q$ & $r$ & $s$ & $t$ \\
\midrule
 & $8$ & $1$ & $0$ \\
 & $6$ & $0$ & $1$ \\
$1$ & $2$ & $1$ & $-1$ \\
$3$ & $0$ & $-3$ & $4$ \\
\end{tabular}
\end{center}
Quindi $\operatorname{MCD}(6, 8) = 2 = 1\cdot 8 - 1\cdot 6$. Moltiplicando per $5$:
\[
  10 = 5\cdot 8 - 5\cdot 6 = 6\cdot(-5) - 8\cdot(-5),
\]
quindi una soluzione particolare è $\bar{x} = -5$, $\bar{y} = -5$. La soluzione generica è
\[
  \begin{cases} x' = -5 - \dfrac{-8}{2}\,k = -5 + 4k \\[2mm] y' = -5 + \dfrac{6}{2}\,k = -5 + 3k \end{cases} \qquad k \in \mathbb{Z}.
\]
\end{esempio}

\subsection{Equazioni modulari}

\begin{definizione}[Equazione congruenziale lineare]
Un'\emph{equazione congruenziale lineare} (o \emph{equazione modulare lineare}) nell'incognita $x$ è un'equazione della forma
\[
  ax \equiv b \pmod n, \qquad a, b \in \mathbb{Z},\ n \in \mathbb{N},\ n > 0.
\]
Si cercano soluzioni intere $x \in \mathbb{Z}$, o soluzioni $[x]_n \in \mathbb{Z}_n$.
\end{definizione}

\begin{esempio}
\begin{enumerate}
  \item $6x \equiv 21 \pmod 4$. Riducendo i coefficienti modulo $4$: $2x \equiv 1 \pmod 4$. Poiché $\operatorname{MCD}(6, 4) = 2 \nmid 21$, l'equazione è \emph{impossibile}.
  \item $11x \equiv -3 \pmod 6$, cioè $5x \equiv 3 \pmod 6$. Si cerca l'inverso di $5$ in $\mathbb{Z}_6$: $5\cdot 1 \equiv 5$, $5\cdot 2 = 10 \equiv 4$, \dots, $5\cdot 5 = 25 \equiv 1 \pmod 6$, quindi $[5]_6^{-1} = [5]_6$. Moltiplicando per $5$:
  \[
    5\cdot 5\,x \equiv 3\cdot 5 \pmod 6 \quad\Rightarrow\quad x \equiv 15 \equiv 3 \pmod 6,
  \]
  quindi $S = [x]_6 = [3]_6 = \{6k + 3 \mid k \in \mathbb{Z}\}$.
  \item $4x \equiv 2 \pmod 6$. Qui $\operatorname{MCD}(4, 6) = 2 \mid 2$. Si usa la tabella di $4\cdot x$ in $\mathbb{Z}_6$:
  \[
    4\cdot 0 = 0,\ 4\cdot 1 = 4,\ 4\cdot 2 = 8 \equiv 2,\ 4\cdot 3 = 12 \equiv 0,\ 4\cdot 4 = 16 \equiv 4,\ 4\cdot 5 = 20 \equiv 2.
  \]
  Le soluzioni sono $x = 2$ e $x = 5$:
  \[
    S = [2]_6 \cup [5]_6 = \{6k + 2 \mid k \in \mathbb{Z}\} \cup \{6k + 5 \mid k \in \mathbb{Z}\}.
  \]
\end{enumerate}
\end{esempio}

\begin{osservazione}
\begin{enumerate}
  \item Se $ax \equiv b \pmod n$ e $\operatorname{MCD}(a, n) = 1$, allora esiste l'inverso di $a$ in $\mathbb{Z}_n$, $\left[\tfrac{1}{a}\right]_n = [a^{-1}]_n$, e moltiplicando si ottiene $a^{-1}a\,x \equiv a^{-1}b$, cioè $x \equiv a^{-1}b \pmod n$.
  \item Se $t \mid n$ (e $t$ divide anche $a$ e $b$), l'equazione $ax \equiv b \pmod n$ equivale a $\dfrac{a}{t}\,x \equiv \dfrac{b}{t} \pmod{\dfrac{n}{t}}$ (si veda la proposizione nella lezione successiva).
  \item Se $ac \equiv bc \pmod n$ e $\operatorname{MCD}(c, n) = 1$, allora $a \equiv b \pmod n$ (legge di cancellazione).
\end{enumerate}
\end{osservazione}

\begin{esempio}[$12x \equiv 8 \pmod{35}$]
Poiché $\operatorname{MCD}(12, 35) = 1 \mid 8$, la soluzione esiste. Usiamo Euclide esteso per trovare l'inverso di $12$ in $\mathbb{Z}_{35}$:
\begin{center}
\begin{tabular}{c|r|r|r}
$q$ & $r$ & $s$ & $t$ \\
\midrule
 & $35$ & $1$ & $0$ \\
 & $12$ & $0$ & $1$ \\
$2$ & $11$ & $1$ & $-2$ \\
$1$ & $1$ & $-1$ & $3$ \\
$11$ & $0$ & $12$ & $-35$ \\
\end{tabular}
\end{center}
Quindi $1 = -1\cdot 35 + 3\cdot 12$, cioè $3\cdot 12 \equiv 1 \pmod{35}$ e $\left[\tfrac{1}{12}\right]_{35} = [3]_{35}$. Moltiplicando l'equazione per $3$: $3\cdot 12\,x \equiv 3\cdot 8 \pmod{35}$, cioè $x \equiv 24 \pmod{35}$, e
\[
  S = [24]_{35} = \{24 + 35k \mid k \in \mathbb{Z}\}.
\]
\end{esempio}

\begin{proposizione}[Condizione di compatibilità]
L'equazione modulare $ax \equiv b \pmod n$ ha soluzione se e solo se $\operatorname{MCD}(a, n) \mid b$.
\end{proposizione}

\begin{proof}
$ax \equiv b \pmod n \iff ax - b = kn \iff ax - kn = b$. L'ultima è un'equazione diofantea (nelle incognite $x$, $k$), che ha soluzione se e solo se $\operatorname{MCD}(a, n) \mid b$.
\end{proof}

\section[Equazioni modulari e sistemi (Lezione 16 -- 10 novembre 2026)]{Equazioni modulari e sistemi}
\noindent\emph{Lezione 16 -- 10 novembre 2026}

\subsection{Equazioni modulari: riduzione e insieme delle soluzioni}

Si consideri l'equazione modulare
\[
  ax \equiv b \pmod n, \qquad a, b \in \mathbb{Z},\ n > 0,\ n \in \mathbb{N},
\]
con incognita $x$. Essa ha soluzione se e solo se $\operatorname{MCD}(a, n) \mid b$.

\begin{proposizione}
Siano $a, b \in \mathbb{Z}$, $n > 0$, e $d = \operatorname{MCD}(a, n) \mid b$. Allora l'equazione
\[
  ax \equiv b \pmod n \tag{1}
\]
è equivalente a
\[
  \frac{a}{d}\,x \equiv \frac{b}{d} \pmod{\frac{n}{d}} . \tag{2}
\]
In particolare, se $\bar{x}$ è soluzione di (1) (o di (2)), allora l'insieme delle soluzioni è
\[
  S = \left[\bar{x}\right]_{n/d} = \left[\bar{x}\right]_n \cup \left[\bar{x} + \tfrac{n}{d}\right]_n \cup \left[\bar{x} + 2\tfrac{n}{d}\right]_n \cup \dots \cup \left[\bar{x} + (d-1)\tfrac{n}{d}\right]_n = \left\{\bar{x} + k\tfrac{n}{d} \;\middle|\; k \in \mathbb{Z}\right\}.
\]
\end{proposizione}

In pratica, se $d \neq 1$ si divide l'equazione per $d$ ottenendo un'equazione modulo $n/d$ in cui il coefficiente di $x$ è invertibile (perché $\operatorname{MCD}\!\left(\tfrac{a}{d}, \tfrac{n}{d}\right) = 1$); la si risolve moltiplicando per l'inverso, e poi si descrivono le soluzioni modulo $n$.

\begin{esempio}[$36x \equiv 24 \pmod{105}$]
Si vogliono trovare tutte le soluzioni intere comprese fra $-100$ e $100$.

\emph{Condizione di compatibilità:} $\operatorname{MCD}(36, 105) = 3 \mid 24$. Vero. $[36]_{105}$ non è invertibile.

\emph{Riduzione:} si divide l'equazione per $d = 3$:
\[
  \frac{36}{3}\,x \equiv \frac{24}{3} \pmod{\frac{105}{3}} \quad\Rightarrow\quad 12x \equiv 8 \pmod{35}.
\]
Si cerca l'inverso di $[12]_{35}$ in $\mathbb{Z}_{35}$ con Euclide esteso:
\begin{center}
\begin{tabular}{c|r|r|r}
$q$ & $r$ & $s$ & $t$ \\
\midrule
 & $35$ & $1$ & $0$ \\
 & $12$ & $0$ & $1$ \\
$2$ & $11$ & $1$ & $-2$ \\
$1$ & $1$ & $-1$ & $1 + 2 = 3$ \\
$11$ & $0$ & $1 + 11 = 12$ & $-2 - 33 = -35$ \\
\end{tabular}
\end{center}
Da cui $1 = (-1)\cdot 35 + 3\cdot 12$: l'inverso di $[12]_{35}$ in $\mathbb{Z}_{35}$ è $[3]_{35}$. Moltiplicando per $3$:
\[
  3\cdot 12\,x \equiv 8\cdot 3 \pmod{35} \quad\Rightarrow\quad x \equiv 24 \pmod{35},
\]
quindi $S = [24]_{35} = \{24 + 35k \mid k \in \mathbb{Z}\}$ (soluzioni in $\mathbb{Z}_{35}$).

\emph{Soluzioni fra $-100$ e $100$:}
\[
  24,\ 24 + 35 = 59,\ 24 + 2\cdot 35 = 94,\ (94 + 35 = 129 > 100);
\]
\[
  24 - 35 = -11,\ -11 - 35 = -46,\ -46 - 35 = -81,\ (-81 - 35 = -116 < -100).
\]
Le soluzioni tra $-100$ e $100$ sono $\{-81, -46, -11, 24, 59, 94\}$.

In $\mathbb{Z}_{105}$ le soluzioni sono tre classi: $[24]_{105} \cup [24 + 35]_{105} \cup [24 + 70]_{105} = [24]_{105} \cup [59]_{105} \cup [94]_{105}$.
\end{esempio}

\subsection{Sistemi modulari}

Un sistema modulare è un insieme di equazioni modulari nella stessa incognita, da risolvere simultaneamente. Prima di tutto si verificano le \emph{condizioni di compatibilità} di ciascuna equazione.

\begin{esempio}
\begin{enumerate}
  \item
  \[
    \begin{cases} 3x \equiv 5 \pmod{12} \\ x \equiv 2 \pmod 4 \end{cases}
  \]
  Condizioni di compatibilità: $\operatorname{MCD}(3, 12) = 3 \nmid 5$ (la prima equazione è impossibile), $\operatorname{MCD}(1, 4) = 1 \mid 2$. Il sistema è impossibile.
  \item
  \[
    \begin{cases} x \equiv 3 \pmod 6 \\ x \equiv 0 \pmod 2 \end{cases}
  \]
  Condizioni di compatibilità: $\operatorname{MCD}(1, 6) = 1 \mid 3$, $\operatorname{MCD}(1, 2) = 1 \mid 0$. Dalla seconda equazione $x = 2k$; sostituendo nella prima: $2k \equiv 3 \pmod 6$, che è \emph{impossibile}: $\operatorname{MCD}(2, 6) = 2 \nmid 3$. Il sistema è impossibile.
  \item
  \[
    \begin{cases} 5x \equiv 3 \pmod 6 \\ x \equiv 1 \pmod 5 \end{cases}
  \]
  Condizioni di compatibilità: $\operatorname{MCD}(5, 6) = 1 \mid 3$, $\operatorname{MCD}(1, 5) = 1 \mid 1$: ok.

  Prima equazione: $5x \equiv 3 \pmod 6$, $\left[\tfrac{1}{5}\right]_6 = [5]_6$; moltiplicando per $5$: $x \equiv 15 \equiv 3 \pmod 6$, cioè $x = 6k + 3$.

  Si sostituisce nella seconda equazione: $6k + 3 \equiv 1 \pmod 5$, cioè $k \equiv 1 - 3 = -2 \equiv 3 \pmod 5$, quindi $k = 5h + 3$.

  Allora $x = 6(5h + 3) + 3 = 30h + 18 + 3 = 30h + 21$, cioè
  \[
    S = [21]_{30} = \{21 + 30h \mid h \in \mathbb{Z}\}.
  \]
\end{enumerate}
\end{esempio}

\section[Teorema cinese dei resti (Lezione 16 -- 10 novembre 2026)]{Teorema cinese dei resti}
\noindent\emph{Lezione 16 -- 10 novembre 2026}

\subsection{Teorema cinese dei resti}

\begin{teorema}[Teorema cinese dei resti]
Siano $n_1, \dots, n_s$ interi positivi \emph{coprimi fra loro}, cioè $\operatorname{MCD}(n_i, n_j) = 1$ per $i \neq j$. Allora il sistema modulare
\[
  \begin{cases}
    x \equiv r_1 \pmod{n_1} \\
    x \equiv r_2 \pmod{n_2} \\
    \ \ \vdots \\
    x \equiv r_s \pmod{n_s}
  \end{cases}
\]
ha un'\emph{unica} soluzione $\bar{x}$ modulo $N = n_1\cdot n_2\cdots n_s$, cioè l'insieme delle soluzioni è
\[
  S = [\bar{x}]_N = \{\bar{x} + kN \mid k \in \mathbb{Z}\}.
\]
\end{teorema}

\begin{proof}
\emph{Esistenza.} Sia $N = n_1 n_2 \cdots n_s$ e $N_i = \dfrac{N}{n_i}$. L'equazione
\[
  N_i\, x \equiv r_i \pmod{n_i}
\]
ha sempre soluzione $\bar{x}_i$ (perché $\operatorname{MCD}(N_i, n_i) = 1 \mid r_i$). Dopo aver risolto ogni equazione in questo modo, si definisce
\[
  \bar{x} = N_1\bar{x}_1 + N_2\bar{x}_2 + \dots + N_s\bar{x}_s .
\]
Allora $\bar{x}$ è soluzione del sistema di partenza: per ogni $i$, $N_j \equiv 0 \pmod{n_i}$ per $j \neq i$, quindi
\[
  \bar{x} = N_1\bar{x}_1 + \dots + N_s\bar{x}_s \equiv N_i\bar{x}_i \equiv r_i \pmod{n_i}.
\]

\emph{Unicità.} Siano $\bar{x}$ e $\bar{y}$ due soluzioni ($\bar{x} \equiv r_i$, $\bar{y} \equiv r_i \pmod{n_i}$). Allora
\[
  \bar{x} - \bar{y} \equiv r_i - r_i \equiv 0 \pmod{n_i},
\]
cioè $\bar{x} - \bar{y}$ è multiplo di $n_i$ per tutti gli $i$. Poiché gli $n_i$ sono coprimi fra loro, $\bar{x} - \bar{y}$ è multiplo di $n_1 n_2 \cdots n_s = N$, cioè $\bar{x} \equiv \bar{y} \pmod N$, cioè $[\bar{x}]_N = [\bar{y}]_N$.
\end{proof}

\begin{esempio}
\[
  \begin{cases} x \equiv 2 \pmod 3 \\ x \equiv 3 \pmod 4 \\ x \equiv 4 \pmod 5 \end{cases}
\]
Condizioni di compatibilità: $\operatorname{MCD}(1, 3) = 1 \mid 2$, $\operatorname{MCD}(1, 4) = 1 \mid 3$, $\operatorname{MCD}(1, 5) = 1 \mid 4$. I moduli $3, 4, 5$ sono coprimi fra loro: si applica il teorema cinese dei resti con $N = 3\cdot 4\cdot 5 = 60$. Si risolvono tre equazioni:
\begin{enumerate}
  \item $N_1 = \dfrac{N}{n_1} = \dfrac{3\cdot 4\cdot 5}{3} = 4\cdot 5 = 20$: $20x_1 \equiv 2 \pmod 3$, cioè $2x_1 \equiv 2 \pmod 3$, $x_1 \equiv 1 \pmod 3$; si prende $x_1 = 1$.
  \item $N_2 = \dfrac{N}{n_2} = 3\cdot 5 = 15$: $15x_2 \equiv 3 \pmod 4$, cioè $3x_2 \equiv 3 \pmod 4$, $x_2 \equiv 1 \pmod 4$; $x_2 = 1$.
  \item $N_3 = \dfrac{N}{n_3} = 3\cdot 4 = 12$: $12x_3 \equiv 4 \pmod 5$, cioè $3x_3 \equiv 1 \pmod 5$ (dividendo per $4$, invertibile modulo $5$); $\left[\tfrac{1}{3}\right]_5 = [2]$, quindi $2\cdot 3\,x_3 \equiv 2 \pmod 5$, $x_3 \equiv 2 \pmod 5$; $x_3 = 2$.
\end{enumerate}
La soluzione generale è
\[
  \bar{x} = N_1x_1 + N_2x_2 + N_3x_3 = 20\cdot 1 + 15\cdot 1 + 12\cdot 2 = 59,
\]
modulo $N = 3\cdot 4\cdot 5 = 60$:
\[
  S = [59]_{60} = \{59 + 60k \mid k \in \mathbb{Z}\}.
\]
\end{esempio}

\subsection{Sistemi modulari generali}

Si consideri un sistema
\[
  \begin{cases} a_1x \equiv b_1 \pmod{n_1} \\ \ \ \vdots \\ a_sx \equiv b_s \pmod{n_s} \end{cases}
\]
Procedimento:
\begin{enumerate}
  \item \emph{Condizioni di compatibilità:} $\operatorname{MCD}(a_i, n_i) \mid b_i$ per ogni $i$. Se almeno una non è soddisfatta, il sistema è impossibile.
  \item Se $d_i = \operatorname{MCD}(a_i, n_i) \neq 1$, si divide l'equazione per $d_i$: $\dfrac{a_i}{d_i}\,x \equiv \dfrac{b_i}{d_i} \pmod{\dfrac{n_i}{d_i}}$.
  \item Si ricava $x$ in ciascuna equazione moltiplicando per $\left[\tfrac{1}{a_i}\right]_{n_i}$: $a_ix \equiv b_i \Rightarrow x \equiv b_i\left[\tfrac{1}{a_i}\right] \pmod{n_i}$, ottenendo un sistema della forma $x \equiv c_i \pmod{n_i}$.
  \item Se gli $n_i$ sono coprimi fra loro, si applica il teorema cinese dei resti; altrimenti si trasforma il sistema in uno equivalente con moduli coprimi (vedi sotto).
\end{enumerate}

\begin{esempio}
\[
  \begin{cases} 2x \equiv 1 \pmod 5 \\ 3x \equiv 2 \pmod 7 \\ 4x \equiv 3 \pmod{11} \end{cases}
\]
\begin{enumerate}
  \item Condizioni di compatibilità: $\operatorname{MCD}(2, 5) = 1$, $\operatorname{MCD}(3, 7) = 1$, $\operatorname{MCD}(4, 11) = 1$: ok.
  \item I MCD sono tutti $1$: ok.
  \item Si ricava $x$ nelle equazioni:
  \begin{itemize}
    \item $2x \equiv 1 \pmod 5$: $\left[\tfrac{1}{2}\right]_5 = [3]_5$, quindi $2\cdot 3\,x \equiv 3\cdot 1$, cioè $x \equiv 3 \pmod 5$;
    \item $3x \equiv 2 \pmod 7$: $\left[\tfrac{1}{3}\right]_7 = [5]_7$, quindi $3\cdot 5\,x \equiv 2\cdot 5$, cioè $x \equiv 10 \equiv 3 \pmod 7$;
    \item $4x \equiv 3 \pmod{11}$: $\left[\tfrac{1}{4}\right]_{11} = [3]_{11}$ (ad esempio con Euclide esteso), quindi $x \equiv 3\cdot 3 = 9 \pmod{11}$.
  \end{itemize}
  Il sistema diventa
  \[
    \begin{cases} x \equiv 3 \pmod 5 \\ x \equiv 3 \pmod 7 \\ x \equiv 9 \pmod{11} \end{cases}
  \]
  \item I moduli sono coprimi: si applica il teorema cinese dei resti con $N = 5\cdot 7\cdot 11 = 385$.
  \begin{itemize}
    \item $N_1 = 7\cdot 11 = 77$: $77x_1 \equiv 3 \pmod 5$, cioè $2x_1 \equiv 3 \pmod 5$, da cui $x_1 \equiv 4 \pmod 5$ (moltiplicando per $\left[\tfrac{1}{2}\right]_5 = [3]$: $x_1 \equiv 9 \equiv 4$).
    \item $N_2 = 5\cdot 11 = 55$: $55x_2 \equiv 3 \pmod 7$, cioè $6x_2 \equiv 3 \pmod 7$; dividendo per $3$: $2x_2 \equiv 1 \pmod 7$, $\left[\tfrac{1}{2}\right]_7 = [4]_7$, quindi $x_2 \equiv 4 \pmod 7$.
    \item $N_3 = 5\cdot 7 = 35$: $35x_3 \equiv 9 \pmod{11}$, cioè $2x_3 \equiv 9 \pmod{11}$; $\left[\tfrac{1}{2}\right]_{11} = [6]_{11}$, quindi $x_3 \equiv 9\cdot 6 = 54 \equiv 10 \pmod{11}$.
  \end{itemize}
  La soluzione generale è
  \[
    \bar{x} = N_1x_1 + N_2x_2 + N_3x_3 = 77\cdot 4 + 55\cdot 4 + 35\cdot 10 = 308 + 220 + 350,
  \]
  modulo $385$: $308 \equiv -77$ e $350 \equiv -35$, quindi $\bar{x} \equiv -77 + 220 - 35 = 108 \pmod{385}$. Dunque
  \[
    S = [108]_{385} = \{108 + 385k \mid k \in \mathbb{Z}\}.
  \]
\end{enumerate}
\end{esempio}

\subsection{Moduli non coprimi}

\begin{esempio}
\[
  \begin{cases} x \equiv 1 \pmod{10} \\ x \equiv 5 \pmod 6 \end{cases}
\]
Le condizioni di compatibilità sono soddisfatte. I moduli \emph{non} sono coprimi: $\operatorname{MCD}(6, 10) = 2 \neq 1$, quindi non si può applicare direttamente il teorema cinese dei resti. Si usa la funzione biunivoca $\psi$ per spezzare ciascuna equazione in potenze di primi.

Prima equazione: $\psi : \mathbb{Z}_{10} \to \mathbb{Z}_2\times\mathbb{Z}_5$,
\[
  x \equiv 1 \pmod{10} \iff \begin{cases} x \equiv 1 \pmod 2 \\ x \equiv 1 \pmod 5. \end{cases}
\]
Seconda equazione: $\psi : \mathbb{Z}_6 \to \mathbb{Z}_2\times\mathbb{Z}_3$,
\[
  x \equiv 5 \pmod 6 \iff \begin{cases} x \equiv 5 \equiv 2 \pmod 3 \\ x \equiv 5 \equiv 1 \pmod 2. \end{cases}
\]
Il sistema di partenza diventa
\[
  \begin{cases} x \equiv 1 \pmod 2 \\ x \equiv 1 \pmod 5 \\ x \equiv 2 \pmod 3 \\ x \equiv 1 \pmod 2 \end{cases}
  \iff
  \begin{cases} x \equiv 1 \pmod 2 \\ x \equiv 1 \pmod 5 \\ x \equiv 2 \pmod 3 \end{cases}
\]
(la quarta equazione è ridondante). Ora i moduli $2, 3, 5$ sono coprimi e si applica il teorema cinese dei resti.
\end{esempio}

\begin{osservazione}
Il sistema
\[
  \begin{cases} x \equiv b \pmod{p^i} \\ x \equiv c \pmod{p^j} \end{cases} \qquad p \text{ primo},\ i < j,
\]
ha soluzione se e solo se $b \equiv c \pmod{p^i}$ (potenza più piccola). In tal caso il sistema è equivalente a $x \equiv c \pmod{p^j}$ (l'equazione con il modulo più alto).
\end{osservazione}

\begin{esercizio}
Risolvere i sistemi
\[
  \begin{cases} x \equiv 3 \pmod 4 \\ x \equiv -1 \pmod{64} \end{cases}
  \qquad\text{e}\qquad
  \begin{cases} x \equiv 1 \pmod{10} \\ 2x \equiv 14 \pmod{24} \end{cases}
\]
\end{esercizio}

\end{document}
